\documentclass[12pt,a4paper,twoside,openright]{report}

% --- PAKIETY JĘZYKOWE I KODOWANIE ---
\usepackage[utf8]{inputenc}
\usepackage[T1]{fontenc}
\usepackage[polish]{babel}

% --- GEOMETRIA STRONY (Zgodna z wytycznymi PK dla druku dwustronnego) ---
\usepackage{geometry}
\geometry{
    a4paper,
    top=2.5cm,
    bottom=2.5cm,
    inner=3.5cm, % Większy margines od wewnątrz na oprawę pracy
    outer=2.5cm, % Margines zewnętrzny
    bindingoffset=0cm
}

% --- INTERLINIA I AKAPITY ---
\usepackage{setspace}
\onehalfspacing % Interlinia 1,5 (wymóg formalny)
\setlength{\parindent}{1.25cm} % Standardowe wcięcie akapitowe

% --- DODATKOWE PAKIETY (Niezbędne w informatyce) ---
\usepackage{graphicx}  % Do wstawiania rysunków
\usepackage{amsmath,amssymb,amsfonts} % Do zaawansowanych wzorów matematycznych
\usepackage{listings}  % Do wklejania kodów źródłowych programów
\usepackage{amsthm}      % twierdzenia, definicje
\usepackage{booktabs}    % tabele
\usepackage{longtable}   % tabela wieloetapowa (może łamać się między stronami)
\usepackage{array}
\usepackage{xcolor}
\usepackage{tikz}
\usetikzlibrary{arrows.meta}
\usepackage{enumitem}
\usepackage{microtype}
\usepackage{hyperref}  % Do aktywnych linków w tekście i spisie treści

% --- LISTINGI: kod programu jako osobny typ obiektów (numeracja Listing 4.1, 4.2, ...; Spis listingów) ---
\usepackage{caption}
\renewcommand{\lstlistingname}{Listing}
\renewcommand{\lstlistlistingname}{Spis listingów}
% (pakiet listings sam numeruje listingi w obrębie rozdziałów: Listing 4.1, 4.2, ...)
\captionsetup[lstlisting]{labelsep=period,singlelinecheck=false,justification=raggedright}
\lstdefinestyle{pykod}{language=Python,basicstyle=\ttfamily\footnotesize\color{black},breaklines=true,
  frame=single,columns=fullflexible,keepspaces=true,showstringspaces=false,
  commentstyle=\itshape,xleftmargin=0.5em,xrightmargin=0.5em,captionpos=t,aboveskip=1em,belowskip=1em}

% --- BLOKADA NUMERACJI DLA STRONY TYTUŁOWEJ ---
% Zapobiega wyświetlaniu numerów stron na pustych kartach
\usepackage{emptypage} 


\sloppy
\hypersetup{colorlinks=true,linkcolor=black,citecolor=black,urlcolor=black,
pdftitle={Metoda identyfikacji modeli losowych mikrostruktur na podstawie obrazów cyfrowych z wykorzystaniem sum strukturalnych},pdfauthor={Tomasz Ligocki}}

\theoremstyle{definition}
\newtheorem{definicja}{Definicja}[chapter]
\newtheorem{twierdzenie}[definicja]{Twierdzenie}
\newtheorem{stwierdzenie}[definicja]{Stwierdzenie}
\theoremstyle{remark}
\newtheorem{uwaga}[definicja]{Uwaga}

\newcommand{\R}{\mathbb{R}}
\newcommand{\C}{\mathbb{C}}
\newcommand{\Z}{\mathbb{Z}}
\newcommand{\T}{\mathbb{T}}
\newcommand{\cc}[1]{\overline{#1}}
\newcommand{\czer}[1]{{\color{red}#1}}
\newcommand{\oran}[1]{{\color{orange}#1}}
\newcommand{\nieb}[1]{{\color{blue}#1}}
\newcommand{\ziel}[1]{{\color{green!50!black}#1}}
\renewcommand{\Re}{\operatorname{Re}}
\renewcommand{\Im}{\operatorname{Im}}

\begin{document}

% =========================================================================
% STRONA TYTUŁOWA ROZPRAWY DOKTORSKIEJ WIiM PK
% =========================================================================
\begin{titlepage}
    \pdfbookmark[0]{Strona tytułowa}{Strona tytułowa}
    \centering
    
    % Górna sekcja: Uczelnia i Wydział
    {\fontsize{14}{16}\selectfont\bfseries POLITECHNIKA KRAKOWSKA \par}
    {\fontsize{13}{15}\selectfont\bfseries im. Tadeusza Kościuszki \par}
    \vspace{0.4cm}
    {\fontsize{12}{14}\selectfont\bfseries WYDZIAŁ INFORMATYKI I MATEMATYKI \par}
    \vspace{0.2cm}
    {\fontsize{11}{13}\selectfont Rada Naukowa Dyscypliny: Informatyka Techniczna i Telekomunikacja \par}
    
    \vspace{2.5cm}
    
    % Autor pracy
    {\fontsize{13}{15}\selectfont mgr inż. Tomasz Ligocki \par}
    
    \vspace{2.0cm}
    
    % Tytuł główny po polsku
    {\fontsize{16}{19}\selectfont\bfseries Metoda identyfikacji modeli losowych mikrostruktur na podstawie obrazów cyfrowych z wykorzystaniem sum strukturalnych \par}
    
    \vspace{0.6cm}
    
    % Tytuł po angielsku (wymóg WIiM PK)
    {\fontsize{12}{14}\selectfont\itshape A method for identifying random models of microstructures from digital images using structural sums \par}
    
    \vspace{2.5cm}
    
    % Rodzaj pracy
    {\fontsize{12}{14}\selectfont\bfseries ROZPRAWA DOKTORSKA \par}
    
    \vfill
    
    % Sekcja promotorów (wyrównana do prawej na dole)
    \begin{flushright}
    \fontsize{11}{13}\selectfont
    \begin{tabular}{l}
        \textbf{Promotor:} \\
        prof. dr hab. Władimir Mitiuszew \\
        \rule{0pt}{3ex} \\ % Odstęp pionowy między promotorami
        \textbf{Promotor pomocniczy:} \\
        -
    \end{tabular}
    \end{flushright}
    
    \vfill
    \vspace{1.5cm}
    
    % Stopka
    {\fontsize{11}{13}\selectfont Kraków, 2026 \par}
\end{titlepage}

% --- CZYSZCZENIE STYLU DLA STRONY TYTUŁOWEJ ---
\newpage
\thispagestyle{empty}

% =========================================================================
% CZĘŚĆ WSTĘPNA (Rzymska numeracja stron dla spisu treści i streszczeń)
% =========================================================================
\pagenumbering{roman}
\setcounter{page}{2} % Strona tytułowa to fizyczna strona 1 (ukryta)

\chapter*{Streszczenie}
\addcontentsline{toc}{chapter}{Streszczenie}
Rozprawa dotyczy identyfikacji modeli losowych mikrostruktur na podstawie obrazów cyfrowych. Głównym celem części badawczej jest opracowanie i sprawdzenie metody, która zamienia obraz na zbiór środków inkluzji, oblicza z niego sumy strukturalne (wielkości geometryczne wyznaczane wprost z położeń środków) i na tej podstawie ocenia, jakim modelem losowym opisać daną mikrostrukturę, w szczególności czy układ rzeczywisty zachowuje się jak układ losowy o tym samym stężeniu. Metoda dotyczy zarówno mikrostruktur statycznych, takich jak materiały kompozytowe, stopy metali i inne materiały z wtrąceniami, jak i mikrostruktur dynamicznych, w których układ zmienia się w czasie, a obraz jest pojedynczą klatką sekwencji.

Analiza ruchu pływających bakterii, przedstawiona w części badawczej, jest przykładem identyfikacji mikrostruktury dynamicznej. Klatki filmu zamieniono na bazę położeń bakterii (detekcja metodami LoG i DoG na torusie, okna pływające o wymiarach $100\times100$ px) i dla każdego okna obliczono sumy strukturalne, współczynnik $A^{(3)}$ oraz efektywną lepkość zawiesiny dwuwymiarowych dysków. Dla $239$ klatek lepkość jest silnie skorelowana ze stężeniem (średni współczynnik Spearmana $0{,}9825$), lecz średnio w $2{,}43\,\%$ par okien o różnym stężeniu okno o mniejszym stężeniu ma większą lepkość. Zgodnie z teorią geometria układu wchodzi do lepkości od wyrazu rzędu $f^3$, więc takie odstępstwa świadczą o strukturze, której sam model losowy o danym stężeniu nie wyjaśnia. Dla zbiorów generowanych komputerowo (protokoły R, T, P) przedstawiono, na podstawie cyklu prac zespołu promotora, klasyfikację protokołu generowania z cech będących sumami strukturalnymi (dokładność $0{,}7125$). Testy statystyczne hipotez o modelu losowym dla bakterii pozostają poza zakresem rozprawy.

Podstawę teoretyczną stanowi analityczny reprezentatywny element objętości (aRVE) dla dwuwymiarowych, losowych kompozytów dyspersyjnych z nienakładającymi się kołowymi inkluzjami. Przedstawiono sieć okresowości, obszar fundamentalny torusa, sumy sieciowe i funkcje Eisensteina, sumy strukturalne, rewizję koncepcji RVE Hilla (wraz z członem zależnym od kształtu brzegu, wynikającym z warunkowej zbieżności sum sieciowych) oraz twierdzenie o dekompozycji, według którego efektywne stałe są kombinacjami liniowymi sum strukturalnych (część geometryczna) o współczynnikach zależnych od parametrów kontrastu (część fizyczna). Opisano też protokoły generowania konfiguracji losowych, resumację szeregów potęgowych metodą iterowanych pierwiastków z transformacją typu Borela i wyznaczanie wykładnika krytycznego efektywnego modułu ścinania oraz, po zamianie $\mu\to\eta$ przy $\nu=\tfrac12$, efektywnej lepkości. Teoria poprzedza część badawczą, a w dodatku zebrano rozbieżności zauważone w tekstach źródłowych.

\noindent\textbf{Słowa kluczowe:} identyfikacja mikrostruktur, obrazy cyfrowe, sumy strukturalne, mikrostruktury dynamiczne, ruch bakterii, lepkość efektywna, reprezentatywny element objętości, RVE, homogenizacja, funkcje Eisensteina, uczenie maszynowe.

\newpage

\tableofcontents
\newpage

\phantomsection
\addcontentsline{toc}{chapter}{Spis listingów}
\lstlistoflistings
\newpage

% =========================================================================
% CZĘŚĆ GŁÓWNA PRACY
% =========================================================================
\pagenumbering{arabic}

\chapter{Wprowadzenie}

Kompozyt dyspersyjny jest ośrodkiem, w którym inkluzje jednej fazy rozmieszczone są w ciągłej matrycy. Szczególnym przypadkiem są jednokierunkowe kompozyty włókniste: ich przekrój poprzeczny do włókien jest układem dysków na płaszczyźnie. Podstawowym zadaniem teorii homogenizacji jest wyznaczenie stałych efektywnych ośrodka makroskopowo jednorodnego, równoważnego (w sensie średnich pól) ośrodkowi o mikrostrukturze \cite{czapla2025computational}.

Praktyka inżynierska opiera się zwykle na koncepcji reprezentatywnego elementu objętości (RVE) sformułowanej przez Hilla \cite{hill1963elastic,hill1965self}. Jej sformułowanie ma charakter postulatu fizycznego, a nie definicji matematycznej, i (jak pokazano w \cite{czapla2025computational}) bez dodatkowych założeń prowadzi do błędnych wartości stałych efektywnych wyższych rzędów względem stężenia inkluzji $f$. Przyczyną jest zależność granicy $\lim\mu_N$ od kształtu obszaru, na którym stawia się zadanie brzegowe (rozdz.~\ref{sec:hill}). Poprawny wynik uzyskuje się, gdy RVE jest obszarem fundamentalnym sieci okresowości, czyli komórką, której translacje wypełniają płaszczyznę bez nakładania się (Def.~\ref{def:fund}); zadanie okresowe na torusie nie ma wówczas brzegu, a stałe efektywne nie zależą od wyboru kształtu komórki fundamentalnej \cite{czapla2025computational}.

Warunkowa zbieżność, o której tu mowa, dotyczy sumy sieciowej Eisensteina--Rayleigha
\begin{equation*}
{S_2=\sum_{(m_1,m_2)\ne(0,0)}\frac{1}{(m_1\omega_1+m_2\omega_2)^2},}
\end{equation*}
gdzie $\omega_1,\omega_2$ są wektorami okresowości sieci (dla sieci kwadratowej $\omega_1=1$, $\omega_2=i$), a wyraz $m_1=m_2=0$ pomija się. Szereg zbiega tylko warunkowo, więc jego wartość zależy od kolejności sumowania: $S_2=\pi$ dla sumowania Eisensteina (antypłaskie ścinanie, część I) i $S_2=0$ dla sumowania symetrycznego, Maxwella (płaskie odkształcenie, część II, a więc także lepkość). Szczegóły: patrz rozdz.~\ref{ssec:S}.
 

Metody czysto numeryczne (MES) są efektywne dla jednej, ustalonej geometrii, lecz nie nadają się do analizy statystycznej: według \cite{czapla2025computational} już dla jednostajnego rozkładu równych dysków potrzeba około $10^3$ realizacji po $N=100$ inkluzji, a dla $N\approx10^6$ obliczenia są nieosiągalne. Wynika to z natury MES: wynikiem jest liczba dla jednej, ustalonej konfiguracji i jednego zestawu stałych materiałowych $\mu,\nu,\mu_1,\nu_1$, bez jawnej zależności od położeń inkluzji, więc obliczenia trzeba powtarzać dla każdej pary zestawów stałych \cite{czapla2025computational}. Dla około 1000 inkluzji w komórce okresowej bezpośrednie wyznaczenie stałych efektywnych jest praktycznie nieosiągalne; podział komórki na podobszary pozwala badać pola lokalne, ale nie stałe efektywne \cite{mityushev2025computationalpart3}.

Alternatywą rozwijaną w cyklu prac \cite{mityushev2025computational,czapla2025computational,mityushev2025computationalpart3} jest strategia analityczno-numeryczna, w której:
\begin{enumerate}[label=(\roman*),nosep]
\item stałe efektywne przedstawia się jako szeregi potęgowe względem $f$;
\item współczynniki tych szeregów rozkłada się na część fizyczną i część geometryczną (twierdzenie o dekompozycji), przy czym część geometryczną wyraża się przez \emph{sumy strukturalne}, obliczane z położeń środków inkluzji algorytmem o złożoności praktycznie liniowej względem liczby inkluzji;
\item wielomiany względem $f$ (szeregi potęgowe obcięte do skończonego rzędu) resumuje się metodą iterowanych pierwiastków z transformacją typu Borela, aby uzyskać wzory ważne aż do progu perkolacji i oszacować wykładnik krytyczny modułu ścinania;
\item sumy strukturalne stosuje się także jako cechy opisujące mikrostrukturę w uczeniu maszynowym.
\end{enumerate}
Część I cyklu dotyczy przewodnictwa (antypłaskiego ścinania) \cite{mityushev2025computational}, część II --- płaskiej sprężystości \cite{czapla2025computational}, część III --- klasyfikacji mikrostruktur metodami uczenia maszynowego \cite{mityushev2025computationalpart3}. W \cite{czapla2025computational} (rozdz.~6.1) zaznaczono możliwość rozszerzenia podejścia na zawiesiny lepkie, a w \cite{mityushev2025computationalpart3} omówiono przykład zbiorowego zachowania bakterii; obliczeń efektywnej lepkości w tych pracach nie przedstawiono. Niniejsza praca przedstawia ujednolicony opis metody z części II i III, z odwołaniami do części I tam, gdzie jest to potrzebne (np.\ wartość $S_2$ przy sumowaniu Eisensteina). Przeniesienie formalizmu na efektywną lepkość, przez analogię ze sprężystością nieściśliwą ($\nu=\tfrac12$), zapisano w p.~\ref{ssec:visc}, w p.~\ref{ssec:homog} i w rozdz.~\ref{sec:resum}. Jest to wniosek z tej analogii, a nie wynik obliczeń z prac \cite{mityushev2025computational,czapla2025computational,mityushev2025computationalpart3}; o roli lepkości w zawiesinach bakterii zob.\ \cite{aranson2022bacterial}.

Układ rozprawy jest następujący. \ziel{Rozdz.~\ref{sec:stanwiedzy} przedstawia obecny stan wiedzy od zagadnień ogólnych do najnowszych kierunków badań.} Rozdz.~\ref{sec:opracowania} zawiera opracowania wybranych prac z literatury. \emph{Część teoretyczna} obejmuje rozdz.~\ref{sec:def} (definicje i aparat pomocniczy: sieć okresowości, funkcje i sumy Eisensteina, sumy strukturalne, lepkość i stężenie), \ref{sec:model} (sformułowanie problemu), \ref{sec:hill} (rewizja koncepcji Hilla i analityczny RVE), \ref{sec:decomp} (twierdzenie o dekompozycji i wzory na stałe efektywne) \ziel{\ref{sec:splot} (problem obliczeniowy: sumy strukturalne jako dyskretny splot łańcuchowy)} oraz \ziel{\ref{sec:obraz} (analiza obrazu cyfrowego: ośrodki losowe w obrazach i problem jakości obrazów, w tym kompresji H.264)}. Teoria jest przedstawiona przed częścią badawczą, aby wszystkie pojęcia użyte w wynikach (sumy strukturalne, współczynnik $A^{(3)}$, lepkość efektywna) były wcześniej zdefiniowane. \emph{Część badawcza} zaczyna się przeglądem z lotu ptaka (rozdz.~\ref{sec:nowosc}: nowatorstwo i wkład własny), po którym następują: obrazy ruchu bakterii i obliczanie sum strukturalnych z położeń bakterii (rozdz.~\ref{sec:real}), narzędzia informatyczne (rozdz.~\ref{sec:kod}), symulacje zbiorów teoretycznych (rozdz.~\ref{sec:sim}), resumacja szeregów i wykładnik krytyczny (rozdz.~\ref{sec:resum}), sumy strukturalne w uczeniu maszynowym (rozdz.~\ref{sec:ml}) oraz statystyczna obróbka struktur obserwowanych (rozdz.~\ref{sec:realstat}). Zawartość odpowiada planowi badań z projektu badawczego. Rozprawę kończą podsumowanie i dodatek z rozbieżnościami zauważonymi w tekstach źródłowych. Przyjęto identyfikację $\R^2\cong\C$, $z=x_1+ix_2$; $\cc{z}$ oznacza sprzężenie zespolone.


\chapter{\texorpdfstring{\ziel{Obecny stan wiedzy}}{Obecny stan wiedzy}}\label{sec:stanwiedzy}
{\color{green!50!black}
Rozdział przedstawia stan wiedzy od zagadnień ogólnych do najnowszych kierunków badań; jest oparty na wprowadzeniu do artykułu, którego autor rozprawy jest współautorem.

\section{Ośrodki niejednorodne i ich opis statystyczny}
Rozpoznawanie struktury ośrodków rozproszonych jest ważnym zagadnieniem badania materiałów niejednorodnych, w których własności makroskopowe wynikają z przestrzennej organizacji faz. Należą do nich kompozyty z cząstkami wzmacniającymi \cite{cherkaev2012variational,milton2002theory}, materiały porowate \cite{adler2013fractures,adler2013porous}, układy ziarniste \cite{giorgio2023geometrically}, zawiesiny i emulsje \cite{mewis2012colloidal,gluzman2022asymptotics}, materiały optyczne i elektroniczne \cite{barczak2024anisotropy,andrianov2023approximate} oraz tkanki biologiczne \cite{czapla2016random,mityushev2017optimal,ligocki2024application}. W ośrodkach wielofazowych własności mechaniczne, cieplne, elektryczne, optyczne i transportowe zależą nie tylko od stężeń faz, lecz także od kształtu, rozmiaru i rozmieszczenia cząstek, ich grupowania się i łączności \cite{roding2020predicting}.

Podstawową ideą obliczeniowych metod rozpoznawania struktury jest zastąpienie złożonej mikrostruktury zbiorem deskryptorów statystycznych \cite{ohser2000statistical}. Proponowano m.in. dwupunktowe funkcje korelacji (autokorelacji) \cite{jiao2013modeling,li2023improved}, rozkłady odległości do najbliższego sąsiada, statystyki skupisk, funkcje kowariancji i miary anizotropii. Mikrostruktury opisywano także trójpunktową funkcją korelacji \cite{cheng2022data}.

\section{Sumy strukturalne}
Metodę sum strukturalnych wprowadzono dla zagadnień przewodnictwa w \cite{GMN}, a następnie rozszerzono na sprężyste i lepkie kompozyty rozproszone w \cite{drygas2019applied}. Przy wyznaczaniu stałych efektywnych podejście to jest równoważne użyciu funkcji korelacji wyższych rzędów, ale jest znacznie wydajniejsze i pozwala uwzględnić korelacje praktycznie dowolnego rzędu. Sumy strukturalne jako deskryptory dają ilościową charakterystykę niejednorodności struktury i pozwalają porównywać różne struktury na wspólnej podstawie.

\section{Skupiska, perkolacja i analiza obrazu cyfrowego}
Szczególną uwagę poświęca się wykrywaniu skupisk cząstek i aglomeratów \cite{rylko2022windows,mityushev2026microstructural}. W kompozytach niejednorodnych powstawanie rozległych połączonych struktur może istotnie wpływać na efektywne przewodnictwo, przepuszczalność, sztywność lub odporność na pękanie. Teoria perkolacji dostarcza szczególnie użytecznych ram do wskazywania przejść krytycznych związanych z pojawieniem się łączności dalekiego zasięgu. Algorytmy obliczeniowe pozwalają wyznaczać progi perkolacji, wskaźniki łączności i rozkłady rozmiarów skupisk bezpośrednio z obrazów doświadczalnych lub z symulacji numerycznych.

Analiza obrazu cyfrowego stała się jednym z podstawowych narzędzi charakteryzowania struktury \cite{terada1997digital}. Mikrostruktury uzyskane z mikroskopii, tomografii lub generowane numerycznie można zamienić na reprezentacje cyfrowe nadające się do automatycznego przetwarzania \cite{rylko2022windows}. Metody teorii grafów wskazują sieci kontaktów i wzorce łączności, a nowoczesne techniki uczenia maszynowego ułatwiają klasyfikację złożonych mikrostruktur i wykrywanie ukrytych cech struktury, które mogą nie być widoczne w tradycyjnej analizie statystycznej.

\section{Związek struktury z własnościami}
Ważnym elementem rozpoznawania struktury jest ustalenie związków struktura--własności. Obliczeniowo wyznaczone deskryptory struktury można włączyć do modeli analitycznych i teorii homogenizacji, aby przewidywać efektywne własności materiału. Takie podejście łączy geometrię mikroskopową z odpowiedzią makroskopową ośrodków niejednorodnych i umożliwia skuteczniejsze projektowanie i optymalizację materiałów.

\section{Najnowsze kierunki}
Ostatnio łączy się geometrię stochastyczną, fizykę statystyczną, analityczną homogenizację i wielkoskalowe symulacje numeryczne \cite{armstrong2019quantitative,telega2004stochastic,bramati2026periodic}. Ciągłe modele matematyczne ujawniają ogólne prawa rządzące ośrodkami niejednorodnymi \cite{bakhvalov2012homogenisation,jikov2012homogenization}, a symulacje dyskretne oddają oddziaływania między cząstkami i szczegóły lokalnej struktury. Ich połączenie jest skuteczną podstawą badania materiałów wieloskalowych i rozpoznawania struktury w szerokim zakresie skal długości. Metody asymptotyczne \cite{andrianov2020mathematical,andrianov2002asymptotology,andrianov2023approximate} wspierają badania konstruktywne przez wyprowadzanie wzorów analitycznych ważnych dla różnych typów ośrodków niejednorodnych.

\section{Zawiesiny bakterii i miejsce rozprawy}
Strukturalnie nieuporządkowane układy biologiczne, zarówno ciała stałe, jak i zawiesiny, wykazują szerokie spektrum anomalnych zachowań fizycznych. Zespoły bakterii nie mieszczą się w klasycznych kategoriach ciała stałego i cieczy używanych w inżynierii i fizyce, ale można je skutecznie modelować jako zawiesiny lepkie, w których pojedyncze bakterie są sztywnymi inkluzjami w lepkim ośrodku wodnym \cite{ligocki2025study}. Do anomalnych zachowań takich układów należy m.in. rozdział faz ciecz--ciecz \cite{alberti2019considerations}. Zbiorowe zachowanie pływających bakterii \emph{Bacillus subtilis} badano na podstawie obrazów w \cite{ariel2018collective}.

Rozprawa łączy te nurty: obrazy cyfrowe zamienia się na zbiory środków inkluzji, z których oblicza się sumy strukturalne i, w przypadku zawiesin, efektywną lepkość, a następnie wskazuje obszary anomalne, w których mniejsze stężenie cząstek idzie w parze z większą lepkością.
}

\chapter*{Przegląd literatury}

\emph{RVE i zadania brzegowe.} Koncepcję reprezentatywnego elementu objętości sformułował Hill \cite{hill1963elastic}, a dwa zadania brzegowe (liniowe przemieszczenie i jednorodna trakcja) podał w \cite{hill1965self}. Mają one charakter postulatów fizycznych. W \cite{czapla2025computational,mityushev2025computationalpart3} wskazano, że bez dodatkowych założeń o kształcie obszaru prowadzą do błędów wyższych rzędów względem stężenia $f$; dotyczy to także metod rozszerzających obszar do nieskończoności przez dodawanie inkluzji \cite{czapla2025computational,mityushev2022effective1}.

\emph{Sumy sieciowe i zależność od kształtu.} Warunkową zbieżność sum sieciowych zauważył Rayleigh \cite{rayleigh1892lvi} przy badaniu przewodnictwa ośrodka z regularnie ułożonymi przeszkodami. Sieci kul badali McPhedran i McKenzie \cite{mcphedran1978conductivity} oraz McKenzie, McPhedran i Derrick \cite{mckenzie1978conductivity}. Związek sumy Rayleigha z własnościami ciał skończonych i nieskończonych przedstawił Mityushev \cite{mityushev1997transport}. Pojęcie reprezentatywnej komórki w mechanice kompozytów i uogólnionych sum Eisensteina--Rayleigha, na którym opiera się analityczny RVE, pochodzi z \cite{mityushev2006representative}.

\emph{Aparat analityczny.} Funkcje eliptyczne i funkcje Eisensteina opisuje monografia Weila \cite{weil1999elliptic}. Potencjały zespolone w płaskiej sprężystości podano w \cite{muskhelishvili1953some}, a struktury okresowe w \cite{grigolyuk1992periodic}. Metodę równań funkcyjnych i przemiennych iteracji Schwarza dla kompozytów omówiono w \cite{drygas2019applied,mityushev2022effective1}.

\emph{Cykl obliczeniowy.} Strategia analityczno-numeryczna została rozwinięta dla przewodnictwa \cite{mityushev2025computational}, dla płaskiej sprężystości \cite{czapla2025computational} oraz dla klasyfikacji mikrostruktur metodami uczenia maszynowego \cite{mityushev2025computationalpart3}.

\emph{Resumacja szeregów.} Aproksymanty iterowanych pierwiastków z transformacją typu Borela wprowadził Gluzman \cite{gluzman2022iterative}, a wybór wykładników krytycznych przez kontrolę sumowania Borela opisał w \cite{Gluzman2024CriticalIndices}.

\emph{Lepkość i zawiesiny.} Standardowym odniesieniem dla lepkości newtonowskiej i równań Stokesa jest monografia Landaua i Lifszyca \cite{landau1987fluid}, a dla sprężystości \cite{landau1986theory}. Przegląd \cite{aranson2022bacterial} omawia zawiesiny bakterii jako przykład materii aktywnej: według jego streszczenia aktywność bakterii wpływa na lepkość efektywną i dyfuzyjność zawiesiny. W cyklu \cite{mityushev2025computational,czapla2025computational,mityushev2025computationalpart3} efektywnej lepkości nie obliczano, a jedynie wskazano taką możliwość; dlatego w p.~\ref{ssec:visc} zapisano analogię, która pozwala zastosować te same wzory dla $\nu=\tfrac12$.


\chapter{Opracowanie wybranych prac z literatury}\label{sec:opracowania}

W rozdziale zebrano krótkie opracowania dwóch prac cytowanych w rozprawie. Opracowania oparto na streszczeniach i dostępnych fragmentach tych prac; szczegółowe wyniki należy sprawdzać w tekstach źródłowych.

\section{Alberti, Gladfelter i Mittag (2019): rozdział faz ciecz--ciecz i kondensaty biomolekularne}\label{ssec:alberti}

Praca \cite{alberti2019considerations} jest zaproszonym artykułem przeglądowym opublikowanym w czasopiśmie \emph{Cell} \textbf{176}(3), 419--434 (2019). Jej autorzy przedstawiają praktyczne wprowadzenie do badania rozdziału faz ciecz--ciecz (ang.\ \emph{liquid--liquid phase separation}, LLPS) i kondensatów biomolekularnych. Zastrzegają przy tym, że nie jest to pełny opis koncepcji rozdziału faz, a po szczegółowy opis biofizyki oraz przegląd metod doświadczalnych i teoretycznych odsyłają do literatury.

\paragraph{Tło.} Coraz więcej danych wskazuje, że LLPS, czyli kondensacja, leży u podstaw powstawania w komórkach struktur bezbłonowych, takich jak jąderka. Przedziały tworzące się w ten sposób nazwano \emph{kondensatami biomolekularnymi}. To skłoniło badaczy do szeroko zakrojonych prac nad rolą kondensatów w prawidłowych komórkach oraz w zjawiskach od rozwoju organizmu po choroby związane z wiekiem.

\paragraph{Zawartość.} Artykuł opisuje, czym jest LLPS, i omawia \emph{diagramy fazowe}, czyli wyniki doświadczeń określających, w jakich warunkach układ pozostaje jedną, dobrze wymieszaną fazą, a w jakich ulega rozdziałowi faz. Do czynników wpływających na rozdział należą m.in.\ stężenie składników, pH oraz objętość wykluczona przez inne makrocząsteczki, dlatego makrocząsteczki ulegają rozdziałowi faz w odpowiedzi na zmianę warunków. Wyznaczanie diagramów fazowych pozwala ocenić, jak wartościowość (liczba miejsc wiążących) i właściwości chemiczne cząsteczek modulują rozdział faz oraz czy może on zachodzić w warunkach fizjologicznych.

\paragraph{Znaczenie dla rozprawy.} Praca dotyczy układów molekularnych, a nie zawiesin bakterii, więc nie jest źródłem wzorów stosowanych w rozprawie. Łączy ją z tematem rozprawy to, że w obu przypadkach zjawiska zależą od stężenia składników i od ich rozmieszczenia w przestrzeni; zawiesiny bakterii i kondensaty różnią się jednak skalą, mechanizmem i opisem teoretycznym.

\section{Ariel i in. (2018): zbiorowa dynamika dwuwymiarowych pływających bakterii}\label{ssec:ariel}

Praca \cite{ariel2018collective} została opublikowana w \emph{Physical Review E} \textbf{98}, 032415 (2018) przez G.~Ariela, M.~Sidortsov, S.\,D.~Ryana, S.~Heidenreicha, M.~Bära i A.~Be'era. Jest to praca doświadczalno-modelowa o zbiorowym pływaniu bakterii.

\paragraph{Problem.} Własności fizyczne zbiorowo pływających bakterii badano już dokładnie doświadczalnie i symulacyjnie. Modele poprawnie przewidują niektóre aspekty obserwowanej dynamiki, lecz zarówno modele, jak i doświadczenia różnią się założeniami i warunkami fizycznymi, więc nie jest jasne, które modele są właściwe dla których układów doświadczalnych.

\paragraph{Metoda.} Autorzy badają statystykę ruchu dwóch szczepów bakterii \emph{Serratia marcescens}: szczepu dzikiego oraz szczepu nietumblującego (niezmieniającego gwałtownie kierunku ruchu). Bakterie pływają w dwuwymiarowej monowarstwie przy różnych stężeniach. Układ doświadczalny pozwala na bezpośrednie porównanie z wynikami symulacji, które wykonano dwoma typami modeli: dyskretnym (opartym na agentach) i ciągłym.

\paragraph{Wyniki.} Oba modele oddają niektóre aspekty dynamiki, ale nie odtwarzają innych, zwłaszcza przy dużych gęstościach. Ponadto, według abstraktu, wpływ tumblingu jest znacznie większy niż samej dyfuzji rotacyjnej (kątowej).

\paragraph{Znaczenie dla rozprawy.} Praca \cite{ariel2018collective} jest cytowana w rozprawie jako przykład badań zbiorowej dynamiki pływających bakterii; nie jest źródłem pliku wideo Movie S1.avi, z którego pochodzą analizowane klatki (rozdz.~\ref{sec:real}). Źródłem tego pliku jest praca \cite{ryan2013kinetic}, zgodnie z projektem badawczym. Wyniki rozprawy dotyczą wyłącznie położeń bakterii na pojedynczych klatkach i nie dotyczą dynamiki ruchu badanej w pracy Ariela i in. Gatunek bakterii badany w pracy \cite{ariel2018collective}, \emph{Serratia marcescens}, jest inny niż \emph{Bacillus subtilis} na filmie Movie S1.avi.

\part{Część teoretyczna}



\chapter{Definicje i aparat pomocniczy}\label{sec:def}

\section{Sieć okresowości, obszar fundamentalny i torus}\label{ssec:lattice}

Niech $\vec\omega_1,\vec\omega_2\in\R^2$ będą wektorami liniowo niezależnymi nad $\R$ (czyli nieleżącymi na jednej prostej przechodzącej przez zero). Wobec identyfikacji $\R^2\cong\C$ traktujemy je jako liczby zespolone $\omega_1,\omega_2\in\C$, liniowo niezależne nad $\R$. Przyjmujemy normalizację $\omega_1>0$, $\Im\omega_2>0$ oraz $\omega_1\times\Im\omega_2=1$ (pole komórki równe jeden). \emph{Siecią okresowości} nazywamy zbiór
\begin{equation}\label{eq:lattice}
\Lambda=\{P=m_1\omega_1+m_2\omega_2:\ m_1,m_2\in\Z\}.
\end{equation}

\begin{definicja}[obszar fundamentalny]\label{def:fund}
Mierzalny zbiór $Q\subset\C$ nazywamy \emph{obszarem fundamentalnym} sieci $\Lambda$, jeżeli translacje $Q+P$, $P\in\Lambda$, pokrywają całą płaszczyznę i parami nie mają wspólnych punktów wewnętrznych. Równoważnie: $Q$ zawiera (z dokładnością do brzegu) dokładnie jednego przedstawiciela każdej klasy $z+\Lambda$.
\end{definicja}

Standardowym obszarem fundamentalnym jest równoległobok $Q_{0,0}=\{t_1\omega_1+t_2\omega_2:\ 0\le t_1,t_2<1\}$; dla $\omega_1=1$, $\omega_2=i$ jest to kwadrat jednostkowy. Translacje $Q_{p,q}=Q_{0,0}+p\omega_1+q\omega_2$ pokrywają płaszczyznę. \emph{Torusem} nazywamy przestrzeń ilorazową $\T^2=\C/\Lambda$ z metryką
\begin{equation}\label{eq:torusmetric}
\|z\|=\min_{P\in\Lambda}|z+P|.
\end{equation}
Wielkość $\|a-b\|$ jest odległością punktów $a,b$ na torusie (odległością klas $a+\Lambda$ i $b+\Lambda$), bo wartość \eqref{eq:torusmetric} nie zmienia się po zastąpieniu $z$ przez $z+P'$, $P'\in\Lambda$. W szczególności dla kwadratu jednostkowego ($\omega_1=1$, $\omega_2=i$) i punktów $a,b\in Q_{0,0}$ różnica $z=a-b$ leży w kwadracie $(-1,1)^2$, więc w \eqref{eq:torusmetric} wystarcza minimum po $P=s_1+is_2$, $s_1,s_2\in\{-1,0,1\}$ \cite{czapla2025computational}.

\begin{uwaga}
Jeden torus ma nieskończenie wiele obszarów fundamentalnych różnych kształtów (rys.~2 w \cite{czapla2025computational}); wszystkie definiują ten sam problem okresowy. Natomiast dowolny obszar $V\subset\C$ (np. koło), który nie spełnia Definicji~\ref{def:fund}, \emph{nie} jest obszarem fundamentalnym. Rozróżnienie to jest kluczowe w rozdz.~\ref{sec:hill}.
\end{uwaga}

\section{Sumy sieciowe Eisensteina--Rayleigha}\label{ssec:S}

\begin{definicja}\label{def:S}
Sumami sieciowymi Eisensteina--Rayleigha nazywamy
\begin{equation}\label{eq:Sm}
S_m=\sum_{(m_1,m_2)\ne(0,0)}\frac{1}{(m_1\omega_1+m_2\omega_2)^m},\qquad m=2,3,\dots
\end{equation}
\end{definicja}

Dla $m\ge3$ szereg jest zbieżny bezwzględnie. Dla $m=2$ jest zbieżny jedynie \emph{warunkowo}, więc jego wartość zależy od kolejności sumowania \cite{mityushev2025computationalpart3}.

Uzasadnienie. Szereg nazywamy zbieżnym warunkowo, gdy jest zbieżny, ale szereg modułów jego wyrazów jest rozbieżny. W polu komórki równym jeden liczba punktów sieci w kole o promieniu $R$ jest równa w przybliżeniu $\pi R^2$, więc $\sum|P|^{-m}$ zachowuje się jak całka $\int^{\infty} r^{1-m}\,dr$. Dla $m\ge3$ jest ona skończona, zatem szereg zbiega bezwzględnie, a przenumerowanie wyrazów nie zmienia sumy. Dla $m=2$ całka rozbiega się logarytmicznie: szereg zbiega tylko dzięki wzajemnemu znoszeniu się wyrazów o różnych argumentach, bo $P^{-2}=|P|^{-2}e^{-2i\theta}$, $\theta=\arg P$. W takiej sytuacji twierdzenie Riemanna o przestawianiu wyrazów (dla szeregów na płaszczyźnie: twierdzenie Lévy'ego--Steinitza) pozwala zmienić sumę przez zmianę kolejności sumowania. Sieć jest dwuwymiarowa, więc kolejność wyznacza kształt rosnących obszarów, po których sumujemy. Wkład pierścienia między obszarami o brzegach $r=R_1(\theta)$ i $r=R_2(\theta)$ jest w przybliżeniu równy $\int_0^{2\pi}e^{-2i\theta}\ln\bigl(R_2(\theta)/R_1(\theta)\bigr)\,d\theta$. Nie znika on przy $R_{1,2}\to\infty$, gdy obszary mają różne kształty (znika dla obszarów podobnych o tym samym środku i orientacji, bo $\int_0^{2\pi}e^{-2i\theta}d\theta=0$). Dla $m\ge3$ analogiczny wkład jest rzędu $R^{2-m}\to0$.

Dla sieci kwadratowej ($\omega_1=1$, $\omega_2=i$), czyli takiej, że $i\Lambda=\Lambda$, odwzorowanie $P\mapsto iP$ jest bijekcją zbioru $\Lambda\setminus\{0\}$ na siebie. Dla $m\ge3$ (zbieżność bezwzględna, więc przenumerowanie jest dozwolone) daje to $S_m=\sum(iP)^{-m}=i^{-m}S_m$, czyli $(1-i^{-m})S_m=0$, zatem $S_m=0$, gdy $4\nmid m$. Dla $m=2$ ten sam rachunek jest uprawniony tylko wtedy, gdy kolejność sumowania jest niezmiennicza względem obrotu o $\pi/2$ (np.\ sumowanie po rosnących kołach lub kwadratach o środku w zerze, czyli sumowanie symetryczne), i wtedy $S_2=0$. Dla sumowania Eisensteina (iterowanego) warunek ten nie jest spełniony i wychodzi $S_2=\pi$ (Uwaga~\ref{rem:sum}). Wartość $S_4=\Gamma^8(1/4)/(2^6\cdot3\cdot5\,\pi^2)\approx3{,}1512$ \cite{czapla2025computational}, gdzie $\Gamma$ oznacza funkcję gamma Eulera.

\begin{uwaga}[reguła sumowania]\label{rem:sum}
W cyklu \cite{mityushev2025computational,czapla2025computational,mityushev2025computationalpart3} stosuje się dwie reguły: \emph{sumowanie Eisensteina} (iterowane) oraz \emph{sumowanie symetryczne (Maxwella)}. W części I (antypłaskie ścinanie) przyjęto $S_2=\pi$ (sumowanie Eisensteina), w części II (płaska sprężystość) $S_2=0$ (sumowanie Maxwella). Reguła sumowania nie jest kwestią konwencji notacyjnej: wpływa na wyniki, bo odpowiada wyborowi warunków w nieskończoności (rozdz.~\ref{sec:hill}).
\end{uwaga}

\section{Funkcje Eisensteina}

\begin{definicja}[funkcje Eisensteina]\label{def:E}
Funkcje Eisensteina (względem $\Lambda$) definiuje szereg
\begin{equation}\label{eq:Em}
E_m(z)=\sum_{m_1,m_2\in\Z}\frac{1}{(z-m_1\omega_1-m_2\omega_2)^m},\qquad m=2,3,\dots
\end{equation}
\end{definicja}

Dla $m\ge3$ szereg zbiega bezwzględnie i niemal jednostajnie, a $E_m$ jest funkcją meromorficzną, okresową względem $\Lambda$, z biegunami w punktach sieci. Funkcje $E_m$ wyrażają się przez funkcję $\wp$ Weierstrassa \cite{weil1999elliptic}:
\begin{equation}\label{eq:EWeier}
E_2(z)=\wp(z)+S_2,\qquad E_l(z)=\frac{(-1)^l}{(l-1)!}\frac{d^{\,l-2}\wp(z)}{dz^{\,l-2}},\quad l\ge3,
\end{equation}
w szczególności $E_3=-\tfrac12\wp'$ i $E_4=\tfrac16\wp''$. Funkcja $E_2$ zależy od reguły sumowania przez $S_2$. W otoczeniu zera, dla sieci kwadratowej, z klasycznych rozwinięć $\wp$ i $\zeta$ \cite{czapla2025computational} otrzymujemy
\begin{align}
E_3(z)&=\frac1{z^3}-3S_4z-\frac12\sum_{j\ge2}(4j-1)(4j-2)S_{4j}z^{4j-3},\label{eq:E3exp}\\
E_4(z)&=\frac1{z^4}+S_4+\frac16\sum_{j\ge2}(4j-1)(4j-2)(4j-3)S_{4j}z^{4j-4}.\label{eq:E4exp}
\end{align}

\section{Funkcje i sumy Eisensteina--Natanzona--Filsztyńskiego (ang. Eisenstein–Natanzon–Filshtinsky functions and sums)}

Płaska sprężystość wymaga uogólnienia sum i funkcji Eisensteina: do licznika każdego wyrazu wprowadza się sprzężenie zespolone położenia, $\cc{P}^{\,j}$ (odpowiednio $\cc{(z-P)}^{\,j}$), przy czym $j=0$ odpowiada przypadkowi klasycznemu. Dla $j=0,1$ i $m\ge j+2$ określamy
\begin{equation}\label{eq:SmjENF}
S_m^{(j)}=\sum_{P\in\Lambda\setminus\{0\}}\frac{\cc{P}^{\,j}}{P^{m}},\qquad
E_m^{(j)}(z)=\sum_{P\in\Lambda}\frac{\cc{(z-P)}^{\,j}}{(z-P)^m}.
\end{equation}
Dla $j=0$ otrzymujemy $S_m^{(0)}=S_m$ i $E_m^{(0)}=E_m$. Szereg jest zbieżny warunkowo, gdy $m=j+2$ (a zatem $S_3^{(1)}$ wymaga reguły sumowania, Uwaga~\ref{rem:sum}); wówczas $E_m^{(j)}$ nie jest funkcją meromorficzną, gdyż zależy od $\cc z$. Dla sieci kwadratowej Fil\-sztyński--Natanzon wyrazili $E_3^{(1)}$ przez funkcje klasyczne \cite{czapla2025computational}:
\begin{equation}\label{eq:E31}
E_3^{(1)}(z)=\cc{z}\,E_3(z)+\frac1\pi E_4(z)-\frac1\pi\zeta(z)E_3(z)-\frac5\pi S_4+S_3^{(1)},
\end{equation}
gdzie $\zeta$ jest funkcją zeta Weierstrassa. Stała $S_3^{(1)}$ (dla sieci kwadratowej) wynosi
\begin{equation}\label{eq:S31}
S_3^{(1)}=\frac{\Gamma^8(1/4)}{384\pi^3}\approx 2{,}5077\quad\text{(sumowanie Maxwella)},
\end{equation}
a przy sumowaniu Eisensteina jest większa o $\pi/2$; dla sieci heksagonalnej wartości te to $0$ i $\pi/2$ \cite{mityushev2025computationalpart3}. Z rozwinięć \eqref{eq:E3exp}--\eqref{eq:E4exp} wynika, że $E_3^{(1)}(z)-\cc zE_3(z)\to S_3^{(1)}$ przy $z\to0$ \cite{czapla2025computational}; uzasadnia to konwencję przyjmowaną niżej dla jednakowych środków.

\czer{\textbf{Do wyjaśnienia z promotorem.} Prosta suma $\sum\cc P/P^3$ po kwadratach $|m_1|,|m_2|\le M$ zbiega do $2{,}0785$ ($M=100$--$800$), a nie do $2{,}5077$; różnica $0{,}4292$ przenosi się na $E_3^{(1)}(z)$. Wartość $2{,}5077$ (i $2{,}5077+\pi/2$ dla Eisensteina) nie jest więc sumą po kwadratach. Stała $S_3^{(1)}$ wchodzi do $E_3^{(1)}$ jako dodatek addytywny i nie wpływa na okresowość. Wybór reguły sumowania wymaga potwierdzenia.}


\section{Sumy strukturalne}\label{ssec:struct}

Niech $A_N=\{a_1,\dots,a_N\}\subset Q$ będzie zbiorem środków inkluzji w komórce okresowej, a $\mathcal C$ operatorem sprzężenia zespolonego ($\mathcal C^0=\mathrm{id}$, $\mathcal C^1(w)=\cc w$).

\begin{definicja}[suma strukturalna]\label{def:struct}
Niech $p\in\mathbb N$, $\mathbf n=(n_1,\dots,n_p)$, $\mathbf j=(j_1,\dots,j_p)\in\{0,1\}^p$, $\mathbf l=(l_1,\dots,l_p)\in\{0,1\}^p$, $n_s\ge j_s+2$. Sumą strukturalną nazywamy
\begin{equation}\label{eq:struct}
e^{(j_1,\dots,j_p)(l_1,\dots,l_p)}_{n_1,\dots,n_p}=\frac{1}{N^{\alpha/2+1}}\sum_{k_0,k_1,\dots,k_p=1}^{N}\ \prod_{s=1}^{p}\mathcal C^{\,l_s}E^{(j_s)}_{n_s}\bigl(a_{k_{s-1}}-a_{k_s}\bigr),\qquad \alpha=\sum_{s=1}^p(n_s-j_s),
\end{equation}
przy czym dla $a_{k_{s-1}}=a_{k_s}$ przyjmuje się $E^{(j_s)}_{n_s}(0):=S^{(j_s)}_{n_s}$.
\end{definicja}

\czer{\textbf{Do wyjaśnienia z promotorem (zgodność z kodem).} Kod (rozdz.~\ref{sec:kod}) przyjmuje $E_3^{(1)}(0):=0$, a nie $S_3^{(1)}$ jak w powyższej definicji. Dla okien o $N=39$--$43$ bakteriach różnica jest mała (zmiana $\eta_e/\eta$ średnio $5\cdot10^{-5}$, najwyżej $7\cdot10^{-4}$ przy rozrzucie $0{,}023$), ale tekst i kod należy ujednolicić; szczegóły w dodatku, pkt.~\ref{it:weryf}.}

Liczbę $\alpha$ nazywamy \emph{wagą} sumy. Trójkę $(\mathbf n,\mathbf j,\mathbf l)$ nazywamy \emph{wielowskaźnikiem}; to on, a nie jakiś „rząd”, identyfikuje sumę. Poszczególne składowe wielowskaźnika mają następujący sens. Liczba $p$ to liczba czynników w iloczynie, czyli długość łańcucha $k_0\to k_1\to\dots\to k_p$ (łańcuch ma $p+1$ środków i $p$ ogniw). Wektory $\mathbf n$, $\mathbf j$, $\mathbf l$ opisują kolejne ogniwa: $n_s$ jest wykładnikiem w mianowniku (rzędem bieguna) funkcji $E^{(j_s)}_{n_s}$ w $s$-tym czynniku. Liczba $j_s\in\{0,1\}$ mówi, czy w liczniku występuje sprzężenie zespolone argumentu (tzn.\ czy czynnik jest klasyczną funkcją Eisensteina, $j_s=0$, czy funkcją Eisensteina--Natanzona--Filsztyńskiego, $j_s=1$). Liczba $l_s\in\{0,1\}$ mówi, czy cały $s$-ty czynnik jest dodatkowo sprzężony operatorem $\mathcal C^{\,l_s}$. Przykładowo $e^{(11)(10)}_{3,3}$ ma $p=2$, $\mathbf n=(3,3)$, $\mathbf j=(1,1)$, $\mathbf l=(1,0)$, a więc sprzężony jest tylko pierwszy czynnik, oraz $\alpha=(3-1)+(3-1)=4$. Z kolei $e_4^{(0)(1)}$ ma $p=1$, $n_1=4$, $j_1=0$, $l_1=1$ i $\alpha=4$. Czynnik normalizujący $N^{-(\alpha/2+1)}$ zapewnia, że suma ma rząd $\mathcal O(1)$ przy $N\to\infty$ dla jednorodnie rozmieszczonych środków. Pojęcie rzędu odnosi się do współczynników szeregów potęgowych względem $f$ (rozdz.~\ref{sec:decomp}), nie do samych sum strukturalnych. Przykłady dla $p\le2$ \cite{czapla2025computational}:
\begin{align}
e_2&=\frac1{N^2}\sum_{k,m}E_2(a_k-a_m),\qquad e_4:=e_4^{(0)(1)}=\frac1{N^3}\sum_{k_0,k_1}\cc{E_4(a_{k_0}-a_{k_1})},\label{eq:ex1}\\
e_3^{(1)(1)}&=\frac1{N^2}\sum_{k,m}\mathcal C E_3^{(1)}(a_k-a_m),\qquad
e_{3,3}^{(11)(10)}=\frac1{N^3}\sum_{k=1}^N\Bigl|\sum_{m=1}^N E_3^{(1)}(a_k-a_m)\Bigr|^2.\label{eq:ex2}
\end{align}
Suma $e_{2,2}^{(00)(10)}=\frac1{N^3}\sum_{k_0,k_1,k_2}\cc{E_2(a_{k_0}-a_{k_1})}E_2(a_{k_1}-a_{k_2})\equiv e_{2,2}$ jest nieujemna \cite{czapla2025computational}. Warunkiem koniecznym makroskopowej izotropii jest $e_3^{(1)(1)}=0$ (przy sumowaniu Maxwella).

Struktury te są wielopunktowymi funkcjami korelacji w zwartej formie: każdy składnik to iloczyn wzdłuż łańcucha $k_0\to k_1\to\dots\to k_p$. Ponieważ punkty pośrednie $k_1,\dots,k_{p-1}$ są sumowane, a argumentami kolejnych czynników są różnice sąsiednich środków, suma strukturalna jest wielokrotnym (iterowanym) splotem jąder $E^{(j_s)}_{n_s}$ względem miary empirycznej $\mu_N=\sum_{k=1}^N\delta_{a_k}$. Na przykład $\sum_{k_1}E_{n_1}(a_{k_0}-a_{k_1})E_{n_2}(a_{k_1}-a_{k_2})$ jest splotem $E_{n_1}*_{\mu_N}E_{n_2}$ obliczonym w punktach $(a_{k_0},a_{k_2})$. Gdyby $\mu_N$ zastąpić miarą Lebesgue'a, otrzymalibyśmy zwykły splot funkcji. Dla miary punktowej wynik zależy od obu punktów końcowych z osobna, a nie tylko od ich różnicy, więc jest to splot w sensie złożenia operatorów całkowych, a nie splot funkcji jednej zmiennej. Algorytm ich obliczania ma teoretycznie złożoność $\le N^2$, a w praktyce niewiele większą niż $N$ \cite{mityushev2025computationalpart3}.

Zbiór sum potrzebnych do wyznaczenia modułów $\mu_e$ i $k_e$ z dokładnością do wyrazów $f^4$ (reszta $\mathcal O(f^5)$) liczy 14 elementów, $E_{14}(N)=\{E_1,\dots,E_{14}\}$ \cite{mityushev2025computationalpart3}:

\begin{center}
\small
\begin{tabular}{@{}c>{\raggedright\arraybackslash}p{0.78\textwidth}@{}}
\toprule
waga $\alpha$ & sumy strukturalne \\ \midrule
2 & $e_3^{(1)(1)}$ \\
4 & $e_4^{(0)(1)},\ e_{22}^{(00)(10)},\ e_{22}^{(00)(11)},\ e_{33}^{(11)(10)}$ \\
6 & $e_{33}^{(00)(10)},\ e_{43}^{(01)(10)},\ e_{34}^{(10)(10)},\ e_{44}^{(11)(10)},\ e_{333}^{(111)(101)},$\\
 & $e_{223}^{(001)(001)},\ e_{223}^{(001)(110)},\ e_{223}^{(001)(101)},\ e_{232}^{(010)(010)}$ \\ \bottomrule
\end{tabular}
\end{center}

\section{Lepkość i stężenie: pojęcia podstawowe}\label{ssec:visc_conc}

\emph{Lepkość dynamiczna} $\eta$ jest miarą oporu, jaki płyn stawia płynięciu, czyli jego wewnętrznego tarcia. Dla płynu newtonowskiego naprężenie styczne $\tau$ jest proporcjonalne do szybkości ścinania $\dot\gamma$:
\begin{equation}\label{eq:newton}
\tau=\eta\,\dot\gamma ,
\end{equation}
gdzie dla płaskiego przepływu ścinającego $v_x=\dot\gamma\,y$. Jednostką lepkości dynamicznej jest paskal razy sekunda (Pa$\cdot$s). Jej iloraz przez gęstość płynu, $\nu=\eta/\varrho$, nazywa się \emph{lepkością kinematyczną}. Płyn o większej lepkości płynie wolniej pod tym samym naprężeniem; w tym sensie lepkość jest dla płynów odpowiednikiem modułu ścinania $\mu$ ciała sprężystego, a analogię tę wykorzystuje się w dalszej części rozprawy (p.~\ref{ssec:visc}).

\emph{Stężenie} cząstek zawieszonych w płynie opisuje się na kilka sposobów, dlatego w pracy zawsze zaznaczamy, o które chodzi. W mechanice zawiesin stężeniem nazywa się \emph{ułamek objętościowy} cząstek, a w zadaniach płaskich ułamek powierzchniowy $f$, czyli stosunek łącznego pola cząstek do pola obszaru. Dla $N$ identycznych dysków o promieniu $r$ w obszarze o polu $|Q|$
\begin{equation}\label{eq:f_def}
f=\frac{N\pi r^2}{|Q|}.
\end{equation}
Drugą, często używaną wielkością jest \emph{stężenie liczbowe} $n=N/|Q|$, czyli liczba cząstek przypadająca na jednostkę powierzchni (objętości); w mikrobiologii stężenie zawiesiny komórek podaje się zwykle właśnie w tej postaci, jako liczbę komórek w jednostce objętości. Dla identycznych dysków $f=n\pi r^2$. Dla okna pływającego o boku $\ell$ wzór \eqref{eq:f_def} przyjmuje postać $f=N\pi(r/\ell)^2$ (por.\ rozdz.~\ref{sec:real}).

\emph{Lepkość efektywna} $\eta_e$ zawiesiny jest lepkością jednorodnego płynu, który w skali makroskopowej zastępuje płyn z zawieszonymi cząstkami. Definiuje się ją przez średnie naprężenie i średni tensor szybkości odkształcenia, $\langle\sigma'\rangle=2\eta_e\langle e\rangle$, gdzie $\sigma'$ jest dewiatorem naprężenia. Ponieważ obecność sztywnych cząstek hamuje płynięcie, lepkość względna $\eta_e/\eta$ rośnie ze stężeniem. W granicy małych stężeń prowadzi to do wzoru Einsteina $\eta_e/\eta\approx1+\tfrac52 f$ dla sztywnych kul w trzech wymiarach, a dla sztywnych dysków w dwóch wymiarach do $\eta_e/\eta\approx1+2f$ (por.\ \eqref{eq:etareal} w części badawczej) \cite{batchelor1972determination,mewis2012colloidal}. Przy większych stężeniach $\eta_e/\eta$ zależy nie tylko od $f$, ale też od wzajemnego rozmieszczenia cząstek, i właśnie tę zależność opisują sumy strukturalne.

\section{\texorpdfstring{Lepkość i analogia ze sprężystością nieściśliwą}{Lepkość i analogia ze sprężystością nieściśliwą}}\label{ssec:visc}

Dla nieściśliwej cieczy newtonowskiej o prędkości $v=v_1+iv_2$ i ciśnieniu $p$ tensor naprężeń ma postać \cite{landau1987fluid}
\begin{equation*}
{\sigma_{ij}=-p\,\delta_{ij}+2\eta\,\dot\epsilon_{ij},\qquad \dot\epsilon_{ij}=\tfrac12(\partial_iv_j+\partial_jv_i),\qquad \partial_iv_i=0,}
\end{equation*}
gdzie $\eta>0$ jest \emph{lepkością dynamiczną} (jednostka Pa$\cdot$s), czyli współczynnikiem proporcjonalności między naprężeniem stycznym a szybkością ścinania: $\sigma_{12}=2\eta\dot\epsilon_{12}$. Dla powolnych przepływów (mała liczba Reynoldsa) pomija się człon bezwładnościowy w równaniu Naviera--Stokesa i otrzymuje równania Stokesa $\eta\Delta v=\nabla p$, $\partial_iv_i=0$ \cite{landau1987fluid}.

Płaska sprężystość nieściśliwa ($\nu=\tfrac12$) ma tę samą strukturę: $\sigma_{ij}=-p\,\delta_{ij}+2\mu\epsilon_{ij}$, $\partial_iu_i=0$, $\mu\Delta u=\nabla p$ \cite{landau1986theory}. Zadanie Stokesa jest więc formalnie tożsame z zadaniem sprężystości po podstawieniu $u\to v$, $\epsilon_{ij}\to\dot\epsilon_{ij}$, $\mu\to\eta$. W płaskim odkształceniu $\kappa=3-4\nu=1$, więc reprezentacja \eqref{eq:KM} przechodzi w klasyczną reprezentację przepływu Stokesa
\begin{equation*}
{2\eta v=\varphi(z)-z\,\cc{\varphi'(z)}-\cc{\psi(z)},\qquad p=-2\Re\varphi'(z).}
\end{equation*}
Warunek nieściśliwości jest spełniony tożsamościowo, bo $\partial_1v_1+\partial_2v_2=2\Re\,\partial_zv\propto\Re(\varphi'-\cc{\varphi'})=0$, a wzór na $p$ wynika z $\sigma_{11}+\sigma_{22}=-2p=4\Re\varphi'$.

Wynikają stąd uwagi dotyczące dalszych rozdziałów:
\begin{itemize}[nosep]
\item analogia dotyczy płaskiego odkształcenia z $\kappa=1$; dla płaskiego stanu naprężenia $\nu=\tfrac12$ daje $\kappa=\tfrac53$, a nie $1$;
\item gdy cząstki i ciecz są nieściśliwe ($\kappa=\kappa_1=1$), parametry kontrastu \eqref{eq:rho} redukują się do jednej wartości $\varrho_1=\varrho_2=\varrho_3=\varrho=(m-1)/(m+1)$, $m=\eta_1/\eta$; dla $m\to\infty$ (cząstki sztywne) $\varrho=1$, a dla $m\to0$ (inkluzje bez lepkości) $\varrho=-1$;
\item efektywną lepkość $\eta_e$ wyznacza się tymi samymi wzorami co $\mu_e$ (p.~\ref{ssec:homog} i rozdz.~\ref{sec:decomp}) po podstawieniu $\mu\to\eta$, $\kappa=1$; moduł objętościowy nie ma odpowiednika: nie znika, lecz $k\to\infty$ przy $\kappa\to1$ (p.~\ref{ssec:bulk});
\item zakres stosowalności: mała liczba Reynoldsa, idealny kontakt faz (ciągłość prędkości i trakcji, bez napięcia powierzchniowego), cząstki bierne, bez ruchów Browna. W zawiesinach bakterii aktywność zmienia lepkość efektywną i dyfuzyjność \cite{aranson2022bacterial}, czego ten opis nie obejmuje.
\end{itemize}


\chapter{Sformułowanie problemu}\label{sec:model}

\section{Model kompozytu i losowość}\label{ssec:modellos}

Rozważamy inkluzje $D_m=\{z:\ |z-a_m|<r\}$ o module ścinania $\mu_m$ i współczynniku Poissona $\nu_m$, zanurzone w matrycy $D$ o stałych $\mu,\nu$. Warunek braku nakładania ma postać $|a_m-a_k|\ge 2r$, $m\ne k$. Stężenie inkluzji wynosi
\begin{equation}
f=N\pi r^2\qquad(\text{pole komórki}=1).
\end{equation}
Kontakt faz jest idealny (ciągłość przemieszczeń i trakcji). Zbiór środków $A=\{a_1,a_2,\dots\}\subset\C$ traktujemy jako realizację procesu punktowego stacjonarnego w sensie ścisłym (rozkład skończenie wymiarowy jest niezmienniczy względem przesunięć) i ergodycznego, co pozwala wybrać typowy element zbioru losowego do obliczeń \cite{czapla2025computational}. Praktycznie przechodzi się do skończonego zbioru $A_N$ na torusie $\T^2$ z metryką \eqref{eq:torusmetric}; warunek braku nakładania to $\|a_k-a_m\|\ge2r$.

Wprowadzamy parametry kontrastu. Niech $\kappa=3-4\nu$ (płaskie odkształcenie; dla płaskiego stanu naprężenia $\kappa=(3-\nu)/(1+\nu)$), $\kappa_1$ analogicznie dla inkluzji, $m=\mu_1/\mu$:
\begin{equation}\label{eq:rho}
{\varrho_1=\frac{m-1}{m+\kappa_1},\qquad
\varrho_2=\frac{\kappa m-\kappa_1}{\kappa m+1},\qquad
\varrho_3=\frac{m-1}{\kappa m+1},}
\end{equation}
a parametr $\varrho_1$ wyraża się przez pozostałe wzorem
\begin{equation}\label{eq:rho1}
{\varrho_1=\frac{\varrho_3}{1-\varrho_2+\varrho_3}.}
\end{equation}
Mianownik tożsamości \eqref{eq:rho1}, $(m+\kappa_1)/(\kappa m+1)$, jest dodatni. Dla nieskończenie sztywnych inkluzji ($m\to\infty$) i $\kappa=1$ (nieściśliwa matryca) mamy $\varrho_1=\varrho_2=\varrho_3=1$; to przypadek rozważany w rozdz.~\ref{sec:resum}.

\section{Potencjały zespolone}

Naprężenia i przemieszczenia $u=u_1+iu_2$ wyraża się przez funkcje analityczne $\varphi,\psi$ (Kołosow--Muskhelishvili \cite{muskhelishvili1953some}):
\begin{equation}\label{eq:KM}
\sigma_{11}+\sigma_{22}=4\Re\varphi'(z),\quad
\sigma_{22}-\sigma_{11}+2i\sigma_{12}=2\bigl(\cc z\varphi''(z)+\psi'(z)\bigr),\quad
2\mu u=\kappa\varphi(z)-z\cc{\varphi'(z)}-\cc{\psi(z)}.
\end{equation}
W inkluzji $D_m$ występują funkcje $\varphi_m,\psi_m$, a w matrycy $\varphi_0,\psi_0$ i stała $\kappa$. Idealny kontakt na $\partial D_m$ ($t\in\partial D_m$) przyjmuje postać
\begin{align}
\varphi_m(t)+t\cc{\varphi_m'(t)}+\cc{\psi_m(t)}&=\varphi_0(t)+t\cc{\varphi_0'(t)}+\cc{\psi_0(t)}+C_m,\\
\kappa_m\varphi_m(t)-t\cc{\varphi_m'(t)}-\cc{\psi_m(t)}&=\frac{\mu_m}{\mu}\bigl[\kappa\varphi_0(t)-t\cc{\varphi_0'(t)}-\cc{\psi_0(t)}\bigr],
\end{align}
(pierwsze równanie: równowaga sił, drugie: ciągłość przemieszczeń; $C_m$ --- stałe). W matrycy $\varphi_0(z)=bz+\varphi(z)$, $\psi_0(z)=\Gamma z+\psi(z)$, gdzie stałe $b\in\R$, $\Gamma\in\C$ są zadane (określają pole w nieskończoności), a $\varphi,\psi$ spełniają warunki quasi-okresowości
\begin{equation}\label{eq:quasi}
[\varphi(z)]_j=\varphi(z+\omega_j)-\varphi(z)=\xi_j\omega_j,\qquad
[\psi(z)]_j=-\cc{\omega_j}\varphi'(z)+\eta_j\omega_j,\quad j=1,2,
\end{equation}
dla pewnych stałych $\xi_j,\eta_j$. Warunki \eqref{eq:quasi} zapewniają okresowość naprężeń, a przemieszczenia są quasi-okresowe \cite{grigolyuk1992periodic}.

\section{Zasady homogenizacji}\label{ssec:homog}

Sformułowanie to jest zgodne z następującymi zasadami \cite{czapla2025computational}:
\begin{enumerate}[nosep]
\item zadanie formułuje się jako okresowe (quasi-okresowe) na płaszczyźnie z siecią $\Lambda$, równoważnie na torusie $\T^2=\C/\Lambda$, który nie ma brzegu zewnętrznego. Komórka okresowości $Q$ (dowolny obszar fundamentalny, Def.~\ref{def:fund}) jest tylko reprezentantem torusa: na jej brzegu $\partial Q$ nie zadaje się żadnych warunków brzegowych, a pola na przeciwległych brzegach łączą warunki quasi-okresowości \eqref{eq:quasi};
\item warunki quasi-okresowości \eqref{eq:quasi} odpowiadają zadanemu w nieskończoności średniemu tensorowi naprężeń (lub odkształceń);
\item stałe efektywne ośrodka makroskopowo izotropowego wyznacza się przez uśrednienie pól lokalnych po $Q$, $\langle g\rangle=\frac1{|Q|}\int_Q g\,dx_1dx_2$, gdzie $|Q|=1$.
\end{enumerate}
Dla efektywnego modułu ścinania
\begin{equation}\label{eq:mue}
\frac{\mu_e}{\mu}=\frac{\langle\sigma_{12}\rangle}{2\langle\epsilon_{12}\rangle},
\end{equation}
gdzie średnie obliczane są ze wzorów \eqref{eq:KM} dla $b=0$, $\Gamma=i$ (czyste ścinanie), a równania kontaktu rozwiązuje się metodą równań funkcyjnych \cite{drygas2019applied}. Wzór \eqref{eq:mue} zapisano dla $\mu=1$; ogólnie $\mu_e=\langle\sigma_{12}\rangle/(2\langle\epsilon_{12}\rangle)$, czyli $\mu_e/\mu=\langle\sigma_{12}\rangle/(2\mu\langle\epsilon_{12}\rangle)$.

\textbf{Obliczenie średniego odkształcenia.} Dla $u=u_1+iu_2$ mamy $\partial_{\bar z}u=\tfrac12(\epsilon_{11}-\epsilon_{22})+i\epsilon_{12}$, zatem $\langle\epsilon_{12}\rangle=\Im\langle\partial_{\bar z}u\rangle$. Przemieszczenie jest ciągłe na granicach faz, więc z wzoru Greena $\int_Q\partial_{\bar z}u\,dA=\frac1{2i}\oint_{\partial Q}u\,dz$ średnią wyznaczamy z wartości na brzegu komórki (przesuniętej tak, by $\partial Q$ nie przecinał inkluzji). Niech $Q$ będzie równoległobokiem o wierzchołkach $0,\omega_1,\omega_1+\omega_2,\omega_2$ (orientacja dodatnia, bo $\Im\omega_2>0$), a $[u]_j=u(z+\omega_j)-u(z)$ skokami przemieszczenia, które są stałe. Sumując całki po czterech bokach, otrzymujemy $\oint_{\partial Q}u\,dz=\omega_2[u]_1-\omega_1[u]_2$, a ponieważ $|Q|=1$,
\begin{equation}\label{eq:avgstrain}
{\langle\partial_{\bar z}u\rangle=\frac1{2i}\bigl(\omega_2[u]_1-\omega_1[u]_2\bigr).}
\end{equation}
W matrycy $\varphi_0=bz+\varphi$, $\psi_0=\Gamma z+\psi$, więc z \eqref{eq:KM} $2\mu u=(\kappa-1)b\,z-\cc\Gamma\,\cc z+\kappa\varphi-z\,\cc{\varphi'}-\cc\psi$. Funkcja $\varphi'$ jest okresowa, a z \eqref{eq:quasi} mamy $[z\,\cc{\varphi'}]_j=\omega_j\cc{\varphi'}$ oraz $[\cc\psi]_j=-\omega_j\cc{\varphi'}+\cc{\eta_j}\,\cc{\omega_j}$. Wyrazy z $\varphi'$ znoszą się i
\begin{equation*}
{2\mu[u]_j=(\kappa-1)b\,\omega_j-\cc\Gamma\,\cc{\omega_j}+\kappa\xi_j\omega_j-\cc{\eta_j}\,\cc{\omega_j},\qquad j=1,2.}
\end{equation*}
Dla $b=0$, $\Gamma=i$, $\omega_1\in\R$ i $\omega_1\Im\omega_2=1$ podstawienie do \eqref{eq:avgstrain} daje $2\mu\langle\partial_{\bar z}u\rangle=i+W$, gdzie
\begin{equation*}
{W=\frac1{2i}\Bigl[\kappa\,\omega_1\omega_2(\xi_1-\xi_2)-\omega_1\bigl(\omega_2\cc{\eta_1}-\cc{\omega_2}\,\cc{\eta_2}\bigr)\Bigr],\qquad 2\mu\langle\epsilon_{12}\rangle=1+\Im W.}
\end{equation*}
(Składnik $i$ pochodzi z $\omega_2\cc{\omega_1}-\omega_1\cc{\omega_2}=2i$.) Dla ośrodka bez inkluzji ($\varphi=\psi=0$, $\xi_j=\eta_j=0$) mamy $W=0$, a z drugiego wzoru \eqref{eq:KM} $\sigma_{11}=\sigma_{22}$, $\sigma_{12}=1$; stąd $\mu_e=\mu$. W ogólnym przypadku, zgodnie z zasadą 2, przyjmujemy normalizację $\langle\sigma_{12}\rangle=1$ i otrzymujemy
\begin{equation*}
{\frac{\mu_e}{\mu}=\frac1{1+\Im W}.}
\end{equation*}
Wyznaczenie $\mu_e$ sprowadza się więc do wyznaczenia stałych $\xi_j,\eta_j$ z równań kontaktu; ich rozwinięcie względem $f$ prowadzi do wzoru \eqref{eq:mue31} i współczynników $A^{(s)}$ w rozdz.~\ref{sec:decomp} \cite{czapla2025computational,drygas2019applied}.

\textbf{Efektywna lepkość.} Dla $\nu=\tfrac12$ (płaskie odkształcenie, $\kappa=1$) cały opis z p.~\ref{ssec:modellos}--\ref{ssec:homog} dotyczy \emph{zawiesiny}, czyli cieczy lepkiej z cząstkami, po podstawieniu $\mu\to\eta$, $u\to v$, $\epsilon_{ij}\to\dot\epsilon_{ij}$ (p.~\ref{ssec:visc}). Zasady homogenizacji przyjmują wtedy postać:
\begin{enumerate}[nosep]
\item pole prędkości $v$ jest quasi-okresowe na płaszczyźnie z siecią $\Lambda$ (równoważnie określone na torusie), tzn.\ jest sumą przepływu liniowego i części okresowej;
\item przepływ w cieczy bez cząstek zadaje się wzorem \eqref{eq:vflow};
\item lepkość efektywną definiuje wzór \eqref{eq:etae}.
\end{enumerate}

Dla $b=0$, $\Gamma=i$ przepływ w cieczy bez cząstek ma postać
\begin{equation}\label{eq:vflow}
{2\eta\,v=i\,\cc z,\qquad\text{czyli}\qquad v_1=\frac{y}{2\eta},\quad v_2=\frac{x}{2\eta},}
\end{equation}
gdzie $z=x+iy$, a $\eta$ jest lepkością cieczy (matrycy). Jest to czysty, bezwirowy przepływ ścinający ($\partial_1v_2-\partial_2v_1=0$) o szybkości ścinania $\dot\epsilon_{12}=1/(2\eta)$ i naprężeniu stycznym $\sigma_{12}=1$.

Efektywną lepkość zawiesiny określa, analogicznie do \eqref{eq:mue}, stosunek średnich
\begin{equation}\label{eq:etae}
{\eta_e=\frac{\langle\sigma_{12}\rangle}{2\langle\dot\epsilon_{12}\rangle},}
\end{equation}
gdzie $\langle\cdot\rangle$ oznacza średnią po komórce $Q$, $|Q|=1$. Rachunek z poprzedniego punktu (średnie odkształcenie przez skoki przemieszczenia, tu: prędkości) pozostaje w mocy dla $\kappa=1$ i prowadzi do
\begin{equation}\label{eq:etaW}
{\frac{\eta_e}{\eta}=\frac{1}{1+\Im W},}
\end{equation}
gdzie $W$ jest zdefiniowane wyżej (obliczenie średniego odkształcenia), a stałe $\xi_j,\eta_j$ wchodzące do $W$ wyznacza się z równań kontaktu.

\emph{Dyssypacja.} Dla pól okresowych (naprężenie okresowe i bezdywergencyjne, $v=\langle\dot\epsilon\rangle x+v'$ z okresowym $v'$) obowiązuje zależność Hilla--Mandela dla zadania okresowego
\begin{equation}\label{eq:HM}
{\langle\sigma_{ij}\dot\epsilon_{ij}\rangle=\langle\sigma_{ij}\rangle\langle\dot\epsilon_{ij}\rangle.}
\end{equation}
Dla czystego ścinania ($\langle\dot\epsilon_{11}\rangle=\langle\dot\epsilon_{22}\rangle=0$) lewa strona jest średnią szybkością dyssypacji energii $\Phi$, więc z \eqref{eq:HM} i \eqref{eq:etae}
\begin{equation}\label{eq:Phi}
{\Phi=2\langle\sigma_{12}\rangle\langle\dot\epsilon_{12}\rangle=\frac{\langle\sigma_{12}\rangle^2}{\eta_e}.}
\end{equation}
Przy normalizacji $\langle\sigma_{12}\rangle=1$ otrzymujemy $\eta_e=1/\Phi$: przy zadanym naprężeniu większa lepkość efektywna oznacza mniejszą dyssypację.

\emph{Parametr kontrastu i wzór na szereg.} Dla nieściśliwych cząstek o lepkości $\eta_1$ wszystkie parametry kontrastu \eqref{eq:rho} są równe
\begin{equation}\label{eq:varrho}
{\varrho_1=\varrho_2=\varrho_3=\varrho=\frac{m-1}{m+1},\qquad m=\frac{\eta_1}{\eta},}
\end{equation}
a wzór \eqref{eq:mue31} przyjmuje postać
\begin{equation}\label{eq:etaser}
{\frac{\eta_e}{\eta}=\frac{1+\Re\sum_{s\ge1}A^{(s)}f^s}{1-\Re\sum_{s\ge1}A^{(s)}f^s}.}
\end{equation}
Ponieważ $\varrho_1=\varrho_3$, wyraz z $e_{22}$ w \eqref{eq:A3} znika, i
\begin{equation}\label{eq:Avisc}
{A^{(1)}=\varrho,\qquad A^{(2)}=0,\qquad A^{(3)}=\frac{\varrho^2}{\pi^2}\Bigl(4\varrho\,e_{33}^{(11)(10)}+6\,e_4^{(0)(1)}\Bigr),}
\end{equation}
gdzie $e_{33}^{(11)(10)}$ i $e_4^{(0)(1)}$ są sumami strukturalnymi (p.~\ref{ssec:struct}). Wzory \eqref{eq:cs} dają wówczas rozwinięcie
\begin{equation}\label{eq:etaexp}
{\frac{\eta_e}{\eta}=1+2\varrho f+2\varrho^2f^2+2\bigl(A^{(3)}+\varrho^3\bigr)f^3+\dots}
\end{equation}
Dla cząstek sztywnych ($m\to\infty$, $\varrho=1$) z \eqref{eq:etaexp} wynika
\begin{equation}\label{eq:etarigid}
{\frac{\eta_e}{\eta}=1+2f+2f^2+\dots,}
\end{equation}
w tym dwuwymiarowy wyraz Einsteina $1+2f$. Wzór $\mu_J=(1+A_J)/(1-A_J)$ z rozdz.~\ref{sec:sim} (część badawcza) jest właśnie przypadkiem \eqref{eq:etaser} dla cząstek sztywnych, więc podane tam wyniki liczbowe są wynikami dla lepkości zawiesiny sztywnych dysków.

\emph{Sumowanie.} Jak dla płaskiego odkształcenia przyjmuje się $S_2=0$ (sumowanie Maxwella), a wynik zależy od tej reguły (Uwaga~\ref{rem:sum}). Lepkość $\eta_e$ zależy od konfiguracji wyłącznie przez sumy strukturalne, zmieniają się tylko współczynniki fizyczne.

\czer{\textbf{Do wyjaśnienia z promotorem.} W \eqref{eq:Avisc} przyjęto $A^{(2)}=0$, a kod nie liczy $e_3^{(1)(1)}$. Dla 300 losowych okien ($N=39$--$43$) $|e_3^{(1)(1)}|$ ma średnią $0{,}56$ (maks.\ $1{,}5$). Roboczo przyjęty składnik $A^{(2)}=-\frac2\pi\Re e_3^{(1)(1)}$ (nie wynika on z prac źródłowych) zmieniałby $\eta_e/\eta$ średnio o $0{,}017$ (maks.\ $0{,}063$), czyli o około $0{,}7$ odchylenia standardowego rozrzutu $\eta_e/\eta$ w tym zbiorze ($0{,}023$). Założenie $A^{(2)}=0$ trzeba więc uzasadnić.}



\chapter{Rewizja koncepcji Hilla i analityczny RVE}\label{sec:hill}

\section{Definicja Hilla i jej ograniczenia}

Hill \cite{hill1963elastic} określił \emph{objętość reprezentatywną} jako próbkę, która (a) jest strukturalnie typowa dla całej mieszaniny w sensie wielkości średnich oraz (b) zawiera dostatecznie wiele inkluzji, by pozorne moduły efektywne były praktycznie niezależne od wartości przemieszczeń i trakcji zadanych na brzegu, o ile są one „makroskopowo jednorodne” (fluktuują wokół średniej z długością fali małą wobec rozmiaru próbki). Wkład warstwy przybrzeżnej można zaniedbać, biorąc próbkę dostatecznie dużą.

Jest to definicja \emph{operacyjna i fizyczna}: nie zawiera kwantyfikatorów określających, jak duża ma być próbka, jakiego kształtu, ani względem jakiego rozkładu prawdopodobieństwa mówimy o „typowości”. Dla kompozytów losowych jest to istotne ograniczenie, bo istnieje nieskończenie wiele rozkładów nienakładających się dysków, a nie jedynie rozkład jednostajny, przyjmowany milcząco w pracach inżynierskich \cite{czapla2025computational,mityushev2025computationalpart3}.

\section{Zadania brzegowe Hilla}

Dla obszaru $V\subset\C$ z brzegiem $\partial V$ i normalną $n$ spełniona jest zależność energetyczna Hilla--Mandela
\begin{equation}\label{eq:HM}
\int_{\partial V}\bigl(u_i-\langle\epsilon_{ij}\rangle x_j\bigr)\bigl(\sigma_{il}-\langle\sigma_{il}\rangle\bigr)n_l\,ds=0 .
\end{equation}
Na jej podstawie Hill \cite{hill1965self} zaproponował dwa równoważne zadania brzegowe:
\begin{align}
&\text{liniowe przemieszczenie:}\quad u_i=\langle\epsilon_{ij}\rangle x_j\ \text{ na }\partial V, \label{eq:Hill1}\\
&\text{jednorodna trakcja:}\quad \sigma_{il}n_l=\langle\sigma_{il}\rangle n_l\ \text{ na }\partial V. \label{eq:Hill2}
\end{align}
Niech $\mathbf C(\mathbf x)$ oznacza lokalny tensor sprężystości. Po rozwiązaniu \eqref{eq:Hill1} tensor podatności efektywnej wynika z
\begin{equation}\label{eq:Hill3}
\mathbf C_e^{-1}:\langle\boldsymbol\sigma\rangle=\langle\boldsymbol\epsilon\rangle=\langle\mathbf C^{-1}(\mathbf x):\boldsymbol\sigma(\mathbf x)\rangle,
\end{equation}
a po rozwiązaniu \eqref{eq:Hill2} tensor sztywności efektywnej z
\begin{equation}\label{eq:Hill4}
\mathbf C_e:\langle\boldsymbol\epsilon(\mathbf x)\rangle=\langle\boldsymbol\sigma\rangle=\langle\mathbf C(\mathbf x):\boldsymbol\epsilon(\mathbf x)\rangle .
\end{equation}
Podejście to nie zostało uzasadnione w sensie zgodności z zasadami homogenizacji (rozdz.~\ref{sec:model}); potwierdzano je jedynie numerycznie w wybranych przypadkach \cite{czapla2025computational}.

Układ jednorodnych warunków brzegowych na $\partial V$ dla skończonego fragmentu $Q_N$ z inkluzjami przedstawia rys.~\ref{fig:hillV}.

\begin{figure}[htbp]
\centering
\begin{tikzpicture}[scale=0.95]
\draw[very thick] (0,0) circle (3);
\draw[dashed,blue!60!black,thick] (0,0) circle (2);
\fill[gray!45,draw=gray!70] (-1.06,-2.09) circle (0.2);
\fill[gray!45,draw=gray!70] (0.91,-2.57) circle (0.2);
\fill[gray!45,draw=gray!70] (0.22,-0.81) circle (0.2);
\fill[gray!45,draw=gray!70] (-2.65,0.04) circle (0.2);
\fill[gray!45,draw=gray!70] (-0.45,1.96) circle (0.2);
\fill[gray!45,draw=gray!70] (2.15,-1.26) circle (0.2);
\fill[gray!45,draw=gray!70] (-1.15,1.90) circle (0.2);
\fill[gray!45,draw=gray!70] (-1.92,0.49) circle (0.2);
\fill[gray!45,draw=gray!70] (0.83,-0.77) circle (0.2);
\fill[gray!45,draw=gray!70] (0.29,-2.62) circle (0.2);
\fill[gray!45,draw=gray!70] (-1.12,0.51) circle (0.2);
\fill[gray!45,draw=gray!70] (-0.28,-1.20) circle (0.2);
\fill[gray!45,draw=gray!70] (1.77,1.19) circle (0.2);
\fill[gray!45,draw=gray!70] (0.15,2.25) circle (0.2);
\fill[gray!45,draw=gray!70] (1.38,-1.27) circle (0.2);
\fill[gray!45,draw=gray!70] (-2.09,-0.07) circle (0.2);
\fill[gray!45,draw=gray!70] (1.59,0.44) circle (0.2);
\fill[gray!45,draw=gray!70] (0.48,-0.26) circle (0.2);
\fill[gray!45,draw=gray!70] (-0.16,0.98) circle (0.2);
\fill[gray!45,draw=gray!70] (-0.69,1.01) circle (0.2);
\fill[gray!45,draw=gray!70] (-2.22,-1.51) circle (0.2);
\fill[gray!45,draw=gray!70] (-1.33,-0.51) circle (0.2);
\fill[gray!45,draw=gray!70] (-0.85,2.31) circle (0.2);
\fill[gray!45,draw=gray!70] (-1.60,-0.09) circle (0.2);
\fill[gray!45,draw=gray!70] (0.53,-1.42) circle (0.2);
\fill[gray!45,draw=gray!70] (-0.65,-0.61) circle (0.2);
\fill[gray!45,draw=gray!70] (-2.38,0.81) circle (0.2);
\fill[gray!45,draw=gray!70] (-1.75,-2.03) circle (0.2);
\fill[gray!45,draw=gray!70] (-2.39,-0.82) circle (0.2);
\fill[gray!45,draw=gray!70] (0.68,-2.11) circle (0.2);
\fill[gray!45,draw=gray!70] (-0.20,-0.10) circle (0.2);
\fill[gray!45,draw=gray!70] (-0.94,-1.41) circle (0.2);
\fill[gray!45,draw=gray!70] (0.17,-2.12) circle (0.2);
\fill[gray!45,draw=gray!70] (-2.00,1.63) circle (0.2);
\fill[gray!45,draw=gray!70] (0.20,1.67) circle (0.2);
\fill[gray!45,draw=gray!70] (-1.44,1.16) circle (0.2);
\fill[gray!45,draw=gray!70] (0.74,2.40) circle (0.2);
\fill[gray!45,draw=gray!70] (2.04,-0.12) circle (0.2);
\fill[gray!45,draw=gray!70] (0.92,1.80) circle (0.2);
\fill[gray!45,draw=gray!70] (1.50,-0.13) circle (0.2);
\fill[gray!45,draw=gray!70] (1.35,-1.98) circle (0.2);
\fill[gray!45,draw=gray!70] (1.96,-1.73) circle (0.2);
\fill[gray!45,draw=gray!70] (-1.49,-1.24) circle (0.2);
\fill[gray!45,draw=gray!70] (-0.36,-1.90) circle (0.2);
\fill[gray!45,draw=gray!70] (2.47,-0.34) circle (0.2);
\fill[gray!45,draw=gray!70] (-0.13,2.65) circle (0.2);
\fill[gray!45,draw=gray!70] (1.20,2.26) circle (0.2);
\fill[gray!45,draw=gray!70] (0.42,1.20) circle (0.2);
\draw[->,thick] (2.6,-1.5) -- (3.45,-1.97) node[right]{$n$};
\draw[->] (4.3,2.3) node[right]{$\partial Q_N$} -- (2.15,2.15);
\draw[->,blue!60!black] (4.3,0.6) node[right,text width=3.2cm]{$Q_N'$ (daleko od brzegu)} -- (2.02,0.1);
\draw[->] (4.3,-0.6) node[right,text width=3.2cm]{warstwa przybrzeżna} -- (2.65,-0.4);
\end{tikzpicture}
\caption{Schemat obszaru $Q_N$ z nienakładającymi się kołowymi inkluzjami (położenia wylosowane tylko do celów rysunku). Na brzegu $\partial Q_N$ zadaje się warunki \eqref{eq:Hill1} lub \eqref{eq:Hill2}. Linia przerywana ogranicza wnętrze $Q_N'$, w którym wpływ brzegu jest pomijalny; zewnętrzna warstwa przybrzeżna go nie spełnia.}\label{fig:hillV}
\end{figure}

\section{\texorpdfstring{Warunkowa zbieżność sum sieciowych i człon zależny od kształtu}{Warunkowa zbieżność sum sieciowych i człon zależny od kształtu}}

Niech $Q_N$ będzie skończonym fragmentem płaszczyzny (obszarem $V$ z poprzedniego punktu), zawierającym $N$ inkluzji, z brzegiem $\partial Q_N\to\infty$ przy $N\to\infty$. Przez $\sigma^{(N)}_{ij}$ oznaczmy naprężenia z zadania brzegowego w $Q_N$, a przez $\sigma_{ij}$ naprężenia z zadania okresowego (rozdz.~\ref{sec:model}); zachodzi $\sigma_{ij}=\sigma^{(N)}_{ij}+\Delta\sigma^{(N)}_{ij}$ w $Q_N$, gdzie $\Delta\sigma^{(N)}_{ij}$ jest pomijalne we wnętrzu $Q_N'\subset Q_N$, daleko od brzegu. Moduł ścinania uśredniony po $Q_N$ określamy jako
\begin{equation}\label{eq:muN}
\frac{\mu_N}{\mu}=\frac{\langle\sigma_{12}\rangle_N}{2\langle\epsilon_{12}\rangle_N}.
\end{equation}
Pytanie, czy istnieje granica
\begin{equation}\label{eq:limit}
\lim_{\partial Q_N\to\infty}\mu_N=\mu_e,
\end{equation}
i czy równa się modułowi efektywnemu, ma odpowiedź zależną od kształtu: granica zależy od kształtu zewnętrznego brzegu $\partial Q_N$. Źródłem jest warunkowa zbieżność sum sieciowych (por.~$S_2$ w Def.~\ref{def:S}), zauważona już przez Rayleigha \cite{rayleigh1892lvi} przy przejściu od skończonego zbioru dysków do nieskończonej sieci; jest ona wyjaśniana dodatkowym ładunkiem polaryzacyjnym na zewnętrznej powierzchni próbki \cite{mcphedran1978conductivity,mckenzie1978conductivity,mityushev1997transport}. W \cite{mityushev1997transport} (dla przewodnictwa) wykazano, że suma $S_2(\gamma)$ obliczona dla kształtu $\gamma$ oraz człon korekcyjny pochodzący z ciała skończonego (oznaczany tam $\mu(\gamma)$, bez związku z modułem ścinania) sumują się do wartości niezależnej od $\gamma$; dla obszaru igłowatego Rayleigha człon korekcyjny znika, a $S_2=\pi$ dla kwadratu jednostkowego. W \cite{czapla2025computational} zależność od kształtu zapisano za pomocą czynnika $\delta$, nazywanego tam czynnikiem kształtu McPhedrana:
\begin{equation}\label{eq:shape5}
\mu_e=\mu_N(\partial Q_N)+\delta(\partial Q_N)\,\Delta\mu_N(\partial Q_N),
\end{equation}
gdzie $\mu_e$ jest stałą ośrodka zhomogenizowanego, a składniki $\mu_N$, $\Delta\mu_N$ zależą nie tylko od $N$, lecz i od kształtu brzegu. Czynnik $\delta$ zależy \emph{wyłącznie od klasy kształtu}, nie od rozmiaru: klasy obejmują obszary podobne $\{\alpha V\}$, $\alpha>0$, bez obrotów (np.\ prostokąty, koła, elipsy). Według \cite{czapla2025computational}, powołującego się na \cite{rayleigh1892lvi,mcphedran1978conductivity,mckenzie1978conductivity}, czynnik $\delta$ ma następujące własności:
\begin{itemize}[nosep]
\item $\delta(\partial Q_N)=0$ dla kształtu prostokątnego;
\item $\delta(\partial Q_N)\neq0$ do rzędu $\mathcal O(f^3)$ dla niemal wszystkich pozostałych kształtów.
\end{itemize}

Mechanizm zależności od kształtu ilustruje rys.~\ref{fig:shapes}: sumowanie po rosnących obszarach o różnych kształtach daje różne granice dla $S_2$.

\begin{figure}[htbp]
\centering
\begin{tikzpicture}[scale=1.0]
\foreach \i in {-6,...,6}{\foreach \j in {-6,...,6}{\fill[gray] (0.4*\i,0.4*\j) circle (0.035);}}
\draw[thick,blue!60!black] (0,0) circle (2.3);
\draw[thick,red!70!black] (0,0) ellipse (2.7 and 1.5);
\draw (0,0) circle (0.07);
\draw[->,blue!60!black] (-3.9,1.9) node[left]{$R_1(\theta)$ (koło)} -- (-1.6,1.65);
\draw[->,red!70!black] (3.9,1.9) node[right]{$R_2(\theta)$ (elipsa)} -- (2.2,1.1);
\end{tikzpicture}
\caption{Węzły sieci kwadratowej i dwa obszary sumowania o różnych kształtach. Wkład pierścienia między brzegami $R_1(\theta)$ i $R_2(\theta)$ w sumie $S_2$ wynosi w przybliżeniu $\int_0^{2\pi}e^{-2i\theta}\ln\bigl(R_2/R_1\bigr)\,d\theta$ i nie znika, gdy kształty są różne. Zwiększanie obszarów podobnie, o stały czynnik, nie zmienia sumy.}\label{fig:shapes}
\end{figure}
Dla ustalonego kształtu $L$ i $N\to\infty$ wzór \eqref{eq:shape5} przyjmuje postać
\begin{equation}\label{eq:shape6}
\mu_e=\mu_L+\delta(L)\,\Delta\mu_L .
\end{equation}
Wynika stąd: obliczanie $\mu_N$ dla dużego $N$ ma sens jako metoda wyznaczania $\mu_e$ tylko wtedy, gdy $\delta(L)=0$. W przeciwnym razie $\mu_L$ daje $\mu_e$ jedynie z dokładnością rzędu $\mathcal O(f^2)$ \cite{czapla2025computational}. Konsekwencją jest to, że metody, które można traktować jako rozszerzanie $\partial Q_N$ do nieskończoności przez przyrostowe dodawanie inkluzji (schemat różniczkowy, Mori--Tanaki), dają stałe efektywne także tylko do $\mathcal O(f^2)$ \cite{czapla2025computational,mityushev2022effective1}. Jawna postać $\delta(L)$ dla konkretnych kształtów nie jest podana w \cite{czapla2025computational,mityushev2025computationalpart3}; jej sprawdzenie w \cite{rayleigh1892lvi,mcphedran1978conductivity,mckenzie1978conductivity} wymaga dostępu do oryginałów. Zależność $S_2$ od kształtu brzegu $R=R_0h(\theta)$ wyrażono całką po $\theta$ z $h'(\theta)/h(\theta)$ w dodatku do \cite{mityushev1997transport}.

\section{Analityczny reprezentatywny element objętości (aRVE)}

\begin{definicja}[RVE jako obszar fundamentalny]\label{def:rve}
Niech $\Lambda$ będzie siecią \eqref{eq:lattice}, a $A_N=\{a_1,\dots,a_N\}\subset\T^2$ realizacją procesu punktowego spełniającą $\|a_k-a_m\|\ge2r$ ($k\ne m$). \emph{Reprezentatywnym elementem objętości} (RVE) tak opisanej mikrostruktury nazywamy parę $(Q,A_N)$, w której $Q$ jest obszarem fundamentalnym sieci $\Lambda$ w sensie Def.~\ref{def:fund}. Stałe efektywne ośrodka definiuje się wyłącznie przez uśrednienie po $Q$ pól rozwiązania problemu quasi-okresowego \eqref{eq:quasi}.
\end{definicja}

\begin{twierdzenie}[stanowisko autorów \cite{czapla2025computational}]
Podejście Hilla staje się zgodne z matematyczną teorią homogenizacji, gdy dowolny obszar $V$ zastąpi się obszarem fundamentalnym $Q$. Dla obszaru $V$ niefundamentalnego (np.\ koła na rys.~3 w \cite{czapla2025computational}) zgodność nie zachodzi: stałe efektywne są poprawne tylko do $\mathcal O(f^2)$.
\end{twierdzenie}

Różnicę między obszarem fundamentalnym a dowolnym obszarem pokazuje rys.~\ref{fig:fund} (por.\ rys.~2 i 3 w \cite{czapla2025computational}).

\begin{figure}[htbp]
\centering
\begin{minipage}[t]{0.32\textwidth}\centering
\begin{tikzpicture}[scale=1.1]
\foreach \m in {-1,0,1}{\foreach \n in {-1,0,1}{
 \pgfmathsetmacro{\x}{\m+0.5*\n}\pgfmathsetmacro{\y}{0.8660254*\n}
 \draw[gray] (\x,\y) -- ++(1,0) -- ++(0.5,0.8660254) -- ++(-1,0) -- cycle;
 \fill (\x,\y) circle (0.03);}}
\fill[blue!20,opacity=0.8] (0,0) -- ++(1,0) -- ++(0.5,0.8660254) -- ++(-1,0) -- cycle;
\draw[thick,blue!60!black] (0,0) -- ++(1,0) -- ++(0.5,0.8660254) -- ++(-1,0) -- cycle;
\end{tikzpicture}\\[2pt]\small (a)
\end{minipage}\hfill
\begin{minipage}[t]{0.32\textwidth}\centering
\begin{tikzpicture}[scale=1.1]
\foreach \m in {-1,0,1}{\foreach \n in {-1,0,1}{
 \pgfmathsetmacro{\x}{\m+0.5*\n}\pgfmathsetmacro{\y}{0.8660254*\n}
 \begin{scope}[shift={(\x,\y)}]
 \draw[gray] (30:0.57735) -- (90:0.57735) -- (150:0.57735) -- (210:0.57735) -- (270:0.57735) -- (330:0.57735) -- cycle;
 \fill (0,0) circle (0.03);
 \end{scope}}}
\begin{scope}
\fill[blue!20,opacity=0.8] (30:0.57735) -- (90:0.57735) -- (150:0.57735) -- (210:0.57735) -- (270:0.57735) -- (330:0.57735) -- cycle;
\draw[thick,blue!60!black] (30:0.57735) -- (90:0.57735) -- (150:0.57735) -- (210:0.57735) -- (270:0.57735) -- (330:0.57735) -- cycle;
\end{scope}
\end{tikzpicture}\\[2pt]\small (b)
\end{minipage}\hfill
\begin{minipage}[t]{0.32\textwidth}\centering
\begin{tikzpicture}[scale=1.1]
\fill[red!40] (0,0) -- (1,0) -- (0.5,0.8660254) -- cycle;
\foreach \m in {-1,0,1}{\foreach \n in {-1,0,1}{
 \pgfmathsetmacro{\x}{\m+0.5*\n}\pgfmathsetmacro{\y}{0.8660254*\n}
 \draw[gray,fill=white] (\x,\y) circle (0.5);
 \fill (\x,\y) circle (0.03);}}
\fill[blue!20] (0,0) circle (0.5);
\draw[thick,blue!60!black] (0,0) circle (0.5);
\fill[red!70!black] (0.5,0.2887) circle (0.07);
\draw[->,thick,red!70!black] (1.9,1.55) node[above]{szczelina} -- (0.56,0.33);
\end{tikzpicture}\\[2pt]\small (c)
\end{minipage}
\caption{Sieć o podstawowych translacjach $\omega_1,\omega_2$ i jej translacje obszaru (zaznaczonego na niebiesko): (a) równoległobok rozpięty na $\omega_1,\omega_2$ i (b) sześciokąt Wignera--Seitza są obszarami fundamentalnymi tej samej sieci, a translacje każdego z nich wypełniają płaszczyznę bez nakładania; (c) koło nie jest obszarem fundamentalnym: jego translacje zostawiają szczeliny (zaznaczone na czerwono; każda leży między trzema stykającymi się kołami) (a koło o większym promieniu nakładałoby się na sąsiednie). Schemat ma charakter ilustracyjny (sieć heksagonalna).}\label{fig:fund}
\end{figure}

Wyjaśnienie intuicyjne: w problemie okresowym na $Q$ ośrodek ma brzeg \emph{pusty} (torus nie ma brzegu), więc nie ma ładunku polaryzacyjnego na powierzchni próbki i zależności od kształtu nie ma; w obszarze niefundamentalnym brzeg istnieje, a razem z nim człon $\delta(L)\Delta\mu_L$ w \eqref{eq:shape6}. Wybór $Q$ pomiędzy różnymi obszarami fundamentalnymi tego samego torusa (Uwaga po Def.~\ref{def:fund}) nie wpływa na wynik.

Mówiąc o \emph{analitycznym} RVE (aRVE) \cite{mityushev2006representative} mamy na myśli następującą konstruktywną wersję tej idei, realizowaną w \cite{mityushev2025computational,czapla2025computational}: poza warunkiem $Q$ fundamentalny, zależność stałych efektywnych od konfiguracji $A_N$ dana jest \emph{jawnie} --- przez skończoną liczbę sum strukturalnych (Def.~\ref{def:struct}) dla zadanego rzędu względem $f$ (rozdz.~\ref{sec:decomp}). Formalna definicja aRVE pochodzi z \cite{mityushev2006representative}; w \cite{czapla2025computational,mityushev2025computationalpart3} podano jedynie jej treść operacyjną, przytoczoną powyżej.


\chapter{Twierdzenie o dekompozycji i wzory na stałe efektywne}\label{sec:decomp}

\section{Metoda Schwarza i dekompozycja}

Niech $u_0$ będzie zadanym polem zewnętrznym, a $u_k$ polem w $k$-tej inkluzji $D_k$. Uogólniona metoda przemiennych iteracji Schwarza dla nienakładających się inkluzji prowadzi do równania operatorowego w przestrzeni Banacha \cite{mityushev2025computationalpart3}:
\begin{equation}\label{eq:schwarz}
u_k=\varrho_kG_ku_k+\sum_{m\ne k}\varrho_mG_mu_m+u_0\quad\text{w }D_k,\ k=1,\dots,N,
\end{equation}
gdzie $G_m$ są operatorami ograniczonymi, zależnymi od geometrii (a niejawnie od $f$), \emph{ale nie od stałych materiałowych}, a $\varrho_m$ parametrami kontrastu. Metoda kolejnych przybliżeń daje szereg potęgowy względem $\varrho_k$ o współczynnikach czysto geometrycznych (złożenia $G_{m_1}\circ\dots\circ G_{m_K}u_0$). Uśrednienie po komórce prowadzi do następującego wyniku.

\begin{twierdzenie}[o dekompozycji \cite{mityushev2025computationalpart3}]
Tensor efektywnych właściwości jest kombinacją liniową czysto geometrycznych parametrów inkluzji o współczynnikach określonych przez lokalne stałe fizyczne.
\end{twierdzenie}

W realizacji dla kołowych inkluzji część geometryczną stanowią sumy strukturalne (Def.~\ref{def:struct}), a część fizyczną --- parametry kontrastu $\varrho_1,\varrho_2,\varrho_3$ z \eqref{eq:rho}.

\section{Moduł ścinania}

Dla makroskopowo izotropowego kompozytu z jednakowych kołowych inkluzji \cite{czapla2025computational,drygas2019applied}
\begin{equation}\label{eq:mue31}
\frac{\mu_e}{\mu}=\frac{1+\Re\sum_{s\ge1}A^{(s)}f^s}{1-\kappa\,\Re\sum_{s\ge1}A^{(s)}f^s},
\end{equation}
gdzie współczynniki $A^{(s)}$ są kombinacjami liniowymi sum strukturalnych. Rzędem współczynnika $A^{(s)}$ nazywamy wykładnik $s$ potęgi $f^s$, przy której występuje. Pierwsze współczynniki mają postać
\begin{align}
A^{(1)}&=\varrho_3,\qquad A^{(2)}=-\frac2\pi\varrho_3^2\,e_3^{(1)(1)}=0,\label{eq:A12}\\
A^{(3)}&=\frac{\varrho_3}{\pi^2}\Bigl[4\varrho_3^2e_{33}^{(11)(10)}+6\varrho_3e_4^{(0)(1)}+\frac{\varrho_1-\varrho_3}{1+\varrho_1}\bigl(e_{22}^{(00)(10)}-e_{22}^{(00)(11)}\bigr)\Bigr],\label{eq:A3}
\end{align}
przy czym $A^{(2)}=0$ wynika z izotropii. Współczynnik $A^{(4)}$ zawiera sumy z $p\le3$ (wzór (41) w \cite{czapla2025computational}), a kolejne, $A^{(5)}$, $A^{(6)}$, są coraz bardziej złożone. Ze sprawdzenia wzorów (34)--(43) z \cite{czapla2025computational} wynika prawidłowość: \emph{współczynnik rzędu $s$, $A^{(s)}$, zawiera wyłącznie sumy strukturalne o wadze $\alpha=2(s-1)$ i czynnik $\pi^{-(s-1)}$}. Np. $A^{(2)}$ zawiera $e_3^{(1)(1)}$ ($\alpha=2$), $A^{(3)}$ --- sumy o $\alpha=4$, $A^{(4)}$ --- sumy o $\alpha=6$ (zob.~tabela w p.~\ref{ssec:struct}).

Rozwinięcie \eqref{eq:mue31} względem $f$ ma postać $\mu_e/\mu=1+c_1f+c_2f^2+c_3f^3+c_4f^4+\mathcal O(f^5)$. Podstawiając $u=\sum A^{(s)}f^s$ i rozwijając $(1+u)/(1-\kappa u)=1+(1+\kappa)u/(1-\kappa u)$, przy $A^{(2)}=0$ otrzymujemy
\begin{equation}\label{eq:cs}
c_1=(1+\kappa)\varrho_3,\qquad c_2=(1+\kappa)\kappa\varrho_3^2,\qquad c_3=(1+\kappa)\bigl(A^{(3)}+\kappa^2\varrho_3^3\bigr),\ \dots
\end{equation}
Zatem $c_1,c_2$ nie zależą od położenia inkluzji (pokrywają się z rozwinięciami dolnych ograniczeń Hashina--Shtrikmana dla $\mu_1>\mu$ \cite{mityushev2025computationalpart3}), a zależność od geometrii zaczyna się od wyrazu $f^3$. Stąd kryterium weryfikacji wzorów empirycznych: formuła zawierająca człon $f^3$ niezależny od geometrii musi zostać odrzucona \cite{czapla2025computational}. (Dla twardych inkluzji, $\kappa=\varrho_i=1$, mamy $c_1=2$, $c_2=2$, $c_3=2(A^{(3)}+1)$.)

\section{Moduł objętościowy}\label{ssec:bulk}

Dla modułu objętościowego $k_e$ otrzymano wzór
\begin{equation}\label{eq:ke44}
\frac{k_e}{k}=\frac{1-f+\Re\sum_{s\ge1}\bigl(B^{(s)}+C^{(s)}\bigr)f^s}{1-f+\Re\sum_{s\ge1}\bigl(B^{(s)}+\frac{k}{k_1}C^{(s)}\bigr)f^s},
\end{equation}
gdzie $B^{(1)}=B^{(2)}=0$, $B^{(3)}=\frac{2}{\pi^2}\frac{\varrho_3-\varrho_1}{1+\varrho_1}\,e_{22}^{(00)(01)}$, $C^{(1)}=\frac{\varrho_1(1+\varrho_3)}{\varrho_3(1+\varrho_1)}$, $C^{(2)}=0$, a wyższe współczynniki podano w \cite{czapla2025computational} (wzory (46)--(53)). Obliczenia numeryczne dla $k_e$ pozostawiono w \cite{czapla2025computational} do osobnej publikacji.

\textbf{Uwaga dla lepkości ($\nu=\tfrac12$).} W przypadku lepkim moduł objętościowy nie znika, lecz rośnie nieograniczenie. Dla $k=\lambda+\mu=2\mu/(\kappa-1)$ (płaskie odkształcenie) mamy $k\to\infty$ przy $\kappa\to1$, a znika podatność objętościowa $1/k$. Ciśnienie $p$ nie jest wtedy określone przez odkształcenie, lecz jest mnożnikiem Lagrange'a warunku $\partial_iv_i=0$ (p.~\ref{ssec:visc}); lepkość objętościowa $\zeta$ \cite{landau1987fluid} nie występuje, bo mnoży $\partial_iv_i=0$. Wzór \eqref{eq:ke44} nie ma więc zastosowania: stosunek $k/k_1=(\mu/\mu_1)(\kappa_1-1)/(\kappa-1)$ jest nieoznaczony ($0/0$), gdy $\kappa,\kappa_1\to1$, a $k_e$ nie jest wielkością niezależną. Dla lepkości pozostaje tylko $\eta_e$ z p.~\ref{ssec:homog}.

\chapter{\texorpdfstring{\ziel{Problem obliczeniowy: sumy strukturalne jako dyskretny splot łańcuchowy}}{Problem obliczeniowy: sumy strukturalne jako dyskretny splot łańcuchowy}}\label{sec:splot}

{\color{green!50!black}
Wzory na stałe efektywne (rozdz.~\ref{sec:decomp}) wymagają obliczenia sum strukturalnych (Def.~\ref{def:struct}) dla zadanego układu $N$ środków. Rozdział ten opisuje, na czym polega trudność obliczeniowa i jak jest rozwiązywana algorytmem macierzowo-wektorowym obliczania sum strukturalnych, opisanym w rozdz.~2.3 monografii \cite{drygas2019applied} (por.\ też \cite{nawalaniec2020effective}).

\section{Suma strukturalna jako splot łańcuchowy}
Dla czynników $F_s(z)=\mathcal C^{\,l_s}E^{(j_s)}_{n_s}(z)$, $s=1,\dots,p$, suma strukturalna \eqref{eq:struct} jest, z dokładnością do czynnika normalizującego $N^{-(\alpha/2+1)}$, wielowymiarowym \emph{dyskretnym splotem łańcuchowym}
\begin{equation}\label{eq:splot}
\Sigma=\sum_{k_0,k_1,\dots,k_p=1}^{N}\ \prod_{s=1}^{p}F_s\bigl(a_{k_{s-1}}-a_{k_s}\bigr).
\end{equation}
Nie jest to splot w sensie przetwarzania sygnałów, $\sum_k f(k)\,g(s-k)$, w którym wynik zależy tylko od różnicy punktów i można go liczyć transformatą Fouriera (FFT). Tutaj punkty $a_k$ leżą nieregularnie, wynik zależy od punktów końcowych $k_0,k_p$ z osobna, a jądra $F_s$ są funkcjami eliptycznymi o biegunach. Książka \cite{drygas2019applied} zaznacza, że (\ref{eq:splot}) jest szczególnym przypadkiem splotu wielowymiarowego, a nie tym samym obiektem.

\section{Naiwne obliczanie i jego koszt}
Bezpośrednie sumowanie po wszystkich krotkach $(k_0,\dots,k_p)$ obejmuje $N^{p+1}$ składników, z których każdy jest iloczynem $p$ czynników. Koszt rośnie wykładniczo z długością łańcucha $p$ (w książce: $\mathcal O(pN^{p})$ mnożeń). Dla pojedynczych sum o małym $p$ (np.\ $e_2$, $e_4$ przy $p=1$ oraz $e_{22}$, $e_{33}$ przy $p=2$) jest to jeszcze akceptowalne, ale wyższe rzędy rozwinięcia po $f$ wymagają sum o większym $p$, więc potrzebny jest algorytm niewykładniczy.

\section{Algorytm macierzowo-wektorowy}
Niech $M_s\in\mathbb C^{N\times N}$, $(M_s)_{kl}=F_s(a_k-a_l)$. Sumowanie po kolejnych indeksach łańcucha można wykonywać po kolei, bo $k_{s}$ występuje tylko w czynnikach $s$ i $s+1$:
\begin{equation}\label{eq:wektor}
\Sigma=\mathbf 1^{\!\top}M_1M_2\cdots M_p\,\mathbf 1,\qquad v_p=\mathbf 1,\quad v_{s-1}=M_s v_s,\quad \Sigma=\mathbf 1^{\!\top}v_0 .
\end{equation}
Reprezentacja wektorowa (Tw.~2.3 w \cite{drygas2019applied}) wymaga $p$ mnożeń macierzy $N\times N$ przez wektor, czyli $\mathcal O(pN^2)$ działań, zamiast $\mathcal O(pN^p)$; koszt jest liniowy względem $p$ przy ustalonym $N$. Reprezentacja macierzowa (iloczyn wszystkich $M_s$) kosztuje $\mathcal O(pN^3)$, ale daje od razu całą macierz splotu, co jest potrzebne, gdy dalsze obliczenia używają wartości dla par punktów końcowych. Jest to w istocie programowanie dynamiczne (przeorganizowanie sumy wielokrotnej w ciąg sum jednokrotnych) i dokładnie ten sam wynik w arytmetyce zmiennoprzecinkowej, z dokładnością do kolejności sumowania. Dla $p=2$ algorytm sprowadza się do policzenia sum wierszowych $\sum_m E_3^{(1)}(a_k-a_m)$, co robi kod z rozdz.~\ref{sec:kod} dla $e_{33}$.

\section{Rekurencja symboliczna a obliczanie liczbowe}
W rozdz.~2.2.1.1 książki (rys.~2.1) pokazano \emph{rekurencję symboliczną} $A_q=\pi^{-1}\beta A_{q-1}$, z której powstają wzory na kolejne współczynniki rozwinięcia. Liczba składników $A_q$ podwaja się z każdym $q$. To wzrost rozmiaru wzoru analitycznego, a nie kosztu liczenia splotu, i nie jest wadą metody, tylko własnością rozwinięcia. Problem obliczeniowy z \eqref{eq:splot} dotyczy dopiero numerycznego wyznaczenia sum dla konkretnego zbioru punktów; rozdz.~2.3.5 książki przenosi ten sam schemat na przypadek sprężysty (przez funkcje $F_m$), nie wprowadzając nowego algorytmu.

\section{Zastrzeżenia}
\begin{itemize}
\item Zysk algorytmu jest istotny dla długich łańcuchów (wysokie rzędy). Dla sum używanych w tej pracy ($p\le2$) jest umiarkowany; bez niego koszt byłby $\mathcal O(N^3)$, a z nim $\mathcal O(N^2)$.
\item Podawane w książce czasy nie obejmują wyznaczenia samych wartości $E_n(a_k-a_l)$. To $N^2$ wywołań funkcji eliptycznych (u nas funkcji $\wp$, procedura \texttt{pfast.py}), które w praktyce dominują czas obliczeń.
\item Pamięć: macierze $M_s$ mają rozmiar $N^2$; rekurencja z zapamiętywaniem wyników pośrednich dodatkowo zużywa pamięć rosnącą z rzędem.
\item W przypadku sprężystym ważne jest, które czynniki są sprzężone ($l_s$) i w jakiej kolejności mnoży się macierze: pominięcie sprzężenia w $e_4$ (p.~\ref{ssec:struct}) jest błędem tego typu.
\end{itemize}
}


\chapter{\texorpdfstring{\ziel{Analiza obrazu cyfrowego: ośrodki losowe w obrazach i problem jakości obrazów}}{Analiza obrazu cyfrowego: ośrodki losowe w obrazach i problem jakości obrazów}}\label{sec:obraz}
{\color{green!50!black}
Teoria z poprzednich rozdziałów działa na zbiorze środków inkluzji $\mathcal A_N=\{a_1,\dots,a_N\}$. W części badawczej zbiór ten nie jest zadany, lecz trzeba go wyznaczyć z obrazu cyfrowego. Rozdział opisuje tę analizę obrazu, wskazuje, jakie ośrodki losowe pojawiają się w obrazach, i stawia problem jakości danych: film \texttt{Movie S1.avi} jest zapisany z kompresją stratną H.264, która wprowadza artefakty wpływające na wykryte położenia.

\section{Obraz jako realizacja ośrodka losowego}
Klatkę filmu traktujemy jako obraz w skali szarości $I(x,y)$, w którym bakterie widać jako jasne plamy na ciemniejszym tle. Model matematyczny jest następujący: położenia środków $a_k$ są realizacją losowego układu nienakładających się dysków o jednakowym promieniu $r$ w płaszczyźnie, a obraz jest ich zniekształconą, zdyskretyzowanej projekcją. Z obrazu odzyskuje się zbiór $\mathcal A_N$ przez detekcję plam, a następnie liczy sumy strukturalne dla okien traktowanych jak komórki okresowe (torus), por.\ rozdz.~\ref{sec:model} i \ref{sec:real}. Układ rzeczywisty porównuje się z układami losowymi o tej samej liczbie $N$ i tym samym stężeniu $f=N\pi(r/\ell)^2$, generowanymi protokołami R, T i P (rozdz.~\ref{sec:sim}), w szczególności metodą RSA z odległością toroidalną. W tym sensie obrazy dostarczają dwóch rodzajów ośrodków losowych: ośrodka obserwowanego (położenia bakterii z klatek) i ośrodków odniesienia (zbiory generowane komputerowo), a sumy strukturalne są wspólną miarą, która pozwala je porównywać.

\section{Użyte metody analizy obrazu}
W rozprawie użyto następujących narzędzi analizy obrazu; ich realizację w programie opisuje rozdz.~\ref{sec:kod}.
\begin{itemize}[nosep]
\item \textbf{Torus.} Klatka nie jest okresowa, a teoria wymaga komórki okresowej. Aby uniknąć efektów brzegowych, do każdego boku klatki (i do narożników) doklejono pasek o szerokości $24$ px skopiowany z przeciwległej krawędzi i zachowuje się tylko środki leżące w oryginalnej klatce. Odległości liczy się w metryce torusa.
\item \textbf{Binaryzacja Otsu i otwarcie morfologiczne.} Próg Otsu wybiera wartość $t^*$ maksymalizującą wariancję międzyklasową, $t^*=\arg\max_t\,\omega_0(t)\,\omega_1(t)\,[\mu_0(t)-\mu_1(t)]^2$, gdzie $\omega_i,\mu_i$ są udziałami i średnimi jasnościami klas tła i plam. Obraz binarny czyści się otwarciem morfologicznym elementem w kształcie krzyża $3\times3$.
\item \textbf{Detekcja plam LoG i DoG.} Dla jasnego dysku o promieniu $r$ znormalizowana skalowo odpowiedź $\sigma^2\nabla^2(G_\sigma*I)$ (Laplacian of Gaussian, LoG) ma ekstremum przy $\sigma=r/\sqrt2$; dla $r=3{,}5$ px daje to $\sigma\approx2{,}5$, więc przeszukuje się zakres $\sigma\in[2,6]$ px. Różnica gaussianów (Difference of Gaussians, DoG), $G_{k\sigma}-G_\sigma\approx(k-1)\sigma^2\nabla^2G_\sigma$, przybliża LoG i jest tańsza obliczeniowo. Wyniki obu metod łączy się.
\item \textbf{Czyszczenie punktów.} Punkty bliższe niż $r_{\min}=8$ px (w metryce torusa) uznaje się za jedną bakterię. Wynikiem jest zbiór środków $\mathcal A_N$ klatki.
\item \textbf{Okna pływające.} Klatkę dzieli się na okna $100\times100$ px, przesuwane o stałe kroki, a każde okno przeskalowuje się do kwadratu jednostkowego i traktuje jak torus.
\end{itemize}

\section{Problem jakości obrazów: kompresja H.264 i inne artefakty}
Obrazy użyte w rozprawie mają niską jakość i wynika to z danych źródłowych, a nie z obróbki. Film \texttt{Movie S1.avi} ma według odczytu metadanych parametry: kodek wideo H.264 (dekoder \texttt{ffh264}), kontener AVI, plik zapisany biblioteką \texttt{Lavf54.63.104} (libavformat, z rodziny FFmpeg), rozdzielczość $640\times652$ px, $25$ klatek na sekundę, czas $9{,}56$ s, przepływność $1\,876\,512$ bit/s. Daje to $239$ klatek. Przy tej przepływności na jedną klatkę przypada średnio ok.\ $1\,876\,512/25\approx75$ kbit, czyli ok.\ $9{,}4$ kB, a ok.\ $0{,}18$ bita na piksel. Dla porównania nieskompresowana klatka $640\times652$ px w skali szarości o $8$ bitach na piksel zajmuje ok.\ $417$ kB, więc stopień kompresji jest rzędu $44:1$ (przy założeniu $8$ bitów na piksel w skali szarości). Zapis klatek do plików PNG jest bezstratny, ale zapisuje już zdekodowane, czyli zniekształcone obrazy: artefaktów nie usuwa. Ponieważ metadane pliku zawierają znacznik biblioteki zapisującej, a nie historii nagrania, nie wiadomo, czy materiał był kompresowany wcześniej; jeśli był, do artefaktów dochodzi utrata jakości przy kolejnym kodowaniu.

\subsection{Skąd biorą się artefakty H.264}
Kodek H.264 jest stratny i działa blokowo. Klatkę dzieli się na makrobloki $16\times16$ px, a resztę przewidywania koduje się całkowitoliczbową transformacją podobną do dyskretnej transformaty kosinusowej na blokach $4\times4$ lub $8\times8$ px, po czym kwantuje się współczynniki; kwantyzacja odrzuca drobne szczegóły. Ponadto klatki są kodowane jako klatki I (niezależne), P i B (przewidywane z kompensacją ruchu z innych klatek). Typowe skutki dla obrazów takich jak nasze to:
\begin{itemize}[nosep]
\item \emph{efekt blokowy:} widoczne granice bloków i skokowe zmiany jasności na ich krawędziach (filtr deblokujący je łagodzi, ale rozmywa szczegóły);
\item \emph{rozmycie i utrata drobnych struktur:} kwantyzacja wysokich częstotliwości spłaszcza małe, słabo kontrastowe obiekty;
\item \emph{pierścieniowanie i szum komarowy (ang.\ mosquito noise):} fałszywe falowanie jasności wokół ostrych krawędzi;
\item \emph{artefakty czasowe:} zmiana jakości między klatkami I, P i B (miganie jasności lub kształtu plamy w kolejnych klatkach) oraz smużenie i duchy wynikające z kompensacji ruchu poruszających się obiektów.
\end{itemize}

\subsection{Dlaczego jest to istotne dla sum strukturalnych}
Bakteria ma na obrazie średnicę $2r=7$ px, a więc rozmiar porównywalny z blokiem transformaty ($4\times4$ lub $8\times8$ px), a mniejszy niż makroblok ($16\times16$ px). Kwantyzacja działa więc na obiekt w całości: może przesunąć środek plamy, zmienić jej kształt i jasność, a przy dwóch bliskich plamach scalić je lub rozdzielić. Do tego dochodzą ograniczenia samego obrazu: mała rozdzielczość ($\approx140$ nm na piksel), skala szarości, rozmycie optyczne układu obrazującego, szum i nierównomierne oświetlenie, a także uproszczenie modelu (bakterie pałeczkowate przedstawiamy jako dyski).

Skutki dla wykrywania możemy zapisać jako model błędu: zamiast $\mathcal A_N$ otrzymujemy $\mathcal A'_{N'}$, w którym część punktów jest przesunięta, $a_k'=a_k+\varepsilon_k$, część pominięta (plamy o małym kontraście), część fałszywa (artefakt wykryty jako plama), a część bliskich par jest scalona. Wrażliwość sum strukturalnych na takie błędy wynika wprost z ich budowy: funkcje $E_3^{(1)}$ i $E_4$ mają bieguny (jak $z^{-3}$ i $z^{-4}$) w punktach $z=0$, więc największy wkład do $e_{33}$ i $e_4$ dają pary najbliższych sąsiadów, a to właśnie odległość między nimi najbardziej cierpi z powodu artefaktów i czyszczenia przy progu $r_{\min}$.

\subsection{Postawienie problemu}
Problem jest więc następujący: \emph{na ile wyniki wyznaczone z obrazów skompresowanych H.264 (sumy strukturalne, współczynnik $A^{(3)}$, lepkość efektywna i wskazane anomalie) odzwierciedlają rzeczywisty układ bakterii, a na ile są skutkiem artefaktów?} W rozprawie ograniczono wpływ artefaktów środkami opisanymi wyżej (łączenie metod LoG i DoG, torus, czyszczenie punktów przy $r_{\min}=8$ px), ale ilościowo go nie oszacowano. Naturalny sposób oceny, nie wykonany w rozprawie, to wygenerować syntetyczne obrazy ze znanymi położeniami środków, zakodować je kodekiem H.264 z podobną przepływnością (np.\ bibliotekami z rodziny FFmpeg), wykryć środki tym samym algorytmem i porównać $e_{33}$, $e_4$, $A^{(3)}$ i $\eta_e/\eta$ z wartościami dla położeń zadanych. Do czasu takiej oceny wnioski o anomaliach lepkości należy traktować jako dotyczące obrazów, a nie bezpośrednio samych bakterii.
}

\part{Część badawcza}


% ===== NOWY ROZDZIAŁ: nowatorstwo i wkład własny (przegląd z lotu ptaka) =====
{
\chapter{Nowatorstwo i wkład własny rozprawy: przegląd z lotu ptaka}\label{sec:nowosc}

Rozdział przedstawia całość badań z lotu ptaka: jaki problem rozwiązuje rozprawa, co jest w niej nowatorskie oraz w jaki sposób zrealizowano poszczególne elementy. Szczegóły techniczne zawierają kolejne rozdziały, a opis programu --- rozdz.~\ref{sec:kod}.

\section{Problem i pomysł}
Zbiorowy ruch bakterii bada się zwykle funkcjami korelacji i symulacjami cząstek (rozdz.~\ref{sec:opracowania} i przegląd literatury). Takie metody opisują wąski zestaw cech mikrostruktury, a wyznaczenie własności makroskopowych zawiesiny, np.\ lepkości efektywnej, z samych położeń bakterii wymaga teorii homogenizacji. Teoria analitycznego reprezentatywnego elementu objętości (aRVE) \cite{mityushev2006representative,drygas2019applied} wyraża własności efektywne dwuwymiarowych zawiesin kołowych inkluzji przez \emph{sumy strukturalne}, które oblicza się wprost z położeń środków (rozdz.~\ref{sec:def}).

Pomysł rozprawy polega na tym, by klatki filmu z pływającymi bakteriami potraktować jak dwuwymiarowy kompozyt losowy: zamienić je na zbiory środków bakterii, policzyć dla każdego okna sumy strukturalne i lepkość efektywną, a następnie zbadać, czy rzeczywisty układ bakterii zachowuje się jak układ losowy o tym samym stężeniu. Daje to ilościowe kryterium zbiorowego ruchu bakterii oparte na sumach strukturalnych \cite{ligocki2025study}.

\section{Co jest nowatorskie}
Nowatorstwo rozprawy ma charakter głównie obliczeniowo-aplikacyjny: znane narzędzie teoretyczne zastosowano do rzeczywistych obrazów bakterii i rozszerzono je o obliczenia, których wcześniej nie wykonywano w tej skali. Wkład można ująć w siedmiu punktach.
\begin{enumerate}[nosep]
\item \textbf{Sumy strukturalne dla rzeczywistych obrazów bakterii.} Metodę aRVE zastosowano do klatek filmu \texttt{Movie S1.avi}, a nie do zbiorów generowanych komputerowo, i oparto na niej kryterium wykrywania zbiorowego ruchu bakterii \cite{ligocki2024application,ligocki2025study}. Zrealizowano to przez zamianę klatek na bazę położeń bakterii (rozdz.~\ref{sec:real}).
\item \textbf{Topologia torusa dla obrazów rzeczywistych.} Teoria aRVE wymaga komórki okresowej, a klatka filmu nie jest okresowa. Topologię torusa stosowano dotąd do zbiorów syntetycznych \cite{nawalaniec2020effective}; w rozprawie rozszerzono ją na obrazy rzeczywiste: doklejono ramki o szerokości $24$ px z przeciwległych krawędzi, detekcję wykonano na torusie, bliskie punkty scalano w metryce torusa, a okna traktowano jako torusy (Listing~\ref{lst:kod-torus}, \ref{lst:kod-clean}, \ref{lst:kod-dup}).
\item \textbf{Łączona detekcja LoG i DoG z czyszczeniem danych.} Wyniki obu metod łączy się i scala punkty leżące bliżej niż $r_{\min}=8$ px \cite{ligocki2025study}. W rozprawie czyszczenie wykonuje się w metryce torusa metodą klastrów (algorytm union--find), a detekcję przeprowadza się raz dla całej klatki (Listing~\ref{lst:kod-logdog}, \ref{lst:kod-clean}).
\item \textbf{Skala obliczeń: okna pływające dla wszystkich klatek.} Dla $239$ klatek i $2500$ okien $100\times100$ px na klatkę, czyli dla $597\,500$ okien, oblicza się sumy $e_2$, $e_{2,2}$, $e_{33}$ i $e_4$. Obliczenia wykonuje się wektorowo (CuPy), a wyniki tworzą bazę danych geometrycznych rozprawy (rozdz.~\ref{sec:kod}).
\item \textbf{Lepkość efektywna z położeń bakterii.} Teorię lepkości dwuwymiarowych zawiesin \cite{drygas2019applied} zastosowano do rzeczywistej struktury: dla każdego okna oblicza się współczynnik $\Re A^{(3)}$ i lepkość $\eta_e/\eta$ (Listing~\ref{lst:kod-A3}). Praca \cite{ligocki2025study} ograniczała się do $e_2$ i $e_{2,2}$.
\item \textbf{Anomalie lepkości.} Zbadano pary okien, w których okno o mniejszym stężeniu ma większą lepkość (rozdz.~\ref{sec:realstat}). Takie pary stanowią średnio $2{,}43\,\%$ wszystkich par o różnym stężeniu, mimo że współczynnik Spearmana między $f$ i $\eta_e/\eta$ wynosi średnio $0{,}9825$. Wynik ten wynika z teorii: wyrazy $f$ i $f^2$ nie zależą od układu bakterii, a geometria wchodzi od wyrazu $f^3$ (p.~\ref{ssec:visc}). Najsilniejszą anomalię znajduje się programowaniem dynamicznym (Listing~\ref{lst:kod-anom}); przykładem jest klatka $233$: okno $B$ ma $43$ bakterie i lepkość $1{,}4100$, a okno $A$ tylko $39$ bakterii i lepkość $1{,}4190$.
\item \textbf{Krytyczne uporządkowanie teorii.} Rozprawa zestawia w jeden wykład teorię z cyklu prac \cite{czapla2025computational,mityushev2025computationalpart3}, sprawdza ją numerycznie i wskazuje rozbieżności w tekstach źródłowych (dodatek, m.in.\ konwencja znaku $e_{33}$).
\end{enumerate}

\section{Procedura krok po kroku}
Tabela~\ref{tab:workflow} zestawia osiemnaście etapów procedury badawczej z tabeli~6 pracy \cite{ligocki2025study} z ich realizacją w rozprawie. Etapy 1--13 tworzą potok, który w \cite{ligocki2025study} wykonano na pierwszej klatce; w rozprawie wykonano go dla wszystkich klatek i rozszerzono o lepkość. Etapy 14--16 (analiza szumu, miary $R$ i $I$, hipotezy $H_1$ i $H_2$) zapowiedziano w \cite{ligocki2025study} jako dalszy plan i nie są objęte rozprawą.
{\small\renewcommand{\arraystretch}{1.15}
\begin{longtable}{@{}r >{\raggedright\arraybackslash}p{3.6cm} >{\raggedright\arraybackslash}p{6.3cm} l@{}}
\caption{Procedura badawcza według tabeli~6 w \cite{ligocki2025study} i jej realizacja w rozprawie.}\label{tab:workflow}\\
\toprule
Nr & Działanie (ECMS 2025) & Realizacja w rozprawie & Status\\
\midrule
\endfirsthead
\toprule
Nr & Działanie (ECMS 2025) & Realizacja w rozprawie & Status\\
\midrule
\endhead
\bottomrule
\endfoot
1 & Wyodrębnienie klatek: \texttt{Movie-S1.avi} $\to$ PNG & Film zamieniono na $239$ klatek PNG (rozdz.~\ref{sec:real}) & tak\\
2 & Wybór klatki: analiza pierwszej klatki & Analiza wszystkich klatek $i=1,\dots,239$ (etap 18) & tak\\
3 & Wstępne przetwarzanie: przycięcie, skala szarości & Odcięcie pasków ($6+8$ px), skala szarości, doklejenie ramek $24$ px (torus) & tak\\
4 & Binaryzacja: optymalizacja parametrów & Próg Otsu i otwarcie morfologiczne (Listing~\ref{lst:kod-otsu}) & tak\\
5--7 & Wykrywanie środków: LoG, DoG, DoH; wybór najlepszej & Połączone LoG i DoG (Listing~\ref{lst:kod-logdog}); DoH odrzucona, jak w \cite{ligocki2025study} & tak\\
8 & Czyszczenie danych: połączenie LoG i DoG, scalanie bliskich par & Scalanie w metryce torusa, $r_{\min}=8$ px (Listing~\ref{lst:kod-clean}) & tak\\
9 & Obliczanie $e_2$, $e_{2,2}$ w oknie pływającym & Okna $100\times100$ px, $2500$ okien na klatkę; dodatkowo $e_{33}$, $e_4$ (Listing~\ref{lst:kod-e33}) & tak\\
10 & Generowanie zbioru: rozkład syntetyczny & Losowe zbiory o tym samym $N$ (rozdz.~\ref{sec:kod}) & tak\\
11 & Analiza zbioru: obliczenie $e_2$, $e_{2,2}$ & Te same funkcje co dla danych rzeczywistych & tak\\
12 & Porównanie powierzchni 3D (zbiór rzeczywisty i losowy) & Mapy $N$ i okna $A$, $B$ dla danych rzeczywistych (rozdz.~\ref{sec:realstat}) & częściowo\\
13 & Statystyka: średnie, odchylenia & Statystyka po $239$ klatkach: udział par anomalnych, $\rho_S$ (rozdz.~\ref{sec:realstat}) & tak\\
14 & Analiza szumu: $\Re e_2$ --- korelacja, $\Im e_2$ --- izotropia & Poza zakresem rozprawy & nie\\
15 & Miary $R=\Re e_2^{\mathrm{real}}-\Re e_2^{\mathrm{rand}}$, $I=\Im e_2^{\mathrm{real}}-\Im e_2^{\mathrm{rand}}$ & Poza zakresem rozprawy & nie\\
16 & Hipotezy: $H_1$ --- $R$ nie jest białym szumem, $H_2$ --- $I$ jest białym szumem & Poza zakresem rozprawy & nie\\
17 & Rozszerzenie: $e_{3,3}$, $e_{4,4}$ itd. & Zrealizowano $e_{33}$ i $e_4$ (lepkość); $e_{4,4}$ i wyższe --- dalsze prace & częściowo\\
18 & Pełna analiza: klatki $2,3,\dots,239$ & Klatki $1,\dots,239$ & tak\\
\end{longtable}}

\section{Podział wkładu}
Podstawy teoretyczne --- teoria aRVE, twierdzenie o dekompozycji, wzory na stałe efektywne, resumacja oraz klasyfikacja mikrostruktur metodami uczenia maszynowego --- pochodzą z cyklu prac zespołu Mityusheva \cite{mityushev2006representative,drygas2019applied,czapla2025computational,mityushev2025computationalpart3}; rozdziały teoretyczne rozprawy je omawiają i porządkują. Wkład własny autora obejmuje: przygotowanie obrazów i bazy położeń bakterii, program przetwarzania obrazu i obliczania sum strukturalnych (rozdz.~\ref{sec:kod}), zastosowanie teorii lepkości do rzeczywistych struktur, statystykę i badanie anomalii dla $239$ klatek (rozdz.~\ref{sec:realstat}) oraz numeryczną weryfikację wzorów teoretycznych.

\section{Ograniczenia i dalsze kierunki}
Rozprawa opiera się na jednym filmie, a wynik zależy od parametrów detekcji ($\sigma\in[2,6]$, próg $0{,}1$, $r_{\min}=8$ px). Okna zawierają kilkadziesiąt bakterii, więc pojedyncze różnice lepkości są małe (np.\ $1{,}4190$ i $1{,}4100$), a ich istotność statystyczną trzeba badać testami. Naturalnym uzupełnieniem są etapy 14--16 z tabeli~\ref{tab:workflow} (analiza szumu i hipotezy $H_1$, $H_2$) oraz sumy wyższych rzędów $e_{4,4}$ i $e_{2,2,2,2}$.
}


\chapter{\texorpdfstring{Obrazy ruchu bakterii i obliczanie sum strukturalnych z ich położeń}{Obrazy ruchu bakterii i obliczanie sum strukturalnych z ich położeń}}\label{sec:real}

Aparat opisany w części teoretycznej (rozdz.~\ref{sec:def}--\ref{sec:decomp}) zastosowano do rzeczywistej struktury, którą są położenia bakterii na klatkach filmu. Pytanie badawcze jest następujące: czy lokalnie, w małych oknach, większe stężenie $f$ zawsze oznacza większą lepkość efektywną $\eta_e/\eta$, czy o wartości $\eta_e/\eta$ decyduje także układ środków bakterii, poprzez współczynnik $A^{(3)}$ (por.~\eqref{eq:Avisc})? Zgodnie z oznaczeniami p.~\ref{ssec:visc} lepkość efektywną oznaczamy $\eta_e/\eta$, ponieważ w tym artykule symbol $\mu_e/\mu$ jest zarezerwowany dla efektywnego modułu ścinania; na rysunkach, generowanych programem obliczeniowym, ta sama wielkość podpisana jest $\mu_e/\mu$.

\section{\texorpdfstring{Dane i procedura obliczeń}{Dane i procedura obliczeń}}

\emph{Dane.} Źródłem jest plik wideo Movie S1.avi uzyskany w pracy \cite{ryan2013kinetic} (por.\ tamże o lepkości efektywnej zawiesin pływających bakterii); przedstawia on rzeczywisty ruch dwuwymiarowych pływających bakterii \emph{Bacillus subtilis}. Film zamieniono na serię 239 klatek w formacie PNG; każda klatka jest obrazem w skali szarości o wymiarach $640\times652$ px z powiększonym widokiem bakterii w roztworze (opis zbioru danych i wyodrębniania klatek podano w \cite{ligocki2025study}). Przed analizą usuwa się czarne paski o wysokości 6~px u góry i 8~px u dołu. Bakterie modelujemy jako identyczne, nienakładające się dyski o promieniu $r=3{,}5$ px; położeniem dysku jest środek bakterii.

\emph{Wykrywanie środków.} Wzory z rozdz.~\ref{sec:def} dotyczą struktury podwójnie okresowej, więc każdą klatkę traktuje się jak torus: do każdego boku i narożnika dołącza się pasek 24~px skopiowany z przeciwległej krawędzi. Obraz binaryzuje się metodą Otsu i oczyszcza otwarciem morfologicznym (element krzyżowy $3\times3$). Środki bakterii wyznacza się dwoma detektorami plam: Laplasjanem z gaussjanu (LoG) i różnicą gaussjanów (DoG) (szerokości $\sigma\in[2,6]$ px, próg względny 0,1, promienie plam 2--8 px). Oba zbiory środków łączy się, pozostawia tylko środki leżące w oryginalnej klatce i scala punkty bliższe niż 8~px (odległość torusowa).

\emph{Okna pływające i obliczenia.} Po klatce przesuwa się kwadratowe okno o boku $\ell=100$ px na siatce $50\times50$ położeń, czyli 2500 okien w klatce. Środki bakterii z okna przeskalowuje się do kwadratu jednostkowego $Q$; powstaje zbiór $N$ punktów $a_k$, przy czym $N$ jest wielkością lokalną (okna z $N<3$ pomija się, a do obliczeń bierze co najwyżej 500 punktów). Stężenie lokalne wynosi
\begin{equation}\label{eq:freal}
{f=N\pi\Bigl(\frac r\ell\Bigr)^2.}
\end{equation}
Dla każdego okna oblicza się sumy strukturalne $e_{33}=e_{33}^{(11)(10)}$ i $e_4=e_4^{(0)(1)}$ (p.~\ref{ssec:S}) z funkcjami $E_3=-\wp'/2$ i $E_4=\wp^2-g_2/12$, $g_2=60S_4\approx189{,}0727$, $g_3=0$ (sieć kwadratowa), używając dla każdej różnicy $a_k-a_m$ najbliższej kopii torusowej; funkcję $\wp$ liczy się wektorowo (GPU lub CPU). Dla cząstek sztywnych ($\varrho=1$) z \eqref{eq:Avisc} otrzymuje się $A^{(3)}=\frac2{\pi^2}\bigl(2e_{33}+3e_4\bigr)$, a lepkość efektywną okna z \eqref{eq:etaser}, przy $A^{(1)}=f$ i $A^{(2)}=0$:
\begin{equation}\label{eq:etareal}
{\frac{\eta_e}\eta=\frac{1+s}{1-s},\qquad s=f+\Re A^{(3)}f^3.}
\end{equation}
W tabelach podano $e_4'=\Re e_4$. Wzór \eqref{eq:etareal} jest stosowalny dla dysków bez nakładania, tj.\ gdy $2r\le 8$ px; tu $2r=7$ px.


% ===== NOWY ROZDZIAŁ: opis kodu (notebook Concentration_function_real_400_CPU_lepkosc.ipynb) =====
% Fragmenty kodu przytoczono z wersji v16 skryptu (eksport notebooka); komentarze z polskimi znakami pominięto.
{
\chapter{Narzędzia informatyczne: kod przetwarzania obrazu i obliczania sum strukturalnych}\label{sec:kod}

Rozdział opisuje program, który zamienia kolejne klatki filmu (rozdz.~\ref{sec:real}) w bazę położeń bakterii i oblicza dla pływających okien sumy strukturalne oraz efektywną lepkość. Program ma postać notebooka Pythona \texttt{Concentration\_function\_real\_400\_CPU\_lepkosc.ipynb} z repozytorium \cite{ligocki_doktorat_github}. Dla każdej z klatek $i=1,\dots,239$ wykonuje on następujące kroki:
\begin{enumerate}[nosep]
\item odcina czarne paski u góry i u dołu klatki i dokleja ramki o szerokości $\delta=24$ px z przeciwległych krawędzi, tak aby obraz miał topologię torusa;
\item binaryzuje obraz (progiem Otsu) i wykrywa środki bakterii dwiema metodami: LoG i DoG;
\item łączy wyniki obu metod, odrzuca punkty spoza oryginalnej klatki i scala punkty leżące zbyt blisko siebie (w metryce torusa);
\item dzieli klatkę na $50\times50$ okien pływających o rozmiarze $100\times100$ px, w każdym oknie usuwa duplikaty torusowe i normalizuje położenia do kwadratu jednostkowego;
\item oblicza dla okna sumy strukturalne $e_2$, $e_{2,2}$, $e_{33}$, $e_4$, współczynnik $\Re A^{(3)}$, stężenie $f$ i lepkość efektywną $\eta_e/\eta$;
\item wyszukuje pary okien anomalnych i generuje losowe zbiory odniesienia o tym samym $N$.
\end{enumerate}
Wszystkie wyniki zapisuje do plików CSV, które tworzą bazę danych geometrycznych rozprawy. Fragmenty kodu omawiane w tekście przytoczono jako listingi (Listing~\ref{lst:kod-param}--\ref{lst:kod-anom}); są numerowane osobno, niezależnie od rysunków i tabel, i zestawione w Spisie listingów.

\section{Parametry programu}
Rozmiary klatki i okna oraz progi czyszczenia zebrano na początku notebooka (Listing~\ref{lst:kod-param}):
\begin{lstlisting}[style=pykod,caption={Parametry programu: rozmiary klatki i okna oraz progi czyszczenia},label={lst:kod-param}]
i_min, i_max = 1, 239
y_w_orig, x_w_orig = 652, 640
Y_CUT_TOP = 6
Y_CUT_BOT = 8
y_w = y_w_orig - Y_CUT_TOP - Y_CUT_BOT    # = 638
x_w = x_w_orig                             # = 640
y_d, x_d = 100, 100
r_min = 8
r_min_detekcji = 2
r_max_detekcji = 8
DIST_SCALANIE = 10
delta = 3 * r_min         # = 24
LICZBA_OKIEN_Y = 50
LICZBA_OKIEN_X = 50
BLOB_MIN_SIGMA = 2
BLOB_MAX_SIGMA = 6
BLOB_NUM_SIGMA = 5
BLOB_THRESHOLD = 0.1
\end{lstlisting}
Klatka ma więc po odcięciu pasków $638\times640$ px, okno $100\times100$ px, a promień wykrywanych plam mieści się w przedziale $2$--$8$ px. Okna przesuwa się co $\lfloor538/49\rfloor=10$ px w pionie i $\lfloor540/49\rfloor=11$ px w poziomie (zmienne \texttt{krok\_y}, \texttt{krok\_x}), co daje $2500$ okien na klatkę.

\section{Odcięcie pasków i doklejanie ramek 24 px (torus)}
Funkcja \texttt{przygotuj\_torus(i)} wczytuje klatkę, odcina czarne paski i dokleja z czterech stron oraz w rogach pasy o szerokości $\delta$ skopiowane z przeciwległych krawędzi (Listing~\ref{lst:kod-torus}):
\begin{lstlisting}[style=pykod,caption={Odcięcie pasków i doklejanie ramek $\delta=24$ px (budowa torusa)},label={lst:kod-torus}]
img_bez_paskow = img[Y_CUT_TOP:y_w_orig - Y_CUT_BOT, :]
H, W = img_bez_paskow.shape
img_torus = np.zeros((H + 2 * delta, W + 2 * delta), dtype=img_bez_paskow.dtype)
img_torus[delta:delta + H, delta:delta + W] = img_bez_paskow
img_torus[delta:delta + H, delta + W:delta + W + delta] = img_bez_paskow[:, 0:delta]
img_torus[delta:delta + H, 0:delta] = img_bez_paskow[:, W - delta:W]
img_torus[delta + H:delta + H + delta, delta:delta + W] = img_bez_paskow[0:delta, :]
img_torus[0:delta, delta:delta + W] = img_bez_paskow[H - delta:H, :]
\end{lstlisting}
Cztery analogiczne instrukcje wypełniają rogi (np.\ \texttt{img\_torus[0:delta, 0:delta] = img\_bez\_paskow[H - delta:H, W - delta:W]}). Powstaje obraz $686\times688$ px, w którym klatka jest otoczona sąsiedztwem przeciwległych brzegów, czyli widziana tak, jak w strukturze okresowej (rozdz.~\ref{sec:def}, p.~\ref{ssec:lattice}). Dzięki temu bakterie leżące przy brzegu klatki są wykrywane z pełnym otoczeniem, a nie z brzegiem sztucznie uciętym.

\section{Detekcja: binaryzacja Otsu oraz metody LoG i DoG}
Detekcja odbywa się \emph{jeden raz} na całym torusie (wyniki są zapamiętywane w pamięci podręcznej \texttt{\_blobs\_cache}), a nie osobno w każdym oknie. Obraz jest najpierw binaryzowany progiem Otsu i oczyszczany otwarciem morfologicznym (Listing~\ref{lst:kod-otsu}):
\begin{lstlisting}[style=pykod,caption={Binaryzacja progiem Otsu i otwarcie morfologiczne},label={lst:kod-otsu}]
thresh_val, thresh = cv2.threshold(image, 0, 255, cv2.THRESH_BINARY_INV + cv2.THRESH_OTSU)
kernel = cv2.getStructuringElement(cv2.MORPH_CROSS, (3, 3))
result = cv2.morphologyEx(thresh, cv2.MORPH_OPEN, kernel, iterations=1)
\end{lstlisting}
Następnie środki bakterii wykrywają dwie metody wykrywania plam z biblioteki scikit-image: LoG (\emph{Laplacian of Gaussian}) i DoG (\emph{Difference of Gaussian}). Obie wywołuje się z tymi samymi progami, a promień plamy przelicza się z parametru skali $\sigma$ jako $r=\sigma\sqrt2$ (Listing~\ref{lst:kod-logdog}):
\begin{lstlisting}[style=pykod,caption={Wykrywanie środków bakterii metodami LoG i DoG},label={lst:kod-logdog}]
blobs_log = blob_log(img_denoised, min_sigma=BLOB_MIN_SIGMA, max_sigma=BLOB_MAX_SIGMA,
                     num_sigma=BLOB_NUM_SIGMA, threshold=BLOB_THRESHOLD)
blobs_dog = blob_dog(img_denoised, min_sigma=BLOB_MIN_SIGMA, max_sigma=BLOB_MAX_SIGMA,
                     threshold=BLOB_THRESHOLD)
blobs_log[:, 2] = blobs_log[:, 2] * sqrt(2)
blobs_dog[:, 2] = blobs_dog[:, 2] * sqrt(2)
\end{lstlisting}
Zostawia się tylko plamy o promieniu z przedziału $[r_{\min,\mathrm{det}},r_{\max,\mathrm{det}}]=[2,8]$ px, a wyniki obu metod łączy się w jeden zbiór środków (po prostym złączeniu tablic; Listing~\ref{lst:kod-join}):
\begin{lstlisting}[style=pykod,caption={Filtrowanie promieni plam i połączenie wyników LoG i DoG},label={lst:kod-join}]
mask = (blobs[:, 2] >= r_min) & (blobs[:, 2] <= r_max)
blobs_sum = np.concatenate([blobs_log_f[:, :2], blobs_dog_f[:, :2]], axis=0)
\end{lstlisting}
Metody LoG i DoG dają podobne, lecz nie identyczne zbiory środków; ta sama bakteria bywa więc wykryta dwukrotnie, a nawet kilkakrotnie. Takie powtórzenia usuwa następny krok.

\section{Odrzucanie punktów spoza klatki i punktów zbyt bliskich}
Punkty wykryte w doklejonych pasach nie należą do oryginalnej klatki, więc są odrzucane, a współrzędne wracają do układu klatki (Listing~\ref{lst:kod-mask}):
\begin{lstlisting}[style=pykod,caption={Odrzucenie punktów spoza oryginalnej klatki},label={lst:kod-mask}]
mask_oryginal = ((blobs_sum[:, 0] >= delta) & (blobs_sum[:, 0] < delta + H) &
                 (blobs_sum[:, 1] >= delta) & (blobs_sum[:, 1] < delta + W))
blobs_sum = blobs_sum[mask_oryginal]
blobs_sum[:, 0] -= delta
blobs_sum[:, 1] -= delta
\end{lstlisting}
Pozostałe punkty, które leżą od siebie bliżej niż $r_{\min}=8$ px, są scalane funkcją \texttt{clean\_kdtree\_torus}. Odległość mierzy się w metryce torusa: zbiór punktów kopiuje się w dziewięciu przesunięciach o okresy $(\pm y_d,\pm x_d)$, pary bliższe niż \texttt{dist} wyszukuje drzewo \texttt{cKDTree}, a punkty połączone w klastry (algorytm \emph{union--find}) zastępuje się średnią (Listing~\ref{lst:kod-clean}):
\begin{lstlisting}[style=pykod,caption={Scalanie punktów zbyt bliskich w metryce torusa (\texttt{clean\_kdtree\_torus})},label={lst:kod-clean}]
shifts = [(0, 0), (0, +x_d), (0, -x_d), (+y_d, 0), (-y_d, 0),
          (+y_d, +x_d), (+y_d, -x_d), (-y_d, +x_d), (-y_d, -x_d)]
kopia[:, 0] = y_0 + (kopia[:, 0] - y_0) % y_d
kopia[:, 1] = x_0 + (kopia[:, 1] - x_0) % x_d
tree = cKDTree(kopie)
pairs = tree.query_pairs(r=dist, output_type='ndarray')
oryginaly = [i for i in idx_lista if etykiety[i] == 0]
sr = kopie[oryginaly].mean(axis=0)
\end{lstlisting}
Do wyniku wchodzą tylko te klastry, które zawierają punkt oryginalny (etykieta $0$), więc para punktów po dwóch stronach brzegu klatki jest scalana poprawnie, a nie liczona podwójnie. Próg $r_{\min}=8$ px jest nieco większy od średnicy bakterii ($2r=7$ px).

\section{Okna pływające i baza położeń}
Okno o lewym górnym narożniku $(y_0,x_0)$ wycina z całej klatki punkty, które w nim leżą, i przesuwa je do układu okna $[0,100)\times[0,100)$. W oknie traktowanym jako torus usuwa się następnie duplikaty: jeśli dwa punkty są bliżej niż $r_{\min}$ w metryce torusa, usuwa się ten o większych współrzędnych $(y,x)$ w porządku leksykograficznym (zamiast uśredniać, co tworzyłoby sztuczne punkty; Listing~\ref{lst:kod-dup}):
\begin{lstlisting}[style=pykod,caption={Reguła usuwania duplikatów torusowych w oknie},label={lst:kod-dup}]
if (p_a[0] > p_b[0]) or (p_a[0] == p_b[0] and p_a[1] > p_b[1]):
    do_usuniecia[orig_a] = True
else:
    do_usuniecia[orig_b] = True
\end{lstlisting}
Potem punkty kwantyzuje się do siatki o kroku \texttt{DIST\_SCALANIE}$=10$ px (punkty w jednej komórce siatki zastępuje się średnią; Listing~\ref{lst:kod-quant}):
\begin{lstlisting}[style=pykod,caption={Kwantyzacja położeń do siatki o kroku \texttt{DIST\_SCALANIE}},label={lst:kod-quant}]
i_idx = np.floor((array[:, 0] - y_0) / krok).astype(np.int64)
j_idx = np.floor((array[:, 1] - x_0) / krok).astype(np.int64)
\end{lstlisting}
Wynikiem jest zbiór $N$ środków bakterii okna, który normalizuje się do kwadratu jednostkowego i zapisuje jako liczby zespolone $a_k=x_k+iy_k$ (Listing~\ref{lst:kod-norm}):
\begin{lstlisting}[style=pykod,caption={Normalizacja położeń do kwadratu jednostkowego i zapis w postaci zespolonej},label={lst:kod-norm}]
blobs_sum2[:, 0] = blobs_roi[:, 0] / y_d
blobs_sum2[:, 1] = blobs_roi[:, 1] / x_d
a = cp.asarray(array[:, 1] + 1j * array[:, 0], dtype=cp.complex128)
\end{lstlisting}
Gdy w oknie jest więcej niż \texttt{MAX\_PUNKTOW\_E2}$=500$ punktów, losuje się podzbiór (to ograniczenie nie jest w praktyce osiągane, bo okna zawierają ok.\ $25$--$50$ bakterii). Dla każdego okna zapisuje się wiersz: numer klatki, położenie okna $(y_0,x_0)$, $N$, $f$, sumy strukturalne, $\Re A^{(3)}$ i lepkość; tak powstają pliki \texttt{result\_lepkosc\_all.csv} (statystyki klatek) i tabele okien, na których opiera się rozdz.~\ref{sec:realstat}.

\oran{Poprawną kolejność potwierdza zgodność z Mathematicą (\texttt{Import-e2-CORRECTED.nb}): dla $182$ punktów $e_2=3{,}32561+0{,}12054\,i$ i $e_{2,2}=15{,}2837$.}

\section{Obliczanie sum strukturalnych i lepkości}
Sumy $e_2$ i $e_{2,2}$ (p.~\ref{ssec:struct}) liczy się z macierzy $\wp(a_k-a_m)$ funkcji Weierstrassa dla sieci kwadratowej ($g_2\approx189{,}0727$, $g_3=0$), a sumy $e_{33}$ i $e_4$ potrzebne w lepkości z funkcji $E_3=-\wp'/2$, $E_4=\wp^2-g_2/12$ i $E_3^{(1)}$ ze wzoru (A7) z Części II \cite{czapla2025computational}. Stałe sieci wylicza się z funkcji gamma (Listing~\ref{lst:kod-const}):
\begin{lstlisting}[style=pykod,caption={Stałe sieci kwadratowej: $S_4$, $g_2$, $S_3^{(1)}$},label={lst:kod-const}]
_G14  = math.gamma(0.25)
_S4   = _G14 ** 8 / (960 * math.pi ** 2)      # 3.151211
_G2   = 60.0 * _S4                             # 189.0727
_S31  = _G14 ** 8 / (384 * math.pi ** 3)      # S3^(1) = 2.50765
\end{lstlisting}
Funkcja \texttt{sumy\_e33\_e4} bierze różnice położeń w najbliższym obrazie torusowym i oblicza obie sumy (Listing~\ref{lst:kod-e33}):
\begin{lstlisting}[style=pykod,caption={Obliczanie sum $e_{33}$ i $e_4$ (\texttt{sumy\_e33\_e4})},label={lst:kod-e33}]
d = a[:, None] - a[None, :]
d = d - xp.round(d.real) - 1j * xp.round(d.imag)
E3 = -dwp / 2.0
E4 = wp ** 2 - _G2 / 12.0
E31 = (xp.conj(d) - zeta / math.pi) * E3 + E4 / math.pi - 5.0 * _S4 / math.pi + _S31
e33 = float(xp.sum(xp.abs(xp.sum(E31, axis=1)) ** 2)) / N ** 3
e4 = complex(xp.sum(E4)).conjugate() / N ** 3
\end{lstlisting}
Funkcje $\zeta$, $\wp$, $\wp'$ liczy się szeregiem theta (\texttt{wp\_zeta\_dp}), a przy zerowych różnicach przyjmuje się $E_3^{(1)}(0):=0$ i $E_4(0):=S_4$. \czer{[Do wyjaśnienia z promotorem: w Definicji~\ref{def:struct} jest $E_3^{(1)}(0):=S_3^{(1)}$; w obliczeniach $e_2$ i $e_{2,2}$ (funkcja \texttt{prepare}) kod przyjmuje $E_2=\wp+\pi$, $E_2(0):=\pi$, czyli $S_2=\pi$ (sumowanie Eisensteina), natomiast p.~\ref{ssec:visc} zakłada $S_2=0$. Dla $e_{33}$ i $e_4$ nie ma to znaczenia, bo $E_2$ w nich nie występuje.]} Następnie funkcja \texttt{lepkosc\_okna} wylicza stężenie, $\Re A^{(3)}$ i lepkość efektywną (Listing~\ref{lst:kod-A3}):
\begin{lstlisting}[style=pykod,caption={Obliczanie $\Re A^{(3)}$, stężenia $f$ i lepkości efektywnej (\texttt{lepkosc\_okna})},label={lst:kod-A3}]
A3 = 2.0 / math.pi ** 2 * (2.0 * e33 + 3.0 * e4r)
f = N * math.pi * (r_px / x_d) ** 2          # R_BAKT = 3.5
s = f + A3 * f ** 3
mu = (1.0 + s) / (1.0 - s) if s < 1.0 else float('nan')
\end{lstlisting}
Odpowiada to wzorom
\[
f=N\pi\Bigl(\frac{r}{\ell}\Bigr)^2,\qquad \Re A^{(3)}=\frac{2}{\pi^2}\bigl(2e_{33}+3\Re e_4\bigr),\qquad \frac{\eta_e}{\eta}=\frac{1+s}{1-s},\quad s=f+A^{(3)}f^3,
\]
gdzie $r=3{,}5$ px (stała \texttt{R\_BAKT}), $\ell=100$ px (p.~\ref{ssec:visc}). W kodzie $e_{33}$ jest liczone jako $\frac1{N^3}\sum_k\bigl|\sum_m E_3^{(1)}(a_k-a_m)\bigr|^2$\ (zgodnie ze wzorem \eqref{eq:ex2} po poprawce znaku).

\section{Wyszukiwanie anomalii i losowe zbiory odniesienia}
Pary okien anomalnych (większe $N$, mniejsza lepkość) wyszukuje funkcja \texttt{znajdz\_skrajne}. Okna jednej klatki grupuje się według liczby bakterii $N$ i szuka łańcuchów okien, w których $N$ rośnie o $1$, a lepkość maleje co najmniej o próg \texttt{d}; łańcuch znajduje się programowaniem dynamicznym (Listing~\ref{lst:kod-anom}):
\begin{lstlisting}[style=pykod,caption={Wyszukiwanie łańcuchów okien anomalnych (programowanie dynamiczne)},label={lst:kod-anom}]
krok = lev[k - 1]['mu_e_pade'].values[:, None] - lev[k]['mu_e_pade'].values[None, :]
wart = np.where(krok >= d, dp[:, None] + krok, -np.inf)
\end{lstlisting}
Pierwsze i ostatnie okno najlepszego łańcucha tworzą parę $A$, $B$. Jako zbiory odniesienia program losuje bakterie w oknie $100\times100$ px dla tych samych $N$ (funkcja \texttt{losuj\_punkty\_okna}), spełniające te same warunki co punkty po czyszczeniu: odległość torusowa każdych dwóch środków $\ge r_{\min}$ i co najwyżej jeden środek w komórce siatki $10\times10$ px. Dla tych zbiorów oblicza się te same wielkości funkcją \texttt{lepkosc\_okna}.

\section{Uwagi o implementacji i odtwarzalności}
Obliczenia macierzy $\wp(a_k-a_m)$ wykonuje się wektorowo (moduł pomocniczy \texttt{pfast}, biblioteka CuPy na GPU albo NumPy bez GPU), a wyniki detekcji i torusa zapamiętuje się w słownikach \texttt{\_torus\_cache} i \texttt{\_blobs\_cache}, więc każda klatka jest przetwarzana raz. Program wykorzystuje biblioteki OpenCV, scikit-image, SciPy, NumPy, pandas i mpmath. Kilka komentarzy w kodzie jest nieaktualnych: np.\ komentarz przy funkcji \texttt{stezenie\_f} mówi o $r=r_{\mathrm{px}}/220$, a rzeczywisty rozmiar okna to $100$ px (\texttt{x\_d}). Wartości liczbowe w rozprawie opierają się na kodzie, a nie na tych komentarzach.
}


\chapter{Protokoły generowania konfiguracji i wyniki symulacji}\label{sec:sim}

Wszystkie trzy protokoły generują $N$ dysków o stałym promieniu $r$ na torusie, teoretycznie zgodnie z rozkładem jednostajnym nienakładających się dysków, lecz różnią się sposobem losowania i akceptacji środków \cite{czapla2025computational}:
\begin{itemize}[nosep]
\item \textbf{R} (random sequential adsorption, RSA): kolejne środki losuje się jednostajnie na $\T^2$; kandydat $a_k$ jest odrzucany, jeśli $\|a_k-a_j\|<2r$ dla pewnego $j<k$, i losowany ponownie;
\item \textbf{T} (czyszczenie): losuje się punkty, a następnie iteracyjnie zastępuje punkty zbyt bliskie ich środkiem geometrycznym, aż wszystkie pary spełnią próg minimalnej odległości $d_{\min}$ dobrany do docelowego $f$;
\item \textbf{P} (osadzanie): zaczyna się od konfiguracji z nakładaniem i iteracyjnie rozsuwa nakładające się dyski ($a_i\mapsto a_i-\tfrac12\delta$, $a_j\mapsto a_j+\tfrac12\delta$), z zawijaniem na torus, aż do zaniku nakładań; metoda osiąga większe gęstości niż R.
\end{itemize}
Symulacje wykonano dla $N=100$, $f\in\{0{,}1;0{,}2;0{,}3;0{,}4\}$, po 100 realizacji dla każdej pary (protokół, $f$); łącznie $100\cdot3\cdot4=1200$ konfiguracji, a sumy strukturalne uśredniono po realizacjach \cite{czapla2025computational,mityushev2025computationalpart3}. Podstawienie do \eqref{eq:mue31} dało wielomiany stopnia 4 względem $f$ dla $A_J(\varrho,f)$, $J\in\{R,T,P\}$, dla każdego $f$ osobno (wzory (56)--(67) w \cite{czapla2025computational}); zależność od $f$ jest też pośrednio zawarta w liczbowych współczynnikach, bo konfiguracje generowane przy różnych $f$ mają różne sumy strukturalne.

Dla twardych inkluzji $\mu_J(1,f)=(1+A_J)/(1-A_J)$, gdzie $A_J(1,f)=f+a_2f^2+a_3f^3+a_4f^4$. Współczynniki z \cite{czapla2025computational} (wzory (69)--(80)) zebrano w tabeli~\ref{tab:coef}. Wartość $a_2$ powinna być równa zeru (izotropia, $A^{(2)}=0$); jej odchylenie od zera mierzy odstępstwo uśrednionej konfiguracji od idealnej izotropii.

\begin{table}[h]
\centering\small
\caption{Współczynniki $a_2,a_3,a_4$ wielomianu $A_J(1,f)$ dla twardych inkluzji \cite{czapla2025computational}.}\label{tab:coef}
\begin{tabular}{@{}cccc|ccc|ccc@{}}
\toprule
 & \multicolumn{3}{c|}{R} & \multicolumn{3}{c|}{T} & \multicolumn{3}{c}{P}\\
$f$ & $a_2$ & $a_3$ & $a_4$ & $a_2$ & $a_3$ & $a_4$ & $a_2$ & $a_3$ & $a_4$\\ \midrule
0,1 & 0,0038 & 8,9608 & $-9{,}4832$ & $-0{,}0224$ & 17,1591 & $-24{,}9094$ & 0,0511 & 16,902 & $-22{,}9362$\\
0,2 & $-0{,}0145$ & 4,2253 & $-0{,}5876$ & $-0{,}0130$ & 7,1840 & $-0{,}8448$ & 0,0018 & 7,6789 & $-0{,}5701$\\
0,3 & 0,0215 & 2,5687 & 0,2891 & $-0{,}0189$ & 4,5616 & $-0{,}1218$ & $-0{,}0048$ & 4,5440 & 0,5868\\
0,4 & 0,0138 & 1,7336 & 0,6695 & $-0{,}0503$ & 3,1748 & 0,4207 & $-0{,}0092$ & 2,9193 & 0,9921\\ \bottomrule
\end{tabular}
\end{table}


\chapter{Resumacja szeregów i wykładnik krytyczny}\label{sec:resum}

Dla twardych inkluzji w nieściśliwej matrycy ($\kappa=\varrho_i=1$; po zamianie $\mu\to\eta$, $\nu=\tfrac12$, jest to efektywna lepkość zawiesiny sztywnych dysków, p.~\ref{ssec:visc}) wielkość $\mu_e$ rozbiega się do nieskończoności przy zbliżaniu się do granicznego stężenia $f_c=\pi/(2\sqrt3)\approx0{,}9069$ (progu perkolacji kontaktujących się dysków):
\begin{equation}\label{eq:crit}
\mu_e\simeq B\,(f_c-f)^{-s},\qquad f\to f_c^-,
\end{equation}
gdzie $B$ jest amplitudą krytyczną, a $s$ wykładnikiem krytycznym; oczekiwana wartość to $s\approx1{,}3$ \cite{czapla2025computational}. Zależność chcemy opisać w całym przedziale $0\le f<f_c$, a znamy tylko kilka pierwszych współczynników małego parametru $f$. Rozumowanie przebiega w pięciu krokach: szereg, wielomian, aproksymanta Padégo, parametry $(f_c,s)$ przez pochodną logarytmiczną i wreszcie metoda Gluzmana.

\section{\texorpdfstring{Krok 1: szereg}{Krok 1: szereg}}

Z \eqref{eq:mue31} otrzymujemy $\mu_e/\mu=\sum_{n\ge0}c^{(n)}f^n$. Dla twardych inkluzji $\mu_J=(1+A_J)/(1-A_J)$, $A_J=f+a_2f^2+a_3f^3+a_4f^4$ (rozdz.~\ref{sec:sim}), więc
\begin{equation*}
{c^{(0)}=1,\quad c^{(1)}=2,\quad c^{(2)}=2(1+a_2),\quad c^{(3)}=2(1+2a_2+a_3),\quad c^{(4)}=2(1+3a_2+a_2^2+2a_3+a_4).}
\end{equation*}
Współczynniki do $n=4$ wymagają sum strukturalnych o wadze do $6$ (tabela w p.~\ref{ssec:struct}); kolejne wymagałyby sum o coraz wyższych wagach.

\section{\texorpdfstring{Krok 2: wielomian (czy faktycznie?)}{Krok 2: wielomian (czy faktycznie?)}}

Obcięcie szeregu po $n=4$ daje
\begin{equation*}
{\mu_{e,4}(f)=\sum_{n=0}^{4}c^{(n)}f^n.}
\end{equation*}
Jest to wielomian stopnia $4$ względem $f$, ale o współczynnikach \emph{numerycznych}: dla danego protokołu i stężenia $c^{(n)}$ są liczbami obliczonymi z sum strukturalnych, a sumy te zależą od konfiguracji, czyli od $f$. Dlatego $\mu_{e,4}$ jest wielomianem numerycznym dla $\mu_e$ (rozdz.~\ref{sec:sim}), a nie rozwinięciem Taylora w ścisłym sensie. Wielomian jest ograniczony na przedziale skończonym, a $\mu_e$ rozbiega się przy $f\to f_c$; np.\ wielomian $\mu^{\{9\}}_{e,4}$ podany niżej daje w $f=f_c$ wartość około $31$. Zatem blisko $f_c$ jest on nieprzydatny.

\section{\texorpdfstring{Krok 3: aproksymacja Padégo, czyli aproksymanta wymierna}{Krok 3: aproksymacja Padégo, czyli aproksymanta wymierna}}

Funkcję wymierną, która może mieć biegun, daje aproksymanta Padégo $[L/M]$: iloraz wielomianów $P_L(f)/Q_M(f)$, $Q_M(0)=1$, którego rozwinięcie zgadza się z szeregiem do wyrazu $f^{L+M}$. Dla danych $c^{(0)},\dots,c^{(4)}$ mamy $L+M=4$. Na przykład $[3/1]$ ma $Q_1=1-\frac{c^{(4)}}{c^{(3)}}f$, a więc biegun w $f_p=c^{(3)}/c^{(4)}$; dla wielomianu $\mu^{\{9\}}_{e,4}$ daje to $f_p\approx0{,}959$ (własne obliczenie ze współczynników podanych niżej), zamiast $f_c\approx0{,}907$. Aproksymanta wymierna ma dwie wady w tym zadaniu: położenie bieguna wynika z niewielu współczynników, a sam biegun jest prosty, czyli odpowiada wykładnikowi $1$, podczas gdy w \eqref{eq:crit} oczekujemy $s\approx1{,}3$. Osobliwość $(f_c-f)^{-s}$ z niecałkowitym $s$ jest punktem rozgałęzienia, a nie biegunem.

\section{\texorpdfstring{Krok 4: parametry $(f_c,s)$ i pochodna logarytmiczna}{Krok 4: parametry (f\_c,s) i pochodna logarytmiczna}}

Wartość $f_c$ znamy z geometrii (Dodatek, pkt~7), więc nieznanymi parametrami są $s$ i $B$. Aby uzyskać $s$ jako wielkość wprost, przechodzimy do funkcji z logarytmem. Jeśli $\mu_e\simeq B(f_c-f)^{-s}$, to
\begin{equation*}
{D(f)=\frac{d}{df}\ln\mu_e=\frac{s}{f_c-f}+\dots,}
\end{equation*}
czyli $D$ ma \emph{prosty} biegun w $f_c$ z residuum $s$ (z minus). Szereg $D$ liczymy ze współczynników: $D(f)=d_0+d_1f+d_2f^2+\dots$, $d_0=c^{(1)}$, $d_1=2c^{(2)}-(c^{(1)})^2$, $d_2=3c^{(3)}-3c^{(1)}c^{(2)}+(c^{(1)})^3$ (dla $c^{(0)}=1$). Aproksymanta Padégo funkcji $D$ daje wtedy parę $(f_c,s)$: położenie i residuum bieguna (tzw.\ metoda $D$-log Padé). Równoważnie można podnieść funkcję do potęgi $1/s$: $\mu_e^{1/s}\simeq B^{1/s}(f_c-f)^{-1}$ ma biegun prosty, dla którego aproksymanty wymierne są odpowiednie. Ten drugi pomysł jest podstawą metody opisanej niżej.

\section{\texorpdfstring{Krok 5: iterowane pierwiastki (Gluzman)}{Krok 5: iterowane pierwiastki (Gluzman)}}

Metodę, w której wykładnik $s$ jest wpisany w postać aproksymanty, rozwinął Gluzman: aproksymanty iterowanych pierwiastków z transformacją typu Borela \cite{gluzman2022iterative} oraz wybór $s$ z warunków na wykładniki krytyczne \cite{Gluzman2024CriticalIndices}. Szereg obcięty przekształca się transformacją typu Borela
\begin{equation}\label{eq:borel}
\psi_{k,b}(f)=\sum_{n=0}^{k}\frac{c^{(n)}}{(\Gamma(1+n))^{b}}f^n,\qquad b\in\{0,1\},\ k=4,
\end{equation}
($b=0$ --- szereg wyjściowy, $b=1$ --- klasyczna transformacja Borela), a następnie aproksymantę iterowanych pierwiastków \cite{gluzman2022iterative}
\begin{equation}\label{eq:root}
\psi^*_{k,b}(f)=\Bigl(\Bigl(\Bigl((1+A_{1,b}z)^2+A_{2,b}z^2\Bigr)^{3/2}+A_{3,b}z^3\Bigr)^{4/3}+\dots+A_{k,b}z^k\Bigr)^{s_{k,b}/k},\qquad z=\frac{f}{f_c-f},
\end{equation}
której „amplitudy” $A_{n,b}$ wyznacza się z dopasowania asymptotycznego do \eqref{eq:borel} przy $f\to0$. Przy $f\to f_c$ zachodzi $\psi^*_{k,b}\simeq f_c^{s}(f_c-f)^{-s}C_{k,b}(s)$, co daje amplitudę $B_{k,b}(s)=f_c^{s}C_{k,b}(s)(\Gamma(1+s))^b$. Wykładnik $s$ wyznacza się z jednego z dwóch warunków \cite{Gluzman2024CriticalIndices}:
\begin{align}
&\text{minimalnej różnicy amplitud:}&&B_{4,0}(s)-B_{3,0}(s)=0,\label{eq:md}\\
&\text{minimalnej różnicy czułości:}&&\frac{\partial B_{4,1}(s)}{\partial s}-\frac{\partial B_{4,0}(s)}{\partial s}=0.\label{eq:msd}
\end{align}

\paragraph{Wyniki.} Dla poszczególnych protokołów wykładnik $s$ wynosi:
\begin{center}\small
\begin{tabular}{@{}lccc@{}}
\toprule
Warunek & R & P & T\\ \midrule
minimalna różnica \eqref{eq:md} & 1,50628 & 1,43399 & 1,39844\\
minimalna różnica czułości \eqref{eq:msd} & 1,25551 & 1,29169 & 1,26844\\ \bottomrule
\end{tabular}
\end{center}
Uśrednianie wielomianów po protokołach i stężeniach daje cztery zestawy: $M=12$ (wszystkie), $M=9$ (bez $f=0{,}4$), $M=6$ (bez $f=0{,}3$ i $0{,}4$), $M=3$ (tylko $f=0{,}1$). Jako kryterium wyboru najlepszej aproksymacji przyjęto minimalne odchylenie $c^{(2)}$ od wartości teoretycznej $2$:
\begin{center}\small
\begin{tabular}{@{}lcccc@{}}
\toprule
 & $M=12$ & $M=9$ & $M=6$ & $M=3$\\ \midrule
$c^{(2)}$ (dokładnie 2) & 1,99315 & 2,00103 & 2,00228 & 2,02171\\
$s$, minimalna różnica \eqref{eq:md} & 1,43417 & 1,40231 & 1,35931 & 1,28867\\
$s$, minimalna różnica czułości \eqref{eq:msd} & 1,26628 & 1,27434 & 1,28509 & 1,29518\\ \bottomrule
\end{tabular}
\end{center}
Kryterium wskazuje $M=9$ z wielomianem $\mu^{\{9\}}_{e,4}(f)=1+2f+2{,}00103f^2+18{,}3986f^3+19{,}1828f^4$ i przedziałem $s\in[1{,}27434;\,1{,}40231]$. Dla $s=\tfrac12(s_{\min}+s_{\max})=1{,}33833$ otrzymano zwarty wzór ważny dla wszystkich $f<f_c$ (wzór (89) w \cite{czapla2025computational}):
\begin{equation}\label{eq:mustar}
\mu_e^*(f)=\Bigl(6{,}6104\,z^4+\Bigl(26{,}55\,z^3+\bigl((1+1{,}3553\,z)^2-0{,}8725\,z^2\bigr)^{3/2}\Bigr)^{4/3}\Bigr)^{0{,}3346},\quad z=\frac{f}{0{,}9069-f},
\end{equation}
z amplitudą krytyczną $B\approx3{,}9486$. Zgodność z szeregiem wyjściowym dla małych $f$ wynika stąd, że $\mu_e^*\approx1+s\,A_1 z\approx1+2f$ (tj.\ $c^{(1)}=2$). Wzór \eqref{eq:mustar} odbiega od $\mu^{\{9\}}_{e,4}$ już około $f\approx0{,}4$, gdzie zaczyna się reżim krytyczny. Bez stosowania kryterium wyboru można przyjąć $s\approx1{,}292$ jako wynik uśrednienia trzech protokołów dla małych stężeń (kolumna $M=3$) \cite{czapla2025computational}.


\chapter{Sumy strukturalne w uczeniu maszynowym}\label{sec:ml}

\section{Cel i dane}

Część III \cite{mityushev2025computationalpart3} wykorzystuje sumy strukturalne jako opis mikrostruktury. Kompozyty uzyskane różnymi protokołami (R, T, P) bywają praktycznie nierozróżnialne wizualnie, wg korelacji dwupunktowej i metodami czysto numerycznymi; sumy strukturalne zawierają informację o korelacjach wyższych rzędów i pozwalają klasyfikować mikrostruktury \emph{bez} obliczania stałych efektywnych. Zbiór danych to 1200 konfiguracji ($3$ protokoły $\times$ $4$ stężenia $\times$ $100$ realizacji, $N=100$) z etykietą protokołu i stężeniem $f$.

\section{Konstrukcja cech}

Dla każdej z 14 sum $z\in E_{14}(N)$ liczy się cztery deskryptory: $\Re z$, $\Im z$, $|z|$ i $\arg z$ (z korektą kwadrantu), a następnie agregaty rodzinowe (średnia i odchylenie standardowe $\Re z$, $|z|$, $\arg z$ w grupach sum o wspólnej strukturze indeksów) oraz wielkości pochodne (stosunki części rzeczywistej do urojonej, znormalizowane różnice amplitud). Daje to około 130 cech, standaryzowanych (z-score). Analiza PCA wskazuje około 36 efektywnie niezależnych składowych (współliniowość cech); PCA służy wyłącznie wizualizacji, a klasyfikator działa na pełnym zbiorze cech \cite{mityushev2025computationalpart3}.

\section{Klasyfikacja R/T/P}

Klasyfikatorem jest zespół typu bagging złożony z 200 nieprzycinanych drzew decyzyjnych (MATLAB \texttt{fitcensemble}, metoda \texttt{Bag}); walidacja: stratyfikowana 5-krotna. Średnia dokładność wynosi $\mathrm{ACC}=0{,}7125$ przy poziomie losowym $\approx0{,}33$. Macierz pomyłek (rys.~2 w \cite{mityushev2025computationalpart3}):
\begin{center}\small
\begin{tabular}{@{}l|ccc|c@{}}
\toprule
prawdziwa $\backslash$ przewidziana & P & R & T & czułość\\ \midrule
P & 271 & 18 & 111 & 0,678\\
R & 23 & 354 & 23 & 0,885\\
T & 136 & 34 & 230 & 0,575\\ \bottomrule
\end{tabular}
\end{center}
Suma przekątnej, $855/1200=0{,}7125$, zgadza się z podaną dokładnością. Z macierzy wynika, że dominują pomyłki P$\leftrightarrow$T (247 z 345 błędów), a klasa R jest rozpoznawana najlepiej, co jest sprzeczne z komentarzem w tekście źródłowym (Dodatek~\ref{app:notes}). Ważność cech liczono metodą permutacyjną (\texttt{oobPermutedPredictorImportance}): najważniejsze są moduły $|z|$ sum wyższych rzędów i części rzeczywiste $\Re z$. Testy odporności (szum gaussowski $\sigma=5\%$ odchylenia standardowego cech, podwojenie liczby realizacji, zmiana promieni o $\pm5\%$) zmieniały ACC o mniej niż $0{,}01$--$0{,}015$. Projekcje PCA i t-SNE (rys.~3--4 w \cite{mityushev2025computationalpart3}) pokazują częściową separację klas; w przestrzeni dwuwymiarowej najlepsze pary cech wybrano klasyfikatorem $k$-NN ($k=5$) przez przeszukanie wszystkich par.

\section{Miary nieregularności}

Dla \emph{zbioru} kompozytów wprowadzono indeks
\begin{equation}\label{eq:eta}
\eta^*=0{,}5\,z_E+0{,}3\,z_\phi+0{,}2\,z_A,
\end{equation}
gdzie $z_E$ (energia wyższych rzędów) jest średnią $|z|$ po rodzinach o wyższych wskaźnikach, normalizowaną min--max, $z_\phi$ (chaos fazowy) odchyleniem standardowym $\arg z$, normalizowanym przez $\pi$, a
\begin{equation}
z_A=\frac{\|\Im z\|_2}{\|\Re z\|_2}
\end{equation}
(asymetria Im/Re) po przekształceniu $x\mapsto x/(1+x)$ i normalizacji min--max. Wagi dobrano na podstawie wariancji i mocy dyskryminacyjnej; alternatywne schematy (wagi równe, proporcjonalne do wariancji, entropijne) dają wartości silnie skorelowane z \eqref{eq:eta} ($r>0{,}97$). Średnia $\eta^*$ maleje ze stężeniem: $\approx0{,}38$ ($f=0{,}1$), $0{,}24$ ($0{,}2$), $0{,}22$ ($0{,}3$), $0{,}20$ ($0{,}4$); średnio $0{,}31$ dla $f\in[0{,}1;0{,}2]$ i $0{,}21$ dla $f\in[0{,}3;0{,}4]$, przy równych liczebnościach klas (po 200 w każdym przedziale), co wyklucza wpływ zmiany proporcji klas.

Dla \emph{pojedynczej} mikrostruktury zdefiniowano miarę strukturalnej nieregularności względem sieci heksagonalnej:
\begin{equation}\label{eq:chi}
\chi=\sum_{\substack{i=1\\E_i(N)\ne0}}^{M}f^{\ell(i)}\Bigl|1-\frac{E_i^{\rm hex}}{E_i(N)}\Bigr|\ln\Bigl|1-\frac{E_i^{\rm hex}}{E_i(N)}\Bigr|,
\end{equation}
gdzie $E_i^{\rm hex}$ to sumy strukturalne regularnej sieci heksagonalnej ($N=1$), a $\ell(i)$ najmniejszy rząd (potęga $f$) współczynnika wielomianu przybliżającego stałą efektywną, w którym występuje suma $E_i$. Miara $\chi$ znika dla sieci heksagonalnej.

\textbf{Wyniki części III \cite{mityushev2025computationalpart3}.} Poniżej zebrano wyniki dotyczące nieregularności; to, co podają autorzy \cite{mityushev2025computationalpart3}, odróżniono od uwag własnych.
\begin{itemize}[nosep]
\item Obie miary mają różne zadania: indeks $\eta^*$ \eqref{eq:eta} jest skalarnym deskryptorem nieregularności \emph{zbioru} kompozytów, a miara $\chi$ \eqref{eq:chi} opisuje pojedynczą mikrostrukturę.
\item Dla 1200 realizacji podano histogram $\eta^*$ (rys.~6 w \cite{mityushev2025computationalpart3}) i zależność od stężenia (rys.~7 w \cite{mityushev2025computationalpart3}: średnia i odchylenie standardowe na każdym poziomie $f$). Średnia maleje ze wzrostem $f$ (wartości podano wyżej), a różnica między przedziałami $f\in[0{,}1;0{,}2]$ i $[0{,}3;0{,}4]$ ($0{,}31$ i $0{,}21$) nie wynika ze zmiany liczebności klas.
\item Miara $\chi$ znika dla regularnej sieci heksagonalnej i mierzy ilościowo odchylenie struktury nieuporządkowanej od sieci regularnej. Według \cite{mityushev2025computationalpart3} jest pojęciowo zbliżona do wcześniejszej miary typu entropijnego, wyprowadzonej jednak dla innych klas kompozytów.
\item Druga składowa główna w analizie PCA (rys.~3 w \cite{mityushev2025computationalpart3}) odzwierciedla według autorów wtórną zmienność związaną z nieregularnością strukturalną.
\end{itemize}
Uwagi własne. (a) Oznaczenie $\eta^*$ pochodzi z \cite{mityushev2025computationalpart3} i nie ma związku z lepkością $\eta$ z p.~\ref{ssec:visc}. (b) Miary $\eta^*$ i $\chi$ zależą wyłącznie od położeń środków, więc nie zależą od tego, czy stałą efektywną jest $\mu_e$ czy $\eta_e$ (twierdzenie o dekompozycji). (c) Sieć heksagonalna jest konfiguracją o najgęstszym upakowaniu, $f=f_c=\pi/(2\sqrt3)$ (rozdz.~\ref{sec:resum}), więc $\chi$ można czytać jako odległość od struktury granicznej. (d) Składnik $|u|\ln|u|$ w \eqref{eq:chi} ma postać wkładu entropijnego. (e) W \cite{czapla2025computational,mityushev2025computationalpart3} nie badano związku $\chi$ z wykładnikiem krytycznym $s$, zależnym od protokołu (rozdz.~\ref{sec:resum}); sprawdzenie takiej zależności byłoby naturalnym rozszerzeniem, ale pozostaje hipotezą.


\chapter{\texorpdfstring{Statystyczna obróbka obserwowanych struktur: statystyka po klatkach i anomalie}{Statystyczna obróbka obserwowanych struktur: statystyka po klatkach i anomalie}}\label{sec:realstat}

W rozdziale przedstawiono statystykę wyników po wszystkich klatkach filmu, przykłady anomalii i wnioski z nich. Dane obserwowane (rozdz.~\ref{sec:real}) odnosi się do zbiorów teoretycznych z rozdz.~\ref{sec:sim} i \ref{sec:ml} oraz do teorii z rozdz.~\ref{sec:def}--\ref{sec:resum}.

\section{\texorpdfstring{Statystyka po wszystkich klatkach}{Statystyka po wszystkich klatkach}}

Dla każdej z 239 klatek porównano wszystkie pary okien o różnym stężeniu (średnio $2{,}87\cdot10^6$ par na klatkę). Parę $(A,B)$ nazwano \emph{dyskordantną}, jeśli $f_A<f_B$, ale $\eta_A/\eta>\eta_B/\eta$, tzn.\ okno o mniejszym stężeniu ma większą lepkość efektywną. Pary o jednakowym $N$ pominięto, bo wtedy różnica $\eta_e/\eta$ wynika wyłącznie z układu bakterii. Udział par dyskordantnych wynosi średnio $2{,}43\,\%$ (od $1{,}55\,\%$ w klatce 230 do $3{,}70\,\%$ w klatce 106), co daje średnio około $69\,900$ par na klatkę. Współczynnik Spearmana między $f$ i $\eta_e/\eta$ wynosi średnio $0{,}9825$ (od $0{,}9706$ w klatce 233 do $0{,}9903$ w klatce 115). Zależność jest więc silnie monotoniczna, ale nie ścisła: w każdej klatce zdarzają się pary okien, w których większe stężenie idzie w parze z mniejszą lepkością. Pięć klatek z największym udziałem takich par (31, 1, 105, 233 i 106) zestawiono w tabeli~\ref{tab:realstat}.
\begin{table}[htbp]
\centering\small
\caption{Statystyka par okien w klatkach z największym udziałem par dyskordantnych (2500 okien w klatce); $\rho_S$ --- współczynnik korelacji Spearmana między $f$ i $\eta_e/\eta$.}\label{tab:realstat}
\begin{tabular}{@{}ccccc@{}}
\toprule
\textbf{$i$} & \textbf{pary o różnym $f$} & \textbf{pary dyskordantne} & \textbf{udział} & \textbf{$\rho_S(f,\eta_e/\eta)$}\\
\midrule
$31$ & $2\,861\,419$ & $97\,068$ & $3{,}39\,\%$ & $0{,}9745$\\
$1$ & $2\,867\,989$ & $97\,590$ & $3{,}40\,\%$ & $0{,}9744$\\
$105$ & $2\,869\,995$ & $99\,046$ & $3{,}45\,\%$ & $0{,}9760$\\
$233$ & $2\,834\,522$ & $99\,351$ & $3{,}51\,\%$ & $0{,}9706$\\
$106$ & $2\,891\,089$ & $106\,962$ & $3{,}70\,\%$ & $0{,}9758$\\
\bottomrule
\end{tabular}
\end{table}

\section{\texorpdfstring{Przykłady anomalii}{Przykłady anomalii}}

W klatkach z największym udziałem par dyskordantnych szukano \emph{łańcuchów} okien o kolejnych liczbach bakterii $N,N+1,\dots$, w których stężenie rośnie, a lepkość efektywna maleje o co najmniej zadany próg w każdym kroku (progi $0{,}004;0{,}003;\dots;0{,}0001$; wybiera się największy możliwy próg, a dla niego łańcuch o największym łącznym spadku). Skrajne okna łańcucha, $A$ i $B$, tworzą najsilniejszą anomalię: okno $B$ ma więcej bakterii, ale mniejszą lepkość niż okno $A$. Przykłady dla klatek 31, 1, 105, 233 i 106 pokazują rys.~\ref{fig:real31}--\ref{fig:real106} i tabele~\ref{tab:real31}--\ref{tab:real106}.

W klatce $i=31$ okno $B$ zawiera więcej bakterii niż okno $A$ ($N_B=42>N_A=38$, stężenie większe o $10{,}5$\,\%), a mimo to ma mniejszą lepkość efektywną: $\eta_B/\eta=1{,}4019<\eta_A/\eta=1{,}4107$ (różnica $0{,}63$\,\%). Przyczyną jest układ środków bakterii: współczynnik $\Re A^{(3)}$ jest w oknie $A$ około $5{,}7$ razy większy niż w oknie $B$ ($7{,}7123$ wobec $1{,}3443$), więc człon $\Re A^{(3)}f^3$ w \eqref{eq:etaser} przeważa nad przyrostem stężenia.

\begin{figure}[htbp]
\centering
\includegraphics[width=0.75\textwidth]{figs/real31_2D_N_interp_anomalia.png}
\caption{Klatka $i=31$: mapa liczby bakterii $N$ w oknach pływających o wymiarach $100\times100$ px z zaznaczonymi oknami $A$ i $B$ --- krańcami łańcucha pięciu okien o kolejnych liczbach $N$ (próg kroku 0,001). Etykiety okien podają położenie narożnika okna $(x_0,y_0)$ w pikselach, liczbę bakterii $N$, stężenie $f$, współczynnik $\Re A^{(3)}$ oraz lepkość efektywną, oznaczoną na rysunku symbolem $\mu_e/\mu$ (w tekście: $\eta_e/\eta$).}
\label{fig:real31}
\end{figure}

\begin{table}[htbp]
\centering\small
\caption{Wyniki dla skrajnych okien $A$ i $B$ w klatce $i=31$ ($r=3{,}5$ px; okno $A$: $(x_0,y_0)=(396,360)$, okno $B$: $(517,430)$).}\label{tab:real31}
\begin{tabular}{@{}ccccccc@{}}
\toprule
\textbf{Okno} & \textbf{$N$} & \textbf{$f$} & \textbf{$e_{33}$} & \textbf{$e_4'$} & \textbf{$\Re A^{(3)}$} & \textbf{$\eta_e/\eta$}\\
\midrule
A & $38$ & $0{,}1462$ & $12{,}3807$ & $4{,}4324$ & $7{,}7123$ & $1{,}4107$\\
B & $42$ & $0{,}1616$ & $4{,}0522$ & $-0{,}4902$ & $1{,}3443$ & $1{,}4019$\\
\bottomrule
\end{tabular}
\end{table}

W klatce $i=1$ okno $B$ zawiera więcej bakterii niż okno $A$ ($N_B=42>N_A=38$, stężenie większe o $10{,}5$\,\%), a mimo to ma mniejszą lepkość efektywną: $\eta_B/\eta=1{,}3960<\eta_A/\eta=1{,}4256$ (różnica $2{,}08$\,\%). Przyczyną jest układ środków bakterii: współczynnik $\Re A^{(3)}$ jest w oknie $A$ około $10{,}8$ razy większy niż w oknie $B$ ($9{,}3470$ wobec $0{,}8657$), więc człon $\Re A^{(3)}f^3$ w \eqref{eq:etaser} przeważa nad przyrostem stężenia.

\begin{figure}[htbp]
\centering
\includegraphics[width=0.75\textwidth]{figs/real1_2D_N_interp_anomalia.png}
\caption{Klatka $i=1$: mapa liczby bakterii $N$ w oknach pływających o wymiarach $100\times100$ px z zaznaczonymi oknami $A$ i $B$ --- krańcami łańcucha pięciu okien o kolejnych liczbach $N$ (próg kroku 0,004). Etykiety okien podają położenie narożnika okna $(x_0,y_0)$ w pikselach, liczbę bakterii $N$, stężenie $f$, współczynnik $\Re A^{(3)}$ oraz lepkość efektywną, oznaczoną na rysunku symbolem $\mu_e/\mu$ (w tekście: $\eta_e/\eta$).}
\label{fig:real1}
\end{figure}

\begin{table}[htbp]
\centering\small
\caption{Wyniki dla skrajnych okien $A$ i $B$ w klatce $i=1$ ($r=3{,}5$ px; okno $A$: $(x_0,y_0)=(506,350)$, okno $B$: $(121,160)$).}\label{tab:real1}
\begin{tabular}{@{}ccccccc@{}}
\toprule
\textbf{Okno} & \textbf{$N$} & \textbf{$f$} & \textbf{$e_{33}$} & \textbf{$e_4'$} & \textbf{$\Re A^{(3)}$} & \textbf{$\eta_e/\eta$}\\
\midrule
A & $38$ & $0{,}1462$ & $16{,}0469$ & $4{,}6773$ & $9{,}3470$ & $1{,}4256$\\
B & $42$ & $0{,}1616$ & $4{,}6459$ & $-1{,}6732$ & $0{,}8657$ & $1{,}3960$\\
\bottomrule
\end{tabular}
\end{table}

W klatce $i=105$ okno $B$ zawiera więcej bakterii niż okno $A$ ($N_B=42>N_A=38$, stężenie większe o $10{,}5$\,\%), a mimo to ma mniejszą lepkość efektywną: $\eta_B/\eta=1{,}3967<\eta_A/\eta=1{,}4055$ (różnica $0{,}63$\,\%). Przyczyną jest układ środków bakterii: współczynnik $\Re A^{(3)}$ jest w oknie $A$ około $7{,}8$ razy większy niż w oknie $B$ ($7{,}1384$ wobec $0{,}9164$), więc człon $\Re A^{(3)}f^3$ w \eqref{eq:etaser} przeważa nad przyrostem stężenia.

\begin{figure}[htbp]
\centering
\includegraphics[width=0.75\textwidth]{figs/real105_2D_N_interp_anomalia.png}
\caption{Klatka $i=105$: mapa liczby bakterii $N$ w oknach pływających o wymiarach $100\times100$ px z zaznaczonymi oknami $A$ i $B$ --- krańcami łańcucha pięciu okien o kolejnych liczbach $N$ (próg kroku 0,001). Etykiety okien podają położenie narożnika okna $(x_0,y_0)$ w pikselach, liczbę bakterii $N$, stężenie $f$, współczynnik $\Re A^{(3)}$ oraz lepkość efektywną, oznaczoną na rysunku symbolem $\mu_e/\mu$ (w tekście: $\eta_e/\eta$).}
\label{fig:real105}
\end{figure}

\begin{table}[htbp]
\centering\small
\caption{Wyniki dla skrajnych okien $A$ i $B$ w klatce $i=105$ ($r=3{,}5$ px; okno $A$: $(x_0,y_0)=(363,480)$, okno $B$: $(176,420)$).}\label{tab:real105}
\begin{tabular}{@{}ccccccc@{}}
\toprule
\textbf{Okno} & \textbf{$N$} & \textbf{$f$} & \textbf{$e_{33}$} & \textbf{$e_4'$} & \textbf{$\Re A^{(3)}$} & \textbf{$\eta_e/\eta$}\\
\midrule
A & $38$ & $0{,}1462$ & $12{,}8699$ & $3{,}1623$ & $7{,}1384$ & $1{,}4055$\\
B & $42$ & $0{,}1616$ & $5{,}5752$ & $-2{,}2094$ & $0{,}9164$ & $1{,}3967$\\
\bottomrule
\end{tabular}
\end{table}

W klatce $i=233$ okno $B$ zawiera więcej bakterii niż okno $A$ ($N_B=43>N_A=39$, stężenie większe o $10{,}3$\,\%), a mimo to ma mniejszą lepkość efektywną: $\eta_B/\eta=1{,}4100<\eta_A/\eta=1{,}4190$ (różnica $0{,}64$\,\%). Przyczyną jest układ środków bakterii: współczynnik $\Re A^{(3)}$ jest w oknie $A$ około $6{,}7$ razy większy niż w oknie $B$ ($6{,}8384$ wobec $1{,}0209$), więc człon $\Re A^{(3)}f^3$ w \eqref{eq:etaser} przeważa nad przyrostem stężenia.

\begin{figure}[htbp]
\centering
\includegraphics[width=0.75\textwidth]{figs/real233_2D_N_interp_anomalia.png}
\caption{Klatka $i=233$: mapa liczby bakterii $N$ w oknach pływających o wymiarach $100\times100$ px z zaznaczonymi oknami $A$ i $B$ --- krańcami łańcucha pięciu okien o kolejnych liczbach $N$ (próg kroku 0,002). Etykiety okien podają położenie narożnika okna $(x_0,y_0)$ w pikselach, liczbę bakterii $N$, stężenie $f$, współczynnik $\Re A^{(3)}$ oraz lepkość efektywną, oznaczoną na rysunku symbolem $\mu_e/\mu$ (w tekście: $\eta_e/\eta$).}
\label{fig:real233}
\end{figure}

\begin{table}[htbp]
\centering\small
\caption{Wyniki dla skrajnych okien $A$ i $B$ w klatce $i=233$ ($r=3{,}5$ px; okno $A$: $(x_0,y_0)=(374,340)$, okno $B$: $(88,100)$).}\label{tab:real233}
\begin{tabular}{@{}ccccccc@{}}
\toprule
\textbf{Okno} & \textbf{$N$} & \textbf{$f$} & \textbf{$e_{33}$} & \textbf{$e_4'$} & \textbf{$\Re A^{(3)}$} & \textbf{$\eta_e/\eta$}\\
\midrule
A & $39$ & $0{,}1501$ & $11{,}0996$ & $3{,}8490$ & $6{,}8384$ & $1{,}4190$\\
B & $43$ & $0{,}1655$ & $5{,}1018$ & $-1{,}7219$ & $1{,}0209$ & $1{,}4100$\\
\bottomrule
\end{tabular}
\end{table}

W klatce $i=106$ okno $B$ zawiera więcej bakterii niż okno $A$ ($N_B=44>N_A=41$, stężenie większe o $7{,}3$\,\%), a mimo to ma mniejszą lepkość efektywną: $\eta_B/\eta=1{,}4229<\eta_A/\eta=1{,}4391$ (różnica $1{,}12$\,\%). Przyczyną jest układ środków bakterii: współczynnik $\Re A^{(3)}$ jest w oknie $A$ około $5{,}3$ razy większy niż w oknie $B$ ($5{,}6622$ wobec $1{,}0770$), więc człon $\Re A^{(3)}f^3$ w \eqref{eq:etaser} przeważa nad przyrostem stężenia.

\begin{figure}[htbp]
\centering
\includegraphics[width=0.75\textwidth]{figs/real106_2D_N_interp_anomalia.png}
\caption{Klatka $i=106$: mapa liczby bakterii $N$ w oknach pływających o wymiarach $100\times100$ px z zaznaczonymi oknami $A$ i $B$ --- krańcami łańcucha czterech okien o kolejnych liczbach $N$ (próg kroku 0,004). Etykiety okien podają położenie narożnika okna $(x_0,y_0)$ w pikselach, liczbę bakterii $N$, stężenie $f$, współczynnik $\Re A^{(3)}$ oraz lepkość efektywną, oznaczoną na rysunku symbolem $\mu_e/\mu$ (w tekście: $\eta_e/\eta$).}
\label{fig:real106}
\end{figure}

\begin{table}[htbp]
\centering\small
\caption{Wyniki dla skrajnych okien $A$ i $B$ w klatce $i=106$ ($r=3{,}5$ px; okno $A$: $(x_0,y_0)=(66,310)$, okno $B$: $(264,70)$).}\label{tab:real106}
\begin{tabular}{@{}ccccccc@{}}
\toprule
\textbf{Okno} & \textbf{$N$} & \textbf{$f$} & \textbf{$e_{33}$} & \textbf{$e_4'$} & \textbf{$\Re A^{(3)}$} & \textbf{$\eta_e/\eta$}\\
\midrule
A & $41$ & $0{,}1578$ & $10{,}4002$ & $2{,}3805$ & $5{,}6622$ & $1{,}4391$\\
B & $44$ & $0{,}1693$ & $3{,}9397$ & $-0{,}8548$ & $1{,}0770$ & $1{,}4229$\\
\bottomrule
\end{tabular}
\end{table}


\section{\texorpdfstring{Wnioski z przykładów}{Wnioski z przykładów}}

We wszystkich pięciu klatkach okno o większym stężeniu ($N_B$ większe od $N_A$ o 3--4 bakterie, stężenie wyższe o $7{,}3$--$10{,}5\,\%$) ma mniejszą lepkość efektywną, przy czym różnica jest niewielka (od $0{,}63$ do $2{,}08\,\%$). W każdym przypadku współczynnik $\Re A^{(3)}$ w oknie $B$ jest wyraźnie mniejszy niż w oknie $A$ (od $5{,}3$ do $10{,}8$ razy). W tabelach widać, że jest to skutek obu sum strukturalnych: $e_{33}$ w oknie $B$ jest $2{,}2$--$3{,}5$ razy mniejsze niż w oknie $A$, a $e_4'$ zmienia znak z dodatniego na ujemny. Oznacza to, że przy stężeniach $f\approx0{,}15$--$0{,}17$ o lepkości efektywnej zawiesiny nie decyduje wyłącznie stężenie, lecz także układ środków bakterii, opisany sumami strukturalnymi $e_{33}$ i $e_4$; wzór \eqref{eq:etareal} ujmuje to przez człon $\Re A^{(3)}f^3$. Zjawisko ma charakter lokalny i statystycznie rzadki (udział par dyskordantnych poniżej 4\,\%), a globalnie $\eta_e/\eta$ rośnie ze stężeniem ($\rho_S\approx0{,}98$). Wyniki dotyczą jednej realizacji rzeczywistej struktury (film \cite{ariel2018collective}) i wybranych klatek, nie wolno ich więc uogólniać na inne warunki doświadczalne.


\chapter{Podsumowanie}

Teoria opisana w \cite{mityushev2025computational,czapla2025computational,mityushev2025computationalpart3} opiera się na trzech elementach: (1) RVE jako obszarze fundamentalnym torusa, co usuwa zależność od kształtu próbki wynikającą z warunkowej zbieżności sum sieciowych; (2) dekompozycji stałych efektywnych na część fizyczną (parametry kontrastu) i geometryczną (sumy strukturalne o wadze $\alpha=2(s-1)$ w współczynniku rzędu $s$); (3) resumacji szeregów obciętych, dającej wzory zwarte w pełnym zakresie stężeń i oszacowanie wykładnika krytycznego $s\approx1{,}27$--$1{,}40$. Sumy strukturalne służą również jako cechy w uczeniu maszynowym, z dokładnością klasyfikacji protokołów $0{,}7125$.

Część badawcza zastosowała ten aparat do rzeczywistej struktury: położeń bakterii \emph{Bacillus subtilis} na 239 klatkach filmu Movie S1.avi, w pływających oknach $100\times100$ px. Dla każdego okna obliczono sumy strukturalne $e_{33}$ i $e_4$, współczynnik $A^{(3)}$ i efektywną lepkość $\eta_e/\eta$ (wzory \eqref{eq:Avisc} i \eqref{eq:etareal}). Lepkość jest silnie skorelowana ze stężeniem (średni współczynnik Spearmana $0{,}9825$), lecz pary okien, w których okno o mniejszym stężeniu ma większą lepkość, stanowią średnio $2{,}43\,\%$ par o różnym stężeniu; przykładem jest klatka $233$ (okno $A$: $39$ bakterii, $1{,}4190$; okno $B$: $43$ bakterie, $1{,}4100$). Zgodnie z teorią wyrazy $f$ i $f^2$ nie zależą od układu bakterii, a geometria wchodzi od wyrazu $f^3$ przez $A^{(3)}$. Wynik dotyczy jednego filmu i wymaga weryfikacji statystycznej (rozdz.~\ref{sec:realstat}).

Ograniczenia: symulacje dotyczą wyłącznie jednostajnego rozkładu równych dysków (trzech jego realizacji), $N=100$ i $f\le0{,}4$; wyniki liczbowe $c^{(3)},c^{(4)}$ zależą od protokołu, a bezpośrednie obliczenia dla modułu objętościowego $k_e$ nie zostały jeszcze opublikowane. Klasyfikacja ML jest umiarkowanie dokładna ($\mathrm{ACC}\approx0{,}71$), a interpretacja różnic między klasami wymaga ostrożności (Dodatek~\ref{app:notes}).

\textbf{Uwaga końcowa: powroty układu i twierdzenie Poincarégo.} Twierdzenie Poincarégo o powrotach \cite{poincare1890probleme} mówi, że jeśli ewolucja zachowuje skończoną miarę (np. objętość fazową w ograniczonej przestrzeni fazowej), to prawie każdy stan początkowy jest odwiedzany ponownie dowolnie blisko, nieskończenie wiele razy. Ma to sens w jednym miejscu tej pracy: dla konserwatywnej (hamiltonowskiej) dynamiki $N$ twardych dysków na torusie $\T^2$ przestrzeń konfiguracji jest zwarta, a energia stała, więc objętość fazowa jest skończona. Wspiera to założenie z p.~\ref{ssec:modellos}, że konfiguracje można traktować jako próbki rozkładu stacjonarnego (ergodyczność jest warunkiem silniejszym niż powrotność). Nie należy jednak rozszerzać tego na resztę teorii. (i) Protokoły R, T, P z rozdz.~\ref{sec:sim} są algorytmami losowymi lub iteracyjnymi, a nie dynamiką zachowującą miarę. (ii) Przepływ Stokesa jest dyssypatywny (energia jest rozpraszana, objętość fazowa się kurczy), więc dla lepkości twierdzenie nie ma zastosowania; równania Stokesa są natomiast kinematycznie odwracalne i odwrócenie ścinania przywraca położenia cząstek (przy braku kontaktów i ruchów Browna), co jest innym zjawiskiem niż powrót Poincarégo. (iii) Twierdzenie nie podaje czasu powrotu (dla $N\approx100$ astronomicznie długiego) i gwarantuje powrót dowolnie blisko, a nie dokładnie. (iv) Stałe efektywne w tej pracy są wielkościami statycznymi i ich wyprowadzenie nie wymaga tego twierdzenia. Wzmianka ma więc charakter uwagi o założeniach probabilistycznych.

\appendix
\chapter{Rozbieżności w tekstach źródłowych}\label{app:notes}
Przy opracowaniu zauważono następujące niespójności; w tekście przyjęto wersje wskazane poniżej.
\begin{enumerate}[nosep]
\item \textbf{Macierz pomyłek.} Tekst \cite{mityushev2025computationalpart3} mówi, że błędy występują głównie między R i T, a P jest najwyraźniejsza; macierz pokazuje pomyłki głównie P$\leftrightarrow$T i najlepsze rozpoznanie R. W pracy przyjęto dane z macierzy.
\item \textbf{Wariancja PCA.} Podpis rys.~3 \cite{mityushev2025computationalpart3} podaje $47{,}6\%+11{,}0\%=58{,}6\%$, a tekst „prawie 55\%”. Tekst mówi też o 14 cechach, a klasyfikator używa około 130.
\item \textbf{Numeracja równań w \cite{mityushev2025computationalpart3}.} Indeks $\eta^*$ opisano jako równanie (14) lub (17); jest to równanie (18). W \cite{mityushev2025computationalpart3} przykład $z_A$ stosuje sumowanie Eisensteina (wartości $e_2=\pi$, $e_3^{(1)(1)}=\pi/2$), podczas gdy dla sprężystości obowiązuje sumowanie Maxwella ($S_2=0$, $e_3^{(1)(1)}=0$).
\item \textbf{Waga a suma wskaźników.} W \cite{mityushev2025computationalpart3} stwierdzono, że współczynnik przy $f^{k+1}$ zawiera sumy z $n_1+\dots+n_p=2k$. Z wzorów (34)--(43) w \cite{czapla2025computational} wynika, że niezmiennikiem jest waga $\alpha=\sum(n_s-j_s)$ (np.\ $e_3^{(1)(1)}$ ma $\sum n_s=3$, lecz $\alpha=2$). Użyto wagi $\alpha$.
\item \textbf{Rząd dokładności przy $\delta\ne0$.} W \cite{czapla2025computational} stwierdzono jednocześnie, że $\delta\neq0$ „do $\mathcal O(f^3)$” i że $\Delta\mu_L=\mathcal O(f^2)$, a wynik jest poprawny „do $\mathcal O(f^2)$”. Ścisłe sformułowanie wymaga odwołania do \cite{rayleigh1892lvi,mcphedran1978conductivity,mckenzie1978conductivity,mityushev1997transport}.
\item \textbf{Stała $\kappa$.} Rozdz.~3 w \cite{czapla2025computational} przyjmuje $\kappa=3-4\nu$, a rozdz.~5 $\kappa=(3-\nu)/(1+\nu)$ (płaski stan naprężenia). Przyjęto $\kappa=3-4\nu$ i odnotowano drugi przypadek.
\item \textbf{Próg perkolacji.} W \cite{czapla2025computational} zapisano $f_c=\sqrt\pi/12$; poprawnie $f_c=\pi/(2\sqrt3)\approx0{,}9069$.
\item \textbf{Etykiety wielomianów.} W \cite{czapla2025computational} zdanie przypisujące równania (73)--(76) do protokołu P, a (77)--(80) do T, jest zamienione względem oznaczeń równań ($\mu_T$, $\mu_P$); wielomiany $\mu^P_{e,4}$ i $\mu^T_{e,4}$ o podanych współczynnikach odpowiadają poprawnie ich nazwom (sprawdzono średnimi z tabeli~\ref{tab:coef}).
\item \textbf{Odwołanie do części III.} W \cite{czapla2025computational} obliczenie $k_e$ odroczono do „Part III”, która okazała się poświęcona uczeniu maszynowemu.
\item \textbf{Źródło pliku Movie S1.avi.} Artykuł ECMS \cite{ligocki2025study} wskazuje jako źródło filmu pracę \cite{ariel2018collective}, której abstrakt dotyczy bakterii \emph{Serratia marcescens}, natomiast projekt badawczy wskazuje na pracę \cite{ryan2013kinetic} i bakterie \emph{Bacillus subtilis}. W rozprawie przyjęto źródło z projektu badawczego; wymaga to potwierdzenia w opisie filmu.

\item\label{it:weryf} \czer{\textbf{Zgodność wzorów na sumy strukturalne z kodem v16 (do wyjaśnienia z promotorem).} Wzory z p.~\ref{ssec:struct} i rozdz.~\ref{sec:kod} porównano z notebookiem v16 (komórki~3d i~8). Sumy policzono wprost z definicji (potrójne sumy po indeksach) dla okien $C$ i $D$ i porównano z wynikami kodu.
\begin{itemize}[nosep]
\item \textbf{Zgodne.} (a) $\wp$, $\wp'$, $\zeta$ z kodu zgadzają się z sumami sieciowymi do $10^{-5}$; $\zeta'=-\wp$; kwazi-okresy $\eta_1=\pi$, $\eta_2=-i\pi$. (b) $S_4=3{,}151212$, $g_2=189{,}0727$, $S_3^{(1)}=2{,}50765$. (c) Zaimplementowana funkcja \eqref{eq:E31} jest okresowa (różnica $10^{-14}$), więc redukcja do najbliższego obrazu torusowego jest poprawna. (d) $e_{33}$ i $\Re e_4$ z kodu równają się wartościom z definicji (okno $C$: $13{,}143685$ i $1{,}486867$; okno $D$: $10{,}407683$ i $-0{,}641933$). (e) $e_2=\pi+N^{-2}\sum\wp$ i $e_{2,2}=N^{-3}\sum_k|\sum_m(\wp+\pi)|^2$ odpowiadają definicji dla $E_2=\wp+\pi$ (okno $C$: $e_2=3{,}0770-0{,}2058\,i$, $e_{2,2}=18{,}3637$); dotyczy to funkcji \texttt{prepare}, a plik \texttt{pfast.py} nie był dostępny. (f) Wzory $A^{(3)}=\frac2{\pi^2}(2e_{33}+3\Re e_4)$ oraz $\eta_e/\eta=(1+s)/(1-s)$ odtwarzają wartości z pliku CSV.
\item \textbf{$E_3^{(1)}(0)$.} Definicja: $S_3^{(1)}=2{,}5077$; kod: $0$. Okno $C$: $e_{33}=13{,}1437$ (kod) i $13{,}2287$ (Definicja), $\Re A^{(3)}=6{,}2308$ i $6{,}2653$, $\eta_e/\eta=1{,}41300$ i $1{,}41334$. Okno $D$: $10{,}4077$ i $10{,}3552$, $3{,}8278$ i $3{,}8066$, $1{,}44747$ i $1{,}44718$. Dla 300 losowych okien: różnica $e_{33}$ maks.\ $0{,}15$ (typowo $e_{33}\approx9{,}8$), różnica $\eta_e/\eta$ średnio $5\cdot10^{-5}$, maks.\ $7{,}4\cdot10^{-4}$, przy odchyleniu standardowym $0{,}023$.
\item \textbf{Konwencja $S_2$.} Kod w $e_2$ i $e_{2,2}$ używa $S_2=\pi$ (Eisenstein), p.~\ref{ssec:visc} zakłada $S_2=0$ (Maxwell), a rozdz.~\ref{sec:kod} nie wspomina o składniku $+\pi$. Dla $e_{33}$ i $e_4$ bez znaczenia; $e_{22}$ nie występuje w \eqref{eq:Avisc}, bo $\varrho_1=\varrho_3$.
\item \textbf{$A^{(2)}$ i $e_3^{(1)(1)}$.} Praca definiuje $e_3^{(1)(1)}$, ale kod go nie liczy i przyjmuje się $A^{(2)}=0$. Dla 300 losowych okien $|e_3^{(1)(1)}|$ ma średnią $0{,}56$ (maks.\ $1{,}5$). Roboczy składnik $A^{(2)}=-\frac2\pi\Re e_3^{(1)(1)}$ (nie z prac źródłowych) zmieniałby $\eta_e/\eta$ średnio o $0{,}017$ (maks.\ $0{,}063$).
\item \textbf{Suma po kwadratach a $S_3^{(1)}$.} Suma $\sum\cc P/P^3$ po kwadratach daje $2{,}0785$ (nie $2{,}5077$); $E_3^{(1)}(0{,}23+0{,}11\,i)$ z sumy po kwadratach to $-1{,}8557-14{,}7804\,i$, a ze wzoru \eqref{eq:E31} $-1{,}4265-14{,}7804\,i$. Różnica $0{,}4292$ jest stałą addytywną, która nie wpływa na okresowość, więc wybór $S_3^{(1)}$ wynika z reguły sumowania i wymaga potwierdzenia.
\item \oran{\textbf{Porównanie z Mathematicą (\texttt{Import-e2-CORRECTED.nb}).} Dla zbioru $182$ punktów z tego pliku (\texttt{im182a}, $a=x+iy$) Mathematica po poprawce daje $e_2=3{,}3256144+0{,}1205367\,i$ i $e_{2,2}=15{,}283710$. Wzory z kodu v16 dają dokładnie te wartości: z $\wp$ z komórki~3d $3{,}3256144+0{,}1205367\,i$ i $15{,}283710$, a z funkcji \texttt{p\_weierstrass\_matrix\_fast} z \texttt{pfast.py} $3{,}3256143+0{,}1205367\,i$ i $15{,}283709$. Stare wyniki zapisane w pliku ($2{,}9575716+0{,}1205385\,i$ i $12{,}971242$) odtwarza zamiana osi, $a=y+ix$ ($2{,}9575710$ i $12{,}971224$); ponieważ $\wp(iz)=-\wp(z)$, zamiana daje $-\overline{\wp}$ zamiast $\wp$ i zmienia $\Re e_2$ oraz $e_{2,2}$, a nie zmienia $\Im e_2$. Wyniki w rozprawie policzono wersją v16 z $a=x+iy$ (Listing~\ref{lst:kod-norm}). Wzory na $e_2$ i $e_{2,2}$ w kodzie i w \texttt{.nb} są takie same ($E_2=\wp+\pi$, $E_2(0)=\pi$, $\wp(0)=0$, $\pi$ dodane raz do $e_2$). Uwaga: \texttt{pfast.py} sprawdzono z własnym, minimalnym zamiennikiem pakietu \texttt{mpmath} (pakietu nie było w środowisku), więc potwierdzenie przebiegiem z prawdziwym \texttt{mpmath} jest wskazane.}
\item \oran{\textbf{Funkcja $E_3^{(1)}$ w \texttt{Import-e2-CORRECTED.nb}.} Plik definiuje \texttt{E31} bez stałej $S_3^{(1)}$ (kod i praca mają $+S_3^{(1)}$ we wzorze \eqref{eq:E31}), a na przekątnej daje $-4S_4/\pi$ (wynika z \texttt{E3(0)=0}, \texttt{E4(0)=S4}, \texttt{zeta(0)=0}), a nie $0$ jak kod ani $S_3^{(1)}$ jak Definicja~\ref{def:struct}. W komórkach z wynikami \texttt{E31} nie jest używana (liczone są tylko $e_2$ i $e_{2,2}$), więc zgodność wyników nie zależy od tej różnicy. Dla $e_{33}$ plik nie stanowi odniesienia; przed użyciem go do $e_{33}$ definicję trzeba ujednolicić.}
\end{itemize}}
\end{enumerate}

% =========================================================================
% BIBLIOGRAFIA (klucze zgodne z plikiem promotora biblioVM_1.bib)
% =========================================================================
\clearpage
\phantomsection
\addcontentsline{toc}{chapter}{Bibliografia}
\bibliographystyle{plain}
\bibliography{biblioVM_1,biblioLigocki}

\end{document}
